% Leçon 17 — Vocabulaire et compréhension de texte : collectivités locales — Chinois Tle A — Cours signé OnBuch+
% Programme officiel MINESEC. Compiler avec : tectonic ZH17-vocabulaire-collectivites-locales.tex
\newcommand{\DOCMATIERE}{Chinois}
\newcommand{\DOCNIVEAU}{Tle A}
\newcommand{\DOCMODULE}{Module 3 : Citoyenneté et ouverture au monde}
\newcommand{\DOCLECON}{ZH17}
\newcommand{\DOCTITRE}{Leçon 17 — Vocabulaire et compréhension de texte : collectivités locales}
% OnBuch+ — préambule « cours signés » ENRICHI (compilé avec tectonic / XeLaTeX)
% Système visuel « Pop » d'OnBuch+ : fond crème (jamais blanc), bordures encre
% épaisses, ombres « dures » décalées, titres Archivo Black en capitales.
% Version MATHÉMATIQUES : graphes (pgfplots/tikz), boîtes pédagogiques
% (définition, propriété, méthode, attention, le savais-tu, exercice résolu,
% exercice, TP, bilan), page de garde et signature OnBuch+ ancrée en bas.
%
% Macros attendues dans main.tex AVANT \input{preamble} :
%   \DOCMATIERE  \DOCNIVEAU  \DOCMODULE  \DOCLECON (numéro)  \DOCTITRE
\documentclass[11pt]{article}

\usepackage{fontspec}
\usepackage[french]{babel} % français = langue principale ; le chinois via \zh{..} / textebox (xeCJK)
\usepackage{xeCJK}
\usepackage{pifont} % \ding{51} \ding{55} utilisés par les modèles
\usepackage[a4paper,margin=2cm,top=3.4cm,bottom=2.8cm,headheight=1.4cm,footskip=1.3cm]{geometry}
\usepackage{amsmath,amssymb,amsfonts}
\usepackage{mathtools} % \xmapsto, \coloneqq... souvent utilisés par les modèles
\usepackage{cancel} % simplifications barrées (bilans de forces)
% \abs{z} (module, valeur absolue) et \norm{u} (norme), utilisés par les modèles
\providecommand{\abs}{}\let\abs\relax\DeclarePairedDelimiter{\abs}{\lvert}{\rvert}
\providecommand{\norm}{}\let\norm\relax\DeclarePairedDelimiter{\norm}{\lVert}{\rVert}
\usepackage[table]{xcolor} % [table] : fournit aussi \rowcolors (alternance de lignes)
\usepackage{pagecolor}
\usepackage{tikz}
\usetikzlibrary{arrows.meta,positioning,calc,shapes.geometric,shapes.misc,shapes.arrows,decorations.pathmorphing,decorations.pathreplacing,decorations.markings,patterns,shadows,mindmap,trees,fit,backgrounds,matrix,intersections,angles,quotes}
\usepackage{pgfplots}
\usepackage[version=4]{mhchem} % équations nucléaires \ce{^{235}_{92}U}
\usepackage[european,siunitx]{circuitikz} % schémas électriques
\pgfplotsset{compat=1.18}
\usepgfplotslibrary{fillbetween,groupplots} % aires entre courbes (calcul intégral)
\usepackage{enumitem}
\usepackage{fancyhdr}
% [most] charge toutes les bibliothèques tcolorbox (listings, minted, raster,
% documentation...) et gonfle inutilement la mémoire TeX sur un document long
% (~300 boîtes) : on ne charge que ce qui est réellement utilisé.
\usepackage[skins,breakable]{tcolorbox}
\usepackage{graphicx}
\usepackage{colortbl}
\usepackage{booktabs}
\usepackage{tabularx}
\usepackage{diagbox} % case d'en-tête diagonale des tableaux à double entrée
% Colonnes extensibles centrée / gauche / droite (C, L, R), souvent utilisées par les modèles
\newcolumntype{C}{>{\centering\arraybackslash}X}
\newcolumntype{L}{>{\raggedright\arraybackslash}X}
\newcolumntype{R}{>{\raggedleft\arraybackslash}X}
\usepackage{array}
\usepackage{multicol}
\usepackage{multirow}
\usepackage{siunitx}
% parse-numbers=false : accepte aussi \SI{2,5 \times 10^{-3}}{...}
\sisetup{locale=FR,output-decimal-marker={,},group-minimum-digits=4,parse-numbers=false}
\DeclareSIUnit\ohm{\ensuremath{\symup{\Omega}}} % U+2126 absent de la police
\DeclareSIUnit\division{div} % réglages d'oscilloscope (ms/div, V/div)
\DeclareSIUnit\tr{tr} % tours (vitesse de rotation, ex. \SI{1500}{\tr\per\minute})
\DeclareSIUnit\molar{M} % concentration molaire (ex. \SI{3}{\molar}, \SI{1}{\micro\molar})
\DeclareSIUnit\an{an} % année (le modèle confond parfois avec \year, un compteur TeX) — ex. \SI{5}{\kilo\meter\per\an}
\DeclareSIUnit\FCFA{FCFA} % franc CFA (ex. \SI{500}{\FCFA\per\kilo\gram})
\DeclareSIUnit\degreeBrix{\textdegree Brix} % teneur en sucre d'un sirop/jus (réfractométrie)
\DeclareSIUnit\month{mois} % le modèle utilise parfois \month, un compteur TeX, comme unité de durée
\usepackage{qrcode}
% \degree : commande fréquemment utilisée par les modèles pour un angle ou une
% température en dehors de \SI{}{} (non fournie par les paquets chargés).
\providecommand{\degree}{^{\circ}}
% \qed (fin de démonstration), utilisé par les modèles sans amsthm
\providecommand{\qed}{\hfill$\square$}
% \barre : double barre des valeurs interdites dans un tableau de variations (tkz-tab)
\providecommand{\barre}{\ensuremath{\|}}
% environnement proof (amsthm non chargé) utilisé par les modèles
\ifdefined\proof\else\newenvironment{proof}[1][Démonstration]{\par\noindent\textit{#1.}\ }{\qed\par}\fi
\usepackage{lastpage}
\usepackage{hyperref}
\hypersetup{hidelinks}

% ============================================================
% Palette « Pop » OnBuch+ — mêmes hex que la classe P (plus_ui.dart)
% ============================================================
\definecolor{poporange}{HTML}{F59321}
\definecolor{popdark}{HTML}{E07A0C}
\definecolor{popcream}{HTML}{FFF6E8}
\definecolor{popink}{HTML}{1C1714}
\definecolor{poppurple}{HTML}{7A5AE0}
\definecolor{popgreen}{HTML}{2E9E6A}
\definecolor{poppink}{HTML}{E0457B}
\definecolor{popblue}{HTML}{2F7DE1}
\definecolor{popgold}{HTML}{C98A12}
\definecolor{popmuted}{HTML}{8A7F76}
\definecolor{popline}{HTML}{EADFD2}
\definecolor{popgray}{HTML}{8A7F76}
\colorlet{popgrey}{popgray}
% Alias tolérants (le modèle nomme parfois les couleurs différemment)
\colorlet{popwhite}{white}
\colorlet{popblack}{popink}
\colorlet{popred}{poppink}
\colorlet{popyellow}{popgold}
% Teintes claires (fonds de tableaux/encadrés) — le modèle nomme parfois
% les couleurs avec un suffixe L (ex. popblueL) sans les avoir définies.
\colorlet{poporangeL}{poporange!20}
\colorlet{popdarkL}{popdark!20}
\colorlet{poppurpleL}{poppurple!20}
\colorlet{popgreenL}{popgreen!20}
\colorlet{poppinkL}{poppink!20}
\colorlet{popblueL}{popblue!20}
\colorlet{popgoldL}{popgold!20}
\colorlet{popredL}{popred!20}
\colorlet{popgrayL}{popgray!20}
% Teintes claires pour fonds de tableaux / zones
\colorlet{popblueL}{popblue!10!white}
\colorlet{poporangeL}{poporange!14!white}
\colorlet{popgreenL}{popgreen!12!white}
\colorlet{poppurpleL}{poppurple!12!white}
\colorlet{poppinkL}{poppink!12!white}
\colorlet{popgoldL}{popgold!14!white}

\pagecolor{popcream}

% ============================================================
% Typographie
% ============================================================
\defaultfontfeatures{Ligatures=TeX}
\setmainfont{PlusJakartaSans-Medium.ttf}[Path=../../../fonts/,BoldFont=PlusJakartaSans-Bold.ttf]
% Symboles, API et emojis absents de la police principale : repli sur DejaVu Sans (caractères actifs)
\newfontface\symfont{DejaVuSans.ttf}[Path=../../../fonts/]
\makeatletter
\newcommand\symactive[1]{\catcode#1=13 \begingroup\lccode`\~=#1 \lowercase{\endgroup\def~}{{\symfont\char#1}}}
\newcommand\symblank[1]{\catcode#1=13 \begingroup\lccode`\~=#1 \lowercase{\endgroup\def~}{}}
\newcommand\symhyph[1]{\catcode#1=13 \begingroup\lccode`\~=#1 \lowercase{\endgroup\def~}{\mbox{-}}}
\newcount\symc
\newcommand\symrange[2]{\symc=#1 \@whilenum\symc<#2 \do{\expandafter\symactive\expandafter{\the\symc}\advance\symc by 1 }}
\newcommand\symblankrange[2]{\symc=#1 \@whilenum\symc<#2 \do{\expandafter\symblank\expandafter{\the\symc}\advance\symc by 1 }}
\makeatother
\symrange{"0250}{"02FF}\symactive{"2713}\symactive{"2714}\symactive{"2717}\symactive{"2718}\symactive{"2610}\symactive{"2611}\symactive{"2612}
\symrange{"2600}{"27BF}
\catcode"2705=13 \begingroup\lccode`\~="2705 \lowercase{\endgroup\def~}{{\symfont\char"2713}}\symblank{"26BD}\symblank{"26A1}
\symrange{"2460}{"257F}\symactive{"223C}\symactive{"2248}\symactive{"2260}
\symhyph{"2011}\symhyph{"2E1A}\symblankrange{"1F300}{"1FAFF}\symblank{"FE0F}
% Caractères chinois : WenQuanYi Zen Hei (sous-ensemble GB2312 + pinyin), gras/italique simulés
\setCJKmainfont{WenQuanYiZenHei-subset.ttf}[Path=../../../fonts/,AutoFakeBold=3,AutoFakeSlant=0.2]
\xeCJKsetup{CJKspace=false}
% Pinyin : voyelles à caron (ǎ ǐ ǒ ǔ ǖ ǘ ǚ ǜ) absentes de la police principale -> caractères actifs
\newfontface\pyfont{WenQuanYiZenHei-subset.ttf}[Path=../../../fonts/]
\newcommand\pyactive[1]{\catcode#1=13 \begingroup\lccode`\~=#1 \lowercase{\endgroup\def~}{{\pyfont\char#1}}}
\pyactive{"01CD}\pyactive{"01CE}\pyactive{"01CF}\pyactive{"01D0}\pyactive{"01D1}\pyactive{"01D2}\pyactive{"01D3}\pyactive{"01D4}\pyactive{"01D5}\pyactive{"01D6}\pyactive{"01D7}\pyactive{"01D8}\pyactive{"01D9}\pyactive{"01DA}\pyactive{"01DB}\pyactive{"01DC}
% Maths en sans-serif (Fira Math) avec chiffres et lettres droites de Plus
% Jakarta, pour que les formules se fondent dans le texte.
\usepackage[math-style=ISO,bold-style=ISO]{unicode-math}
\setmathfont{FiraMath-Regular.otf}[Path=../../../fonts/]
\setmathfont{PlusJakartaSans-Medium.ttf}[Path=../../../fonts/,range={up/num,up/{Latin,latin},\mathup}]
\usepackage{icomma} % virgule décimale sans espace en mode math
% Caractères mathématiques Unicode absents des polices de texte (Plus Jakarta,
% Archivo Black) : rendus par leur équivalent mathématique, en texte comme en
% formule — sinon ils s'affichent en carré vide (ex. « Arithmétique dans ℤ »).
\usepackage{newunicodechar}
\newunicodechar{ℝ}{\ensuremath{\mathbb{R}}}
\newunicodechar{ℤ}{\ensuremath{\mathbb{Z}}}
\newunicodechar{ℂ}{\ensuremath{\mathbb{C}}}
\newunicodechar{ℕ}{\ensuremath{\mathbb{N}}}
\newunicodechar{ℚ}{\ensuremath{\mathbb{Q}}}
\newunicodechar{𝔻}{\ensuremath{\mathbb{D}}}
\newunicodechar{α}{\ensuremath{\alpha}}
\newunicodechar{β}{\ensuremath{\beta}}
\newunicodechar{γ}{\ensuremath{\gamma}}
\newunicodechar{δ}{\ensuremath{\delta}}
\newunicodechar{ε}{\ensuremath{\varepsilon}}
\newunicodechar{θ}{\ensuremath{\theta}}
\newunicodechar{λ}{\ensuremath{\lambda}}
\newunicodechar{μ}{\ensuremath{\mu}}
\newunicodechar{σ}{\ensuremath{\sigma}}
\newunicodechar{τ}{\ensuremath{\tau}}
\newunicodechar{φ}{\ensuremath{\varphi}}
\newunicodechar{ψ}{\ensuremath{\psi}}
\newunicodechar{ω}{\ensuremath{\omega}}
\newunicodechar{Σ}{\ensuremath{\Sigma}}
% majuscules produites par \MakeUppercase dans les titres (α-aminés -> Α)
\newunicodechar{Α}{\ensuremath{\alpha}}
\newunicodechar{Β}{\ensuremath{\beta}}
\newunicodechar{Γ}{\ensuremath{\Gamma}}
\newunicodechar{Δ}{\ensuremath{\Delta}}
\newunicodechar{Ε}{\ensuremath{\varepsilon}}
\newunicodechar{Θ}{\ensuremath{\Theta}}
\newunicodechar{Λ}{\ensuremath{\Lambda}}
\newunicodechar{Μ}{\ensuremath{\mu}}
\newunicodechar{Τ}{\ensuremath{\tau}}
\newunicodechar{Φ}{\ensuremath{\Phi}}
\newunicodechar{Ψ}{\ensuremath{\Psi}}
\newunicodechar{Ω}{\ensuremath{\Omega}}
\newunicodechar{↦}{\ensuremath{\mapsto}}
\newunicodechar{∈}{\ensuremath{\in}}
\newunicodechar{∉}{\ensuremath{\notin}}
\newunicodechar{∧}{\ensuremath{\wedge}}
\newunicodechar{⇔}{\ensuremath{\Leftrightarrow}}
\newunicodechar{⟺}{\ensuremath{\Longleftrightarrow}}
\newunicodechar{⇒}{\ensuremath{\Rightarrow}}
\newunicodechar{⟹}{\ensuremath{\Longrightarrow}}
\newunicodechar{∘}{\ensuremath{\circ}}
\newunicodechar{∪}{\ensuremath{\cup}}
\newunicodechar{∩}{\ensuremath{\cap}}
\newunicodechar{≡}{\ensuremath{\equiv}}
\newunicodechar{⊂}{\ensuremath{\subset}}
\newunicodechar{⊥}{\ensuremath{\perp}}
\newunicodechar{∃}{\ensuremath{\exists}}
\newunicodechar{∀}{\ensuremath{\forall}}
\newunicodechar{∛}{\ensuremath{\sqrt[3]{\,}}}
\newunicodechar{∜}{\ensuremath{\sqrt[4]{\,}}}
\newunicodechar{⃗}{\ensuremath{{}^{\to}}}
\newunicodechar{ⁿ}{\ifmmode^{n}\else\textsuperscript{\ensuremath{n}}\fi}
\newunicodechar{ˣ}{\ifmmode^{x}\else\textsuperscript{\ensuremath{x}}\fi}
\newunicodechar{ᵇ}{\ifmmode^{b}\else\textsuperscript{\ensuremath{b}}\fi}
\newunicodechar{ʳ}{\ifmmode^{r}\else\textsuperscript{\ensuremath{r}}\fi}
\newunicodechar{ᵘ}{\ifmmode^{u}\else\textsuperscript{\ensuremath{u}}\fi}
\newunicodechar{ʸ}{\ifmmode^{y}\else\textsuperscript{\ensuremath{y}}\fi}
\newunicodechar{ᵃ}{\ifmmode^{a}\else\textsuperscript{\ensuremath{a}}\fi}
\newunicodechar{ᵉ}{\ifmmode^{e}\else\textsuperscript{\ensuremath{e}}\fi}
\newunicodechar{ᵏ}{\ifmmode^{k}\else\textsuperscript{\ensuremath{k}}\fi}
\newunicodechar{ᵖ}{\ifmmode^{p}\else\textsuperscript{\ensuremath{p}}\fi}
\newunicodechar{ᴺ}{\ifmmode^{N}\else\textsuperscript{\ensuremath{N}}\fi}
\newunicodechar{ᶜ}{\ifmmode^{c}\else\textsuperscript{\ensuremath{c}}\fi}
\newunicodechar{ʰ}{\ifmmode^{h}\else\textsuperscript{\ensuremath{h}}\fi}
\newunicodechar{ᵠ}{\ifmmode^{q}\else\textsuperscript{\ensuremath{q}}\fi}
\newunicodechar{ₙ}{\ifmmode_{n}\else\textsubscript{\ensuremath{n}}\fi}
\newunicodechar{ₐ}{\ifmmode_{a}\else\textsubscript{\ensuremath{a}}\fi}
\newunicodechar{ᵢ}{\ifmmode_{i}\else\textsubscript{\ensuremath{i}}\fi}
\newunicodechar{ᵣ}{\ifmmode_{r}\else\textsubscript{\ensuremath{r}}\fi}
\newunicodechar{ₖ}{\ifmmode_{k}\else\textsubscript{\ensuremath{k}}\fi}
\newunicodechar{ᵦ}{\ifmmode_{\beta}\else\textsubscript{\ensuremath{\beta}}\fi}
\newunicodechar{ₓ}{\ifmmode_{x}\else\textsubscript{\ensuremath{x}}\fi}
\newfontfamily\popdisplay{ArchivoBlack-Regular.ttf}[Path=../../../fonts/]
\newfontfamily\poplogo{SpaceGrotesk-Bold.ttf}[Path=../../../fonts/]
\newfontfamily\popsemi{PlusJakartaSans-SemiBold.ttf}[Path=../../../fonts/]
\newfontfamily\popxbold{PlusJakartaSans-ExtraBold.ttf}[Path=../../../fonts/]
\newfontfamily\popxboldls{PlusJakartaSans-ExtraBold.ttf}[Path=../../../fonts/,LetterSpace=3.0]

\setlist[enumerate]{leftmargin=1.7em,itemsep=5pt,topsep=4pt}
\setlist[itemize]{leftmargin=1.7em,itemsep=4pt,topsep=4pt,label={\color{poporange}\textbullet}}
\setlength{\parindent}{0pt}
\setlength{\parskip}{5pt}
\renewcommand{\arraystretch}{1.25}

% pgfplots : style Pop par défaut
\pgfplotsset{
  every axis/.append style={axis line style={popink,line width=1pt},tick style={popink},
    grid style={popline,line width=0.5pt},label style={font=\small},tick label style={font=\footnotesize}},
  popcurve/.style={poporange,line width=1.8pt,no markers,smooth},
}
% Même style disponible dans un \draw TikZ simple (hors pgfplots)
\tikzset{popcurve/.style={poporange,line width=1.8pt}}
\tikzset{popgrid/.style={popline,very thin}}
\colorlet{popcurve}{poporange} % le nom du style est parfois utilisé comme couleur

% ============================================================
% Pill / squiggle
% ============================================================
\newcommand{\poppill}[3]{%
  \tikz[baseline=-0.55ex]\node[draw=popink,line width=1.2pt,rounded corners=2.6pt,%
    fill=#2,inner xsep=6.5pt,inner ysep=3pt]{%
    \popxboldls\fontsize{7}{7}\selectfont\color{#3}\MakeUppercase{#1}};%
}
\newcommand{\popsquiggle}[1]{%
  \tikz[baseline=-0.4ex]{
    \draw[#1,line width=2.6pt,line cap=round]
      plot [smooth] coordinates {(0,0.09) (0.34,0.01) (0.68,0.21) (1.02,0.01) (1.36,0.10)};
  }%
}
% Mot-clé surligné (marqueur orange sous le texte)
\newcommand{\cle}[1]{{\popxbold\color{popdark}#1}}

% ============================================================
% En-tête / pied
% ============================================================
\pagestyle{fancy}
\fancyhf{}
\renewcommand{\headrulewidth}{0pt}
\renewcommand{\footrulewidth}{0pt}
\lhead{\raisebox{-0.15ex}{\poplogo\fontsize{13}{13}\selectfont\color{popink}OnBuch\color{poporange}+}}
\rhead{\poppill{\DOCMATIERE\ \textbullet\ \DOCNIVEAU\ \textbullet\ Leçon \DOCLECON}{white}{popink}}
\lfoot{\popxbold\fontsize{7.6}{7.6}\selectfont\color{popmuted}\MakeUppercase{Cours signé OnBuch+}\ \textbullet\ {\popsemi\DOCTITRE}}
\rfoot{\tikz[baseline=-0.6ex]\node[rounded corners=8.5pt,fill=popink,minimum height=17pt,minimum width=17pt,inner xsep=5pt,inner sep=0]{\popxbold\fontsize{8}{8}\selectfont\color{popcream}\thepage\,/\,\pageref*{LastPage}};}

% ============================================================
% Titres
% ============================================================
\newcommand{\chaptitre}[2]{%
  \begin{tcolorbox}[enhanced,colback=poporange,colframe=popink,boxrule=2.6pt,arc=11pt,
      drop shadow=popink,left=15pt,right=15pt,top=12pt,bottom=11pt,before skip=3pt,after skip=18pt]
    \poppill{#2}{popink}{popcream}\par\vspace{9pt}
    {\raggedright\hyphenpenalty=10000\exhyphenpenalty=10000\popdisplay\fontsize{22}{24}\selectfont\color{popcream}\MakeUppercase{#1}\par}
  \end{tcolorbox}
}
% Section numérotée : pastille orange avec le numéro + titre Archivo
\newcounter{popsec}
\newcommand{\coursec}[1]{%
  \stepcounter{popsec}\setcounter{popsub}{0}%
  \par\vspace{16pt}\noindent
  \begin{minipage}{\linewidth}
  \tikz[baseline=-0.7ex]\node[draw=popink,line width=1.6pt,fill=poporange,rounded corners=4pt,minimum size=22pt,inner sep=0]{\popdisplay\fontsize{12}{12}\selectfont\color{popcream}\thepopsec};%
  \hspace{8pt}{\popdisplay\fontsize{15}{17}\selectfont\color{popink}\MakeUppercase{#1}}\par
  \vspace{4pt}\noindent{\color{poporange}\rule{36pt}{3.6pt}}
  \end{minipage}\par\nopagebreak\vspace{8pt}
}
\newcounter{popsub}
\newcommand{\courssub}[1]{%
  \stepcounter{popsub}%
  \par\vspace{9pt}\noindent{\popxbold\fontsize{11.5}{13}\selectfont\color{popink}{\color{poporange}\thepopsec.\thepopsub}\hspace{6pt}#1}\par\nopagebreak\vspace{4pt}
}

% ============================================================
% Boîtes pédagogiques — même recette : fond blanc, bordure épaisse colorée,
% ombre dure encre, badge plein. Titre optionnel [#1].
% ============================================================
\tcbset{popbase/.style={breakable,enhanced,colback=white,boxrule=2.3pt,arc=9pt,
  shadow={3pt}{-3pt}{0pt}{popink},left=14pt,right=14pt,top=11pt,bottom=11pt,before skip=11pt,after skip=15pt,
  fontupper=\popsemi\fontsize{10.6}{14.5}\selectfont\color{popink},
  fontlower=\popsemi\fontsize{10.6}{14.5}\selectfont\color{popink}}}
% Coche dessinée (le glyphe ✓ n'existe pas dans les polices du document)
\newcommand{\popcheck}[1]{\tikz[baseline=-0.5ex]\draw[#1,line width=1.6pt,line cap=round,line join=round](-0.13,0.0)--(-0.04,-0.09)--(0.13,0.1);}
\providecommand{\coche}{\popcheck{popgreen}}
\providecommand{\resultat}[1]{\par\smallskip\noindent\fcolorbox{popgreen}{popgreenL}{\parbox{\dimexpr\linewidth-2\fboxsep-2\fboxrule\relax}{\textbf{Résultat.} #1}}\par\smallskip}
\newcommand{\popboxhead}[3]{\poppill{#1}{#2}{white}\ifx\relax#3\relax\else\hspace{6pt}{\popxbold\fontsize{10.6}{12}\selectfont\color{popink}#3}\fi\par\vspace{7pt}}

\newtcolorbox{aretenir}[1][]{popbase,colframe=popgreen,before upper={\popboxhead{À retenir}{popgreen}{#1}}}
% Texte en chinois (document de lecture, dialogue, extrait) : cadre
% d'étude. Un titre optionnel (source, auteur) s'affiche dans l'en-tête.
\newtcolorbox{textebox}[1][]{popbase,colframe=popblue,before upper={\popboxhead{文本 · Texte}{popblue}{#1}},after upper={}}
% Marques ✓ / ✗ (forme correcte / fautive) souvent utilisées par les modèles
\providecommand{\cross}{\textcolor{poppink}{\ensuremath{\times}}}
\providecommand{\xmark}{\cross}\providecommand{\crossmark}{\cross}\providecommand{\cmark}{\ensuremath{\checkmark}}
% Ligne de réponse / justification à compléter (inventée par les modèles)
\providecommand{\justification}[1]{\par\noindent\textit{Justification~:} \dotfill\par}
% \textipa (package tipa non chargé) : la police ne couvre pas l'API -> simple passage
\providecommand{\textipa}[1]{#1}
% Soulignement (ulem non chargé) : \uline, \ul
\providecommand{\uline}[1]{\underline{#1}}\providecommand{\ul}[1]{\underline{#1}}
% \glossaire{\item ...} inventée par les modèles : liste de glossaire
\providecommand{\glossaire}[1]{\begin{itemize}#1\end{itemize}}
% \alert (beamer) inventée par les modèles : texte mis en évidence
\providecommand{\alert}[1]{\textbf{\textcolor{poppink}{#1}}}
% variantes ASCII d'ordinaux écrites en macro
\providecommand{\iere}{\textsuperscript{re}}\providecommand{\ieme}{\textsuperscript{e}}\providecommand{\ere}{\textsuperscript{re}}\providecommand{\eme}{\textsuperscript{e}}\providecommand{\er}{\textsuperscript{er}}\providecommand{\ie}{\textsuperscript{e}}
% Ordinaux français écrits en macro par les modèles : 1\ière, 3\ième
\newcommand{\ière}{\textsuperscript{re}}\newcommand{\ième}{\textsuperscript{e}}
% \exemple (inventée par les modèles) : étiquette « Exemple : »
\providecommand{\exemple}{\textbf{Exemple~:}\ }
% \square utilisable aussi hors mode math (cases à cocher)
\AtBeginDocument{\let\squareMath\square\renewcommand{\square}{\ensuremath{\squareMath}}}
% Mot ou phrase en chinois à l'intérieur d'un paragraphe français
\newcommand{\zh}[1]{\ifmmode\text{#1}\else{#1}\fi}
\let\eng\zh \let\esp\zh \let\all\zh % alias tolérés
% flèches d'intonation inventées par les modèles
\providecommand{\textsearrow}{}\renewcommand{\textsearrow}{\ensuremath{\searrow}}\providecommand{\textnearrow}{}\renewcommand{\textnearrow}{\ensuremath{\nearrow}}
% \fr{..} : mot français (inventé par les modèles)
\providecommand{\fr}[1]{\textit{#1}}
% zhbox : encadré de texte chinois inventé par les modèles
\newenvironment{zhbox}{\begin{quote}}{\end{quote}}
% \vs : « versus » inventé par les modèles
\providecommand{\vs}{\textit{vs}\ }
% \blank : case à compléter inventée par les modèles
\providecommand{\blank}[1][]{\underline{\hspace{1.6em}}}
\newtcolorbox{exemplebox}[1][]{popbase,colframe=poporange,before upper={\popboxhead{Exemple}{poporange}{#1}}}
\newtcolorbox{definition}[1][]{popbase,colframe=popblue,before upper={\popboxhead{Définition}{popblue}{#1}}}
\newtcolorbox{propriete}[1][]{popbase,colframe=poppurple,before upper={\popboxhead{Propriété}{poppurple}{#1}}}
\newtcolorbox{methode}[1][]{popbase,colframe=popgold,before upper={\popboxhead{Méthode}{popgold}{#1}}}
\newtcolorbox{attention}[1][]{popbase,colframe=poppink,before upper={\popboxhead{Attention}{poppink}{#1}}}
\newtcolorbox{experience}[1][]{popbase,colframe=popblue,colback=popblueL,before upper={\popboxhead{Expérience}{popblue}{#1}}}
% « Le savais-tu ? » : carte violette pleine (contexte, culture, Cameroun)
\newtcolorbox{savaistu}[1][]{popbase,colback=poppurple,colframe=popink,
  fontupper=\popsemi\fontsize{10.4}{14.2}\selectfont\color{white},
  before upper={\poppill{Le savais-tu\,?}{white}{poppurple}\ifx\relax#1\relax\else\hspace{6pt}{\popxbold\color{white}#1}\fi\par\vspace{7pt}}}
% Exercice résolu : énoncé en haut, \tcblower puis la solution sur fond crème
\newcounter{popexo}\newcounter{popexr}
\newtcolorbox{exoresolu}[1][]{popbase,colframe=poporange,colbacklower=popcream,
  segmentation style={popink,line width=1.2pt,dashed},
  before upper={\stepcounter{popexr}\popboxhead{Exercice résolu \thepopexr}{poporange}{#1}},
  before lower={\poppill{Solution}{popink}{popcream}\par\vspace{6pt}}}
% Exercice (non résolu en place) — la correction est dans la partie Corrigés
\newtcolorbox{exercice}[1][]{popbase,colframe=popink,
  before upper={\stepcounter{popexo}\popboxhead{Exercice \thepopexo}{popink}{#1}}}
% Corrigé
\newtcolorbox{corrige}[1][]{popbase,colframe=popgreen,colback=popgreenL,
  before upper={\popboxhead{Corrigé}{popgreen}{#1}}}
% Objectifs / prérequis / bilan
\newtcolorbox{objectifs}{popbase,colframe=popink,colback=poporangeL,
  before upper={\popboxhead{Objectifs de la leçon}{popink}{}}}
\newtcolorbox{prerequis}{popbase,colframe=popink,colback=popblueL,
  before upper={\popboxhead{Prérequis}{popblue}{}}}
\newtcolorbox{bilan}{popbase,colframe=popink,colback=white,boxrule=2.8pt,
  before upper={\popboxhead{Fiche bilan}{poporange}{L'essentiel à connaître pour l'examen}}}
% Figure encadrée avec légende
\newtcolorbox{popfigure}[1][]{enhanced,breakable=false,colback=white,colframe=popline,boxrule=1.4pt,arc=8pt,
  left=10pt,right=10pt,top=10pt,bottom=8pt,before skip=10pt,after skip=12pt,halign=center,
  lower separated=false,halign lower=center,
  fontlower=\popsemi\fontsize{9}{11.5}\selectfont\color{popmuted},
  #1}
\newcommand{\legende}[1]{\par\vspace{6pt}{\popsemi\fontsize{9}{11.5}\selectfont\color{popmuted}#1}}

% ============================================================
% Signature OnBuch+
% ============================================================
\newcommand{\poplogoinline}[1]{{\poplogo\fontsize{#1}{#1}\selectfont\color{popink}OnBuch\color{poporange}+}}
\newcommand{\signatureonbuch}{%
  \begin{tcolorbox}[enhanced,colback=popink,colframe=popink,boxrule=2.2pt,arc=10pt,
      drop shadow=poporange,left=14pt,right=14pt,top=10pt,bottom=10pt,before skip=0pt,after skip=0pt]
    \tikz[baseline=-0.7ex]\node[circle,fill=popgreen,draw=popcream,line width=1.3pt,minimum size=22pt,inner sep=0]{\popcheck{white}};%
    \hspace{9pt}\begin{minipage}[c]{0.62\linewidth}
      {\popxbold\fontsize{12}{13}\selectfont\color{popcream}Signé\ \textbullet\ {\poplogo OnBuch\color{poporange}+}}\par\vspace{2pt}
      {\raggedright\popsemi\fontsize{8}{10.5}\selectfont\color{popline}\MakeUppercase{Équipe pédagogique OnBuch+}\\\MakeUppercase{Conforme au programme officiel MINESEC}\par}
    \end{minipage}\hfill
    {\poplogo\fontsize{22}{22}\selectfont\color{popcream}OnBuch\color{poporange}+}
  \end{tcolorbox}%
}

% ============================================================
% Page de garde
% #1 titre · #2 matière · #3 niveau · #4 module · #5 n° leçon · #6 durée ·
% #7 description / objectifs (texte ou itemize)
% ============================================================
\newcommand{\pagegarde}[7]{%
  \thispagestyle{empty}
  \newgeometry{a4paper,margin=1.7cm,top=1.6cm,bottom=1.4cm}%
  \begin{tikzpicture}[remember picture,overlay]
    \fill[poporange!18] ([xshift=-3.2cm,yshift=-3.6cm]current page.north east) circle (4.6cm);
    \fill[poppurple!14] ([xshift=2.4cm,yshift=7.6cm]current page.south west) circle (3.8cm);
  \end{tikzpicture}%
  {\poplogo\fontsize{30}{30}\selectfont\color{popink}OnBuch\color{poporange}+}\hfill
  \raisebox{0.6ex}{\poppill{Cours signé \textbullet\ Premium}{popink}{popcream}}\par
  \vspace{1pt}\popsquiggle{poppurple}\par
  \vspace{8pt}
  \begin{tcolorbox}[enhanced,colback=poporange,colframe=popink,boxrule=2.8pt,arc=13pt,
      drop shadow=popink,left=18pt,right=18pt,top=12pt,bottom=13pt,after skip=10pt]
    \poppill{#2 \textbullet\ #3}{popink}{popcream}\hspace{5pt}\poppill{#4}{white}{popink}\par\vspace{8pt}
    {\popdisplay\fontsize{14}{14}\selectfont\color{popink}LEÇON #5}\par\vspace{4pt}
    {\raggedright\hyphenpenalty=10000\exhyphenpenalty=10000\popdisplay\fontsize{25}{27}\selectfont\color{popcream}\MakeUppercase{#1}\par}
    \vspace{8pt}
    {\popxbold\fontsize{9.5}{9.5}\selectfont\color{popink}\MakeUppercase{Durée conseillée : #6}}
  \end{tcolorbox}
  \begin{tcolorbox}[enhanced,colback=white,colframe=popink,boxrule=2.2pt,arc=10pt,
      drop shadow=popink,left=16pt,right=16pt,top=10pt,bottom=10pt,after skip=8pt]
    \poppill{Dans cette leçon}{poppurple}{white}\par\vspace{5pt}
    \popsemi\fontsize{9.3}{12.6}\selectfont\color{popink}#7
  \end{tcolorbox}
  \noindent\begin{minipage}[t]{0.487\linewidth}
    \begin{tcolorbox}[enhanced,colback=popgreen,colframe=popink,boxrule=2.2pt,arc=10pt,
        drop shadow=popink,left=11pt,right=11pt,top=10pt,bottom=10pt,height=3.1cm,valign=center]
      \tcbox[colback=white,colframe=popink,boxrule=1.3pt,arc=5pt,boxsep=0pt,nobeforeafter,
        left=3pt,right=3pt,top=3pt,bottom=3pt,valign=center]{\qrcode[height=1.9cm]{https://play.google.com/store/apps/details?id=com.luvvix.onbuch&hl=fr}}%
      \hspace{8pt}\begin{minipage}[c]{0.5\linewidth}
        {\popdisplay\fontsize{11}{12.5}\selectfont\color{white}\MakeUppercase{Télécharge OnBuch}}\par\vspace{4pt}
        {\raggedright\popsemi\fontsize{8.4}{11}\selectfont\color{white}Gratuit \textbullet\ Google Play\\Tuteur IA, annales, concours, résultats\par}
      \end{minipage}
    \end{tcolorbox}
  \end{minipage}\hfill
  \begin{minipage}[t]{0.487\linewidth}
    \begin{tcolorbox}[enhanced,colback=poppurple,colframe=popink,boxrule=2.2pt,arc=10pt,
        drop shadow=popink,left=11pt,right=11pt,top=10pt,bottom=10pt,height=3.1cm,valign=center]
      \tikz[baseline=-0.7ex]\node[circle,fill=white,draw=popink,line width=1.3pt,minimum size=1.6cm,inner sep=0]{\popdisplay\fontsize{24}{24}\selectfont\color{poppurple}+};%
      \hspace{8pt}\begin{minipage}[c]{0.6\linewidth}
        {\popdisplay\fontsize{11}{12.5}\selectfont\color{white}\MakeUppercase{Passe à OnBuch+}}\par\vspace{4pt}
        {\raggedright\popsemi\fontsize{8.4}{11}\selectfont\color{white}Tous les cours signés, Tuteur IA illimité, fiches PDF — onglet OnBuch+ de l'app\par}
      \end{minipage}
    \end{tcolorbox}
  \end{minipage}
  \vspace{8pt}
  \popcontenu
  \vfill
  \signatureonbuch
  \restoregeometry
  \newpage
}

% Rangée « Ce cours contient » (page de garde)
\newcommand{\poptile}[3]{% #1 couleur #2 gros texte #3 légende
  \begin{tcolorbox}[enhanced,colback=#1,colframe=popink,boxrule=2pt,arc=9pt,shadow={2.5pt}{-2.5pt}{0pt}{popink},
      left=6pt,right=6pt,top=9pt,bottom=9pt,halign=center,valign=center,height=1.75cm,width=\linewidth,nobeforeafter]
    {\popdisplay\fontsize{12.5}{13}\selectfont\color{white}\MakeUppercase{#2}}\par\vspace{3pt}
    {\centering\popsemi\fontsize{7.8}{9.5}\selectfont\color{white}#3\par}
  \end{tcolorbox}}
\newcommand{\popcontenu}{%
  \par\noindent{\popxboldls\fontsize{7.5}{7.5}\selectfont\color{popmuted}CE COURS SIGNÉ CONTIENT}\par\vspace{7pt}
  \noindent\begin{minipage}[t]{0.235\linewidth}\poptile{popblue}{Cours}{complet, illustré, pas à pas}\end{minipage}\hfill
  \begin{minipage}[t]{0.235\linewidth}\poptile{poporange}{Méthodes}{et exercices résolus}\end{minipage}\hfill
  \begin{minipage}[t]{0.235\linewidth}\poptile{poppink}{Exercices}{type examen + corrigés}\end{minipage}\hfill
  \begin{minipage}[t]{0.235\linewidth}\poptile{popgreen}{Fiche bilan}{l'essentiel à retenir}\end{minipage}\par}

% Sommaire
\newtcolorbox{sommaire}{enhanced,colback=white,colframe=popink,boxrule=2.2pt,arc=10pt,
  drop shadow=popink,left=18pt,right=18pt,top=13pt,bottom=13pt,before skip=6pt,after skip=20pt,
  before upper={\poppill{Sommaire}{poppurple}{white}\par\vspace{9pt}\popsemi\fontsize{10.6}{14.5}\selectfont\color{popink}}}

% Signature de fin de cours — poussée tout en bas de la dernière page
\newcommand{\findecours}{%
  \par\vfill\nopagebreak
  \begin{center}\popsquiggle{poporange}\end{center}\vspace{4pt}
  \signatureonbuch
}
\AtBeginDocument{\def\llbracket{\mathopen{[\mkern-3.5mu[}}\def\rrbracket{\mathclose{]\mkern-3.5mu]}}} % intervalles d'entiers (glyphe ⟦ absent de la police)
% Glyphes absents de Fira Math : redéfinis à partir de glyphes disponibles
\AtBeginDocument{%
  \def\setminus{\mathbin{\backslash}}\let\smallsetminus\setminus
  \def\vdots{\mathord{\vbox{\baselineskip3.2pt\lineskiplimit0pt\kern4pt\hbox{.}\hbox{.}\hbox{.}}}}%
  \def\ddots{\mathinner{\mkern1mu\raise6.5pt\hbox{.}\mkern2mu\raise3.6pt\hbox{.}\mkern2mu\raise0.7pt\hbox{.}\mkern1mu}}%
  \def\triangle{\mathord{\scalebox{0.9}{$\symup{\Delta}$}}}%
  \def\nabla{\mathord{\raisebox{\depth}{\scalebox{1}[-1]{$\symup{\Delta}$}}}}%
  \def\checkmark{\popcheck{popdark}}%
  \def\star{\mathbin{\ast}}\def\bigstar{\mathord{\ast}}%
  \def\diamond{\mathbin{\scalebox{0.6}{$\Diamond$}}}%
  \def\bigcup{\mathop{\mathchoice{\raisebox{-0.3em}{\scalebox{1.7}{$\cup$}}}{\scalebox{1.25}{$\cup$}}{\cup}{\cup}}\displaylimits}%
  \def\bigcap{\mathop{\mathchoice{\raisebox{-0.3em}{\scalebox{1.7}{$\cap$}}}{\scalebox{1.25}{$\cap$}}{\cap}{\cap}}\displaylimits}%
  \def\lesseqqgtr{\mathrel{\lessgtr}}\def\bigoplus{\mathop{\oplus}\displaylimits}%
}
\providecommand{\ssi}{si et seulement si} % abréviation fréquente
\AtBeginDocument{\providecommand{\wideparen}{\overparen}} % arc de cercle

%% FIGFIX-BEGIN v3 — correctifs globaux des figures (docs/onbuch-generation/tools/figfix)
\usepackage{adjustbox}\usepackage{etoolbox}
% toute figure TikZ/pgfplots plus large que la ligne est réduite à la largeur disponible
\BeforeBeginEnvironment{tikzpicture}{\begin{adjustbox}{max width=\linewidth}}
\AfterEndEnvironment{tikzpicture}{\end{adjustbox}}
% légendes pgfplots lisibles par-dessus les courbes
\pgfplotsset{every axis/.append style={legend style={fill=white,fill opacity=0.92,text opacity=1,draw=black!15,inner sep=3pt}}}
% étiquettes d'arêtes (syntaxe edge["..."]) avec fond blanc
\tikzset{every edge quotes/.append style={fill=white,inner sep=1.2pt}}
% bulles de cartes mentales : texte à la ligne dans la bulle, bulle un peu plus grande
\tikzset{every concept/.append style={text width=2.0cm,align=center,minimum size=2.3cm,inner sep=2pt}}
%% FIGFIX-END
\begin{document}

\pagegarde{Leçon 17 — Vocabulaire et compréhension de texte : collectivités locales}{Chinois}{Tle A}{Module 3 : Citoyenneté et ouverture au monde}{ZH17}{2 h}{Cette leçon fait découvrir le vocabulaire chinois des collectivités locales (ville, commune, quartier, village, habitants, services) à travers un texte support original. L'élève apprend à lire, comprendre et réutiliser ce lexique pour décrire sa propre localité et répondre à des questions de compréhension type examen.
\vspace{4pt}
\begin{multicols}{2}\small
\begin{itemize}
\item À la fin de cette leçon, je sais lire à voix haute un court texte chinois sur les collectivités locales en respectant le pinyin et les tons.
\item À la fin de cette leçon, je sais reconnaître et traduire au moins 15 mots du thème (ville, village, quartier, habitant, maire, service...).
\item À la fin de cette leçon, je sais identifier les clés (radicaux) et tracer correctement l'ordre des traits de caractères comme 城, 区, 村, 民.
\item À la fin de cette leçon, je sais répondre par écrit à des questions de compréhension portant sur le texte support.
\item À la fin de cette leçon, je sais employer les collocations usuelles du thème (住在..., 市政府, 居委会...) dans des phrases complètes.
\item À la fin de cette leçon, je sais décrire ma commune ou mon quartier camerounais en 4 à 6 phrases chinoises.
\end{itemize}\end{multicols}}

\begin{sommaire}
\begin{multicols}{2}
\begin{enumerate}
\item Mise en route et découverte du thème
\item Texte support : lecture et compréhension globale
\item Glossaire du thème en tableau
\item Clés (radicaux) et ordre des traits
\item Questions de compréhension et collocations
\item Production guidée : ma collectivité locale
\item Activité d'intégration
\item Méthodes et exercices résolus
\item Exercices
\item Corrigés des exercices
\item Fiche bilan
\end{enumerate}
\end{multicols}
\end{sommaire}

\begin{objectifs}
\begin{itemize}
\item À la fin de cette leçon, je sais lire à voix haute un court texte chinois sur les collectivités locales en respectant le pinyin et les tons.
\item À la fin de cette leçon, je sais reconnaître et traduire au moins 15 mots du thème (ville, village, quartier, habitant, maire, service...).
\item À la fin de cette leçon, je sais identifier les clés (radicaux) et tracer correctement l'ordre des traits de caractères comme 城, 区, 村, 民.
\item À la fin de cette leçon, je sais répondre par écrit à des questions de compréhension portant sur le texte support.
\item À la fin de cette leçon, je sais employer les collocations usuelles du thème (住在..., 市政府, 居委会...) dans des phrases complètes.
\item À la fin de cette leçon, je sais décrire ma commune ou mon quartier camerounais en 4 à 6 phrases chinoises.
\item À la fin de cette leçon, je sais comparer brièvement l'organisation locale du Cameroun et de la Chine.
\end{itemize}
\end{objectifs}

\begin{prerequis}
\begin{itemize}
\item Vocabulaire des lieux et de la ville vu en Première (在, 有, 去, 学校, 家...)
\item Structure de base sujet-verbe-complément et verbe 在 pour la localisation
\item Verbe 有 pour exprimer l'existence et la possession
\item Notions de radicaux et d'ordre des traits acquises depuis la Seconde
\item Adjectifs simples de description (大, 小, 多, 少, 漂亮)
\end{itemize}
\end{prerequis}

\newpage

\coursec{Mise en route et découverte du thème}

\courssub{Situation d'accroche : le maire visite ton lycée}

Imagine la scène suivante. Tu es délégué de classe au lycée de Yaoundé, de Douala ou de Bafoussam. Ce matin, le maire de ta commune visite ton établissement, accompagné d'une délégation de correspondants chinois venus d'une ville jumelée. Le proviseur te désigne : « À toi de présenter notre quartier à nos invités ! » Comment ferais-tu ? Comment dirais-tu, en chinois, « mon quartier », « ma commune », « mon village », « les habitants », « le gouvernement local » ?

En Chine comme au Cameroun, la vie quotidienne s'organise autour des \cle{collectivités locales} : la commune, la communauté urbaine, le quartier, le village, la chefferie, le comité de résidents. Savoir nommer ces réalités en chinois, c'est pouvoir parler de chez soi, comparer les administrations camerounaises et chinoises, et comprendre un texte authentique sur la vie locale — exactement le type de texte qui tombe à l'épreuve de chinois du Baccalauréat. C'est ce que nous allons apprendre aujourd'hui.

\textbf{Problématique de la leçon :} comment lire et comprendre un texte chinois qui décrit une collectivité locale, et quel vocabulaire faut-il maîtriser pour parler de sa propre commune, de son quartier ou de son village ?

\courssub{Brainstorming : qu'est-ce qu'une collectivité locale ?}

Avant d'ouvrir le chinois, mobilisons ce que tu sais déjà en français. Une \cle{collectivité locale} est un territoire organisé, avec une population et une autorité chargée de gérer la vie quotidienne : l'eau, les marchés, les écoles, les routes, la propreté, l'état civil.

Au Cameroun, tu connais déjà plusieurs niveaux :

\begin{itemize}
\item la \textbf{commune} (\zh{市} shì / \zh{区} qū, ex. : la commune de Yaoundé III, la commune d'arrondissement de Douala IV) ;
\item la \textbf{communauté urbaine} (\zh{城市} chéngshì, ex. : la Communauté urbaine de Douala, qui regroupe plusieurs communes) ;
\item le \textbf{village} (\zh{村} cūn), souvent dirigé par un \textbf{chef traditionnel} (\zh{酋长} qiúzhǎng, la chefferie \zh{酋长领地} qiúzhǎng lǐngdì) ;
\item le \textbf{quartier} (\zh{社区} shèqū / \zh{小区} xiǎoqū), avec son chef de quartier, ses réunions de jeunesse, son marché.
\end{itemize}

En chinois, chacune de ces réalités a son mot. Observons d'abord quelques correspondances avant de les organiser :

\begin{center}
\begin{tabularx}{\linewidth}{|X|X|X|X|}
\hline
\rowcolor{popblueL} Caractères & Pinyin & Nature & Traduction \\
\hline
地方 & dìfang & nom & le lieu, la localité, l'endroit \\
\hline
城市 & chéngshì & nom & la ville \\
\hline
农村 & nóngcūn & nom & la campagne, le monde rural \\
\hline
社区 & shèqū & nom & la communauté, le quartier résidentiel \\
\hline
政府 & zhèngfǔ & nom & le gouvernement, l'administration \\
\hline
居民 & jūmín & nom & l'habitant, le résident \\
\hline
\end{tabularx}
\end{center}

Observe bien : le caractère \zh{地} (dì, la terre) revient dans \zh{地方} (dìfang) et \zh{地点} (dìdiǎn) ; le caractère \zh{民} (mín, le peuple) revient dans \zh{居民} (jūmín, habitant) et \zh{人民} (rénmín, le peuple). Le chinois construit son vocabulaire par \cle{combinaison de caractères} : comprendre les caractères, c'est deviner le sens des mots nouveaux.

\begin{propriete}[Clés (radicaux) et ordre des traits — caractères du thème]
\begin{itemize}
\item \textbf{\zh{地} (dì) — clé \zh{土} (tǔ, « terre »), 6 traits.} Ordre : horizontale (一), verticale (丨), horizontale (一), point (丶), horizontale (一), trait final oblique (㇏). Idée : la terre (\zh{土}) + aussi/phonétique (\zh{也}) → le sol, le lieu.
\item \textbf{\zh{民} (mín) — clé \zh{氏} (shì, « clan »), 5 traits.} Ordre : horizontale (一), verticale (丨), horizontale (一), horizontale (一), trait final oblique (㇏). Idée : l'œil (\zh{目} anciennement) percé → le peuple qui travaille la terre.
\item \textbf{\zh{居} (jū) — clé \zh{尸} (shī, « corps/maison »), 8 traits.} Ordre : horizontale (一), pliée (㇆), horizontale (一), point (丶), horizontale (一), verticale (丨), horizontale (一), horizontale (一). Idée : une maison (\zh{尸}) + ancien \zh{古} (gǔ) → habiter, résider.
\item \textbf{\zh{委} (wěi) — clé \zh{禾} (hé, « grain/riz »), 8 traits.} Ordre : oblique (丿), horizontale (一), verticale (丨), oblique (丿), point (丶), horizontale (一), verticale (丨), horizontale (一). Idée : confier une tâche (comme confier la rizière).
\end{itemize}
\end{propriete}

\courssub{Carte mentale au tableau : autour de 地方}

Construisons ensemble la carte mentale du thème. Au centre, le mot-clé \zh{地方} (dìfang, le lieu, la localité) ; autour, les cinq grandes familles de mots que nous allons rencontrer dans le texte du jour.

\begin{popfigure}
\begin{tikzpicture}[mindmap, grow cyclic, every node/.style={concept, align=center},
root concept/.append style={concept color=popblue, font=\Large},
level 1/.append style={level distance=3.4cm, sibling angle=72},
level 1 concept/.append style={concept color=poporange, font=\small}]
\node[root concept, align=center]{地方\\ dìfang\\ \emph{le lieu}}
child { node[align=center]{城市\\ chéngshì\\ \emph{la ville}} }
child { node[align=center]{农村\\ nóngcūn\\ \emph{la campagne}} }
child { node[align=center]{社区\\ shèqū\\ \emph{la communauté}} }
child { node[align=center]{政府\\ zhèngfǔ\\ \emph{le gouvernement}} }
child { node[align=center]{居民\\ jūmín\\ \emph{les habitants}} };
\end{tikzpicture}
\legende{Carte mentale lexicale : le champ des collectivités locales autour de 地方 (dìfang).}
\end{popfigure}

Recopie cette carte mentale dans ton cahier : elle te servira de plan pour mémoriser le vocabulaire et pour ta production écrite « ma collectivité locale » à la fin de la leçon.

\begin{savaistu}[Le 居委会, le « chef de quartier » chinois]
En Chine, chaque quartier urbain est géré par un \zh{居委会} (jūwěihuì), le \emph{comité de résidents} : \zh{居} (habiter) + \zh{委} (comité/délégué) + \zh{会} (association/réunion). Ce comité s'occupe des petites affaires du quotidien : propreté, médiation entre voisins, informations aux habitants, aide aux personnes âgées. C'est un peu l'équivalent de nos \textbf{chefs de quartier} au Cameroun, qui connaissent tout le monde et règlent les problèmes de proximité. À la campagne chinoise, on trouve de même le \zh{村委会} (cūnwěihuì), le comité du village — comparable à nos chefferies traditionnelles (\zh{酋长} qiúzhǎng).
\end{savaistu}

\courssub{Objectifs de la leçon et annonce du texte support}

À la fin de cette leçon, tu seras capable de :

\begin{enumerate}
\item lire et comprendre globalement un court texte chinois sur une collectivité locale ;
\item reconnaître, prononcer (pinyin et tons) et employer le vocabulaire du thème ;
\item répondre à des questions de compréhension, comme à l'examen.
\end{enumerate}

Le texte support du jour s'intitule \zh{我们的社区} (Wǒmen de shèqū, « Notre communauté ») : un habitant décrit son quartier, sa ville et l'action du gouvernement local. Mais avant de le lire en détail, apprenons la bonne démarche.

\begin{methode}[Aborder un texte chinois en trois étapes]
\begin{enumerate}
\item \textbf{Lire le titre.} Le titre annonce le thème : \zh{我们的社区} contient \zh{社区} (communauté) — tu sais déjà de quoi on va parler.
\item \textbf{Repérer les mots connus.} Souligne mentalement les mots que tu as déjà appris (\zh{我们}, \zh{的}, \zh{城市}…). Ils forment le squelette du sens.
\item \textbf{Deviner le sens global avant de traduire mot à mot.} Pose-toi la question : « De quoi parle-t-on ? » La traduction précise vient ensuite, jamais avant.
\end{enumerate}
Erreur à bannir : traduire caractère par caractère dès la première lecture. On lit d'abord pour \emph{comprendre}, ensuite pour \emph{traduire}.
\end{methode}

\courssub{Première écoute : repérage global}

Ferme ton cahier de vocabulaire. L'enseignant lit le texte une première fois, à vitesse normale. Ta seule mission : répondre à la question « \emph{De quoi parle-t-on ?} ». Écoute le texte support complet (il sera étudié en détail dans la section suivante) :

\begin{textebox}[我们的社区 — première écoute]
我住在一个很大的社区。这里有学校、医院和市场，生活很方便。\\
Wǒ zhù zài yí ge hěn dà de shèqū. Zhèlǐ yǒu xuéxiào, yīyuàn hé shìchǎng, shēnghuó hěn fāngbiàn.\\
« J'habite dans une très grande communauté. Il y a ici une école, un hôpital et un marché ; la vie est très pratique. »\\
我们的城市很美丽，居民都很友好。地方政府也很关心居民的生活。\\
Wǒmen de chéngshì hěn měilì, jūmín dōu hěn yǒuhǎo. Dìfāng zhèngfǔ yě hěn guānxīn jūmín de shēnghuó.\\
« Notre ville est très belle, les habitants sont tous très amicaux. Le gouvernement local se préoccupe aussi beaucoup de la vie des habitants. »
\end{textebox}

Après cette première écoute, tu dois pouvoir dire, en français : le narrateur parle de \textbf{sa communauté} (son quartier), de ses \textbf{équipements} (école, hôpital, marché), de sa \textbf{ville} et du \textbf{gouvernement local}. Si tu as saisi cela, ta compréhension globale est réussie — les détails viendront à la deuxième lecture.

\begin{attention}[地方 ou 地点 ?]
Ne confonds pas \zh{地方} (dìfang, le lieu, la localité, sens général) et \zh{地点} (dìdiǎn, l'endroit précis, le site exact d'un événement). Compare : \zh{这个地方很美} (Zhège dìfang hěn měi, « cet endroit est beau » — sens général, localité) et \zh{集合的地点} (jíhé de dìdiǎn, « le lieu exact du rassemblement » — point GPS, adresse précise). Au Bac, c'est le \cle{contexte} qui tranche : si l'on parle d'un point précis de rendez-vous, c'est \zh{地点} ; si l'on parle d'une localité, d'une région, d'un cadre de vie, c'est \zh{地方}.
\end{attention}

\begin{exoresolu}[Premier repérage dans un texte]
Énoncé : voici une phrase extraite du texte du jour : \zh{地方政府也很关心居民的生活。} Sans dictionnaire, (a) souligne les mots que tu connais déjà grâce à la carte mentale ; (b) devine le sens global de la phrase.
\tcblower
Solution : (a) Mots connus repérés : \zh{地方} (local), \zh{政府} (gouvernement), \zh{居民} (habitants), \zh{生活} (vie, connu depuis le module « vie quotidienne »). (b) Sens global deviné : « Le gouvernement local (地方政府) se préoccupe aussi (也很关心) de la vie des habitants (居民的生活). » Traduction complète : « Le gouvernement local se préoccupe aussi beaucoup de la vie des habitants. » Remarque la méthode : quatre mots connus suffisent à comprendre l'essentiel de la phrase avant toute traduction mot à mot.
\end{exoresolu}

\begin{aretenir}[Ce qu'il faut retenir de cette mise en route]
\begin{itemize}
\item Une collectivité locale (commune, quartier, village, chefferie) se dit en chinois avec la famille de mots autour de \zh{地方} : \zh{城市}, \zh{农村}, \zh{社区}, \zh{政府}, \zh{居民}.
\item En Chine, le quartier est géré par le \zh{居委会} (comité de résidents), proche de nos chefs de quartier ; le village par le \zh{村委会}.
\item Pour lire un texte chinois : titre d'abord, mots connus ensuite, sens global avant la traduction.
\item Distinction cruciale : \zh{地方} (localité, zone) $\neq$ \zh{地点} (point précis, coordonnées).
\end{itemize}
\end{aretenir}

\coursec{Mise en route et découverte du thème}

Avant d'aborder le texte, activons nos connaissances sur l'organisation du territoire au Cameroun et en Chine.

\begin{experience}[Observer et comparer : mon quartier]
\begin{enumerate}
\item \textbf{Au Cameroun.} Citez le nom de votre quartier ou arrondissement (ex. : \emph{Bastos, Mvog-Ada, Akwa, Bepanda, Molyko}). Quels services publics y trouvez-vous ? (École, centre de santé, marché, terrain de football, mairie d'arrondissement). Notez 5 mots en français.
\item \textbf{En Chine.} L'unité de base urbaine est le \zh{社区} (shèqū), géré par un \zh{居委会} (jūwěihuì). Regardez l'image ci-dessous (panneau d'affichage typique d'un \zh{社区} à Pékin ou Shanghai) et repérez les caractères que vous connaissez déjà.
\end{enumerate}
\end{experience}

\begin{popfigure}
\begin{center}
\begin{tikzpicture}[
node distance=0.6cm,
box/.style={rectangle, rounded corners, draw=popblue, fill=popblue!10, thick, align=center, minimum width=4cm, minimum height=1.2cm, font=\small},
arrow/.style={-{Stealth[length=2.5mm]}, thick, draw=popdark}
]
\node[box, align=center] (comm){\textbf{Commune / Arrondissement}\\(Cameroun)\\\zh{街道 / 区}};
\node[box, right=of comm, align=center] (jwh){\textbf{Comité de quartier}\\\zh{居委会 (jūwěihuì)}};
\node[box, right=of jwh, align=center] (sheq){\textbf{Quartier / Communauté}\\\zh{社区 (shèqū)}};
\node[box, below=of jwh, align=center] (ju){\textbf{Habitants}\\\zh{居民 (jūmín)}};
\draw[arrow] (comm) -- (jwh);
\draw[arrow] (jwh) -- (sheq);
\draw[arrow] (sheq) -- (ju);
\draw[arrow] (ju) -- (sheq);
\end{tikzpicture}
\end{center}
\legende{Organisation administrative locale : comparaison Cameroun / Chine (simplifiée).}
\end{popfigure}

\begin{methode}[Découvrir le vocabulaire par l'image et le contexte]
\begin{enumerate}
\item Associez chaque service à son caractère chinois : \zh{学校} (école), \zh{医院} (hôpital), \zh{商店} (magasin), \zh{公园} (parc), \zh{市场} (marché).
\item Devinez le sens de \zh{居民} (indice : \zh{居} = habiter, \zh{民} = peuple/population) et \zh{居委会} (indice : \zh{委} = confier/mandater, \zh{会} = réunion/association).
\item Échangez avec votre voisin : « Dans mon quartier, il y a \ldots » / « \zh{我社区里有\ldots} ».
\end{enumerate}
\end{methode}

\coursec{Texte support : lecture et compréhension globale}

Dans la section précédente, nous avons découvert le thème des collectivités locales à travers des images et des mots déjà connus : \zh{社区} (shèqū, quartier/communauté), \zh{居民} (jūmín, habitant), \zh{医院} (yīyuàn, hôpital). Il est temps maintenant de rencontrer ces mots \emph{en situation}, c'est-à-dire dans un texte authentique. Lire un texte chinois, ce n'est pas déchiffrer caractère par caractère : c'est suivre une démarche en trois étapes que nous allons apprendre et appliquer ensemble.

\courssub{Le texte support : 《我住的社区》}

Voici le texte de la leçon. Il s'agit du récit d'un habitant qui présente son quartier. Ne cherchez pas encore à tout comprendre : lisez d'abord les caractères, puis le pinyin, enfin la traduction.

\begin{textebox}[《我住的社区》 — Wǒ zhù de shèqū — Le quartier où j'habite]
我住在一个很大的社区。社区里有学校、医院、商店和公园。\\
Wǒ zhù zài yí ge hěn dà de shèqū. Shèqū lǐ yǒu xuéxiào, yīyuàn, shāngdiàn hé gōngyuán.\\
居民们都很友好。居委会经常组织活动，大家都参加。\\
Jūmínmen dōu hěn yǒuhǎo. Jūwěihuì jīngcháng zǔzhī huódòng, dàjiā dōu cānjiā.\\
我喜欢我的社区。\\
Wǒ xǐhuan wǒ de shèqū.\\[0.5em]
\emph{Traduction :} J'habite dans un très grand quartier. Dans le quartier, il y a une école, un hôpital, des magasins et un parc. Les habitants sont tous très amicaux. Le comité de résidents organise souvent des activités, et tout le monde y participe. J'aime mon quartier.
\end{textebox}

\textbf{Version longue pour l'entraînement à la lecture (texte complémentaire).}\par

Pour les élèves rapides ou pour un second passage en classe, voici une version enrichie du même récit (environ 120 caractères), à lire avec la même méthode :

\begin{textebox}[《我们的社区》 — Wǒmen de shèqū — Notre quartier (texte d'approfondissement)]
我们的社区在市中心，交通很方便。\\
Wǒmen de shèqū zài shì zhōngxīn, jiāotōng hěn fāngbiàn.\\
社区里有一所小学、一家医院和一个大市场。\\
Shèqū lǐ yǒu yì suǒ xiǎoxué, yì jiā yīyuàn hé yí ge dà shìchǎng.\\
每天早上，很多居民在公园里锻炼身体。\\
Měi tiān zǎoshang, hěn duō jūmín zài gōngyuán lǐ duànliàn shēntǐ.\\
居委会的工作人员很热情，常常帮助老人和孩子。\\
Jūwěihuì de gōngzuò rényuán hěn rèqíng, chángcháng bāngzhù lǎorén hé háizi.\\
上个月，社区组织了一次清洁活动，两百多个居民参加了。\\
Shàng ge yuè, shèqū zǔzhī le yí cì qīngjié huódòng, liǎng bǎi duō ge jūmín cānjiā le.\\
大家都说，住在这里很幸福。\\
Dàjiā dōu shuō, zhù zài zhèlǐ hěn xìngfú.\\[0.5em]
\emph{Traduction :} Notre quartier est au centre-ville, les transports sont très pratiques. Dans le quartier, il y a une école primaire, un hôpital et un grand marché. Chaque matin, beaucoup d'habitants font de l'exercice dans le parc. Les employés du comité de résidents sont très chaleureux et aident souvent les personnes âgées et les enfants. Le mois dernier, le quartier a organisé une activité de nettoyage, à laquelle plus de deux cents habitants ont participé. Tout le monde dit qu'on est heureux de vivre ici.
\end{textebox}

\courssub{La lecture en trois étapes}

Lire un texte chinois efficacement demande une méthode. Voici celle que nous utiliserons toute l'année.

\begin{methode}[Lire un texte chinois en trois étapes]
\begin{enumerate}
\item \textbf{Écoute silencieuse.} Le professeur lit le texte à voix haute (ou on écoute l'enregistrement). Les élèves \emph{suivent des yeux} les caractères sans lire. Objectif : saisir la mélodie des tons, repérer les pauses (，et 。) et identifier les mots déjà connus.
\item \textbf{Lecture chorale avec le pinyin.} Toute la classe lit ensemble, phrase par phrase, en s'aidant du pinyin. Le professeur corrige immédiatement les tons. On répète chaque phrase deux fois : une fois lentement, une fois à vitesse normale.
\item \textbf{Lecture individuelle.} Chaque élève lit une phrase à voix haute, sans le pinyin si possible. Les autres vérifient. C'est à cette étape que l'on contrôle la compréhension globale : qui parle ? de quoi parle-t-il ?
\end{enumerate}
\end{methode}

\begin{attention}[Les tons avant la vitesse]
L'erreur classique des francophones est de lire vite en « aplatissant » les tons. Or \zh{mā} (maman), \zh{má} (chanvre), \zh{mǎ} (cheval) et \zh{mà} (gronder) sont quatre mots différents. Dans notre texte, faites la différence entre \zh{shèqū} (4e ton + 1er ton) et \zh{jūmín} (1er ton + 2e ton). Mieux vaut lire lentement et juste que vite et faux.
\end{attention}

\courssub{Repérage des mots du thème dans le texte}

Avant de traduire, on \emph{surligne} dans le texte tous les mots qui appartiennent au champ lexical des collectivités locales. C'est un réflexe de lecteur : les mots du thème portent le sens principal.

\begin{exemplebox}[Les mots du thème surlignés dans le texte]
Dans 《我住的社区》, on surligne : \cle{社区} (shèqū, quartier), \cle{学校} (xuéxiào, école), \cle{医院} (yīyuàn, hôpital), \cle{商店} (shāngdiàn, magasin), \cle{公园} (gōngyuán, parc), \cle{居民} (jūmín, habitant), \cle{居委会} (jūwěihuì, comité de résidents), \cle{活动} (huódòng, activité).\\
Ces huit mots suffisent à deviner le sujet du texte : \emph{la vie dans un quartier et ses services}.
\end{exemplebox}

\coursec{Glossaire du thème en tableau}

\begin{center}
\begin{tabularx}{\linewidth}{|X|X|X|X|}
\hline
\rowcolor{popblueL} Caractères & Pinyin & Nature & Traduction \\
\hline
社区 & shèqū & nom & quartier, communauté \\
\hline
居民 & jūmín & nom & habitant, résident \\
\hline
居委会 & jūwěihuì & nom & comité de résidents \\
\hline
商店 & shāngdiàn & nom & magasin, boutique \\
\hline
公园 & gōngyuán & nom & parc public \\
\hline
市场 & shìchǎng & nom & marché \\
\hline
小学 & xiǎoxué & nom & école primaire \\
\hline
市中心 & shì zhōngxīn & nom & centre-ville \\
\hline
交通 & jiāotōng & nom & transport, circulation \\
\hline
组织 & zǔzhī & verbe & organiser \\
\hline
活动 & huódòng & nom & activité \\
\hline
参加 & cānjiā & verbe & participer à \\
\hline
友好 & yǒuhǎo & adjectif & amical, sympathique \\
\hline
热情 & rèqíng & adjectif & chaleureux, enthousiaste \\
\hline
锻炼 & duànliàn & verbe & faire de l'exercice, s'entraîner \\
\hline
清洁 & qīngjié & adjectif/verbe & propre / nettoyer \\
\hline
幸福 & xìngfú & adjectif & heureux, bonheur \\
\hline
\end{tabularx}
\end{center}

\coursec{Clés (radicaux) et ordre des traits}

L'écriture des caractères suit une logique interne : le \textbf{radical (clé)} donne souvent un indice sémantique, l'\textbf{ordre des traits} assure la lisibilité et la vitesse. Règle générale : haut $\to$ bas, gauche $\to$ droite, horizontal avant vertical, trait central avant les ailes, cadre avant contenu, fermeture en dernier.

\begin{center}
\begin{tabularx}{\linewidth}{|X|X|X|X|}
\hline
\rowcolor{popblueL} Caractère & Radical (Clé) & Nb traits & Ordre des traits (description) \\
\hline
\zh{社} & \zh{示} (shì, autel/rite) & 7 & 1. \zh{示} (点, 横, 撇, 竖, 横) 2. \zh{土} (横, 竖) \\
\hline
\zh{区} & \zh{匚} (fāng, boîte) & 4 & 1. 左竖 2. 横折 (haut + droite) 3. 中竖 4. 底横 \\
\hline
\zh{居} & \zh{尸} (shī, toit/maison) & 8 & 1. \zh{尸} (横折, 横, 撇) 2. \zh{古} (十, 口) \\
\hline
\zh{委} & \zh{禾} (hé, grain/riz) & 8 & 1. \zh{禾} (撇, 横, 竖, 撇, 点) 2. \zh{委} partie droite (撇, 竖弯钩) \\
\hline
\zh{市} & \zh{巾} (jīn, tissu) & 5 & 1. \zh{巾} (竖, 横折钩, 竖, 横, 横) — \emph{attention : 5 traits seulement} \\
\hline
\zh{中} & \zh{丨} (gǔn, trait vertical) & 4 & 1. 竖 2. 横折 3. 横 4. 中竖 (traverse le milieu) \\
\hline
\zh{心} & \zh{心} (xīn, cœur) & 4 & 1. 点 (gauche) 2. 点 (droite) 3. 竖 (milieu) 4. 捺 (long trait final) \\
\hline
\zh{活} & \zh{氵} (shuǐ, eau) & 9 & 1. \zh{氵} (点, 点, 提) 2. \zh{舌} (千 + 口 : 横, 竖, 横折, 横, 横, 竖, 横) \\
\hline
\zh{动} & \zh{力} (lì, force) & 6 & 1. \zh{云} (横, 撇, 横折钩, 点) 2. \zh{力} (横折钩, 点) \\
\hline
\end{tabularx}
\end{center}

\begin{experience}[Atelier écriture : « Mon panneau de quartier »]
Sur une feuille A5, écrivez en grand les caractères \zh{社区服务中心} (Centre de services du quartier). Respectez l'ordre des traits. Affichez les productions en classe : votez pour l'écriture la plus équilibrée (structure, proportion, trait).
\end{experience}

\coursec{Questions de compréhension et collocations}

\courssub{Compréhension globale : trois questions essentielles}

Après la lecture, on répond d'abord aux questions \emph{globales}, sans entrer dans les détails.

\begin{enumerate}
\item \textbf{Qui parle ?} Un habitant du quartier (\zh{我}, « je »). C'est un récit à la première personne.
\item \textbf{Où habite-t-il ?} Dans un très grand quartier : \zh{我住在一个很大的社区。}
\item \textbf{Quels services existent dans son quartier ?} Une école, un hôpital, des magasins et un parc : \zh{社区里有学校、医院、商店和公园。} On peut ajouter un « service » humain : le comité de résidents (\zh{居委会}), qui organise des activités.
\end{enumerate}

\begin{exoresolu}[Questions de compréhension globale]
Énoncé. Répondez en chinois, avec une phrase complète : (1) 他住在什么地方？(Tā zhù zài shénme dìfang ? — Où habite-t-il ?) (2) 社区里有什么？(Shèqū lǐ yǒu shénme ? — Qu'y a-t-il dans le quartier ?) (3) 居民们怎么样？(Jūmínmen zěnmeyàng ? — Comment sont les habitants ?)
\tcblower
Solution détaillée. (1) \zh{他住在一个很大的社区。} (Tā zhù zài yí ge hěn dà de shèqū. — Il habite dans un très grand quartier.) On reprend la structure \zh{住在} + lieu. (2) \zh{社区里有学校、医院、商店和公园。} (Shèqū lǐ yǒu xuéxiào, yīyuàn, shāngdiàn hé gōngyuán. — Dans le quartier, il y a une école, un hôpital, des magasins et un parc.) On utilise la structure d'existence \zh{有}. (3) \zh{居民们都很友好。} (Jūmínmen dōu hěn yǒuhǎo. — Les habitants sont tous très amicaux.) Attention : \zh{都} se place \emph{avant} l'adjectif-verbe \zh{友好}.
\end{exoresolu}

\courssub{Compréhension détaillée : texte d'approfondissement}

Appliquez la même méthode au texte 《我们的社区》.

\begin{exercice}[Questions détaillées]
Répondez en chinois (phrase complète) :
\begin{enumerate}
\item 我们的社区在什么地方？(Où se trouve notre quartier ?)
\item 社区里有哪三个地方？(Quels sont les trois lieux dans le quartier ?)
\item 每天早上居民在公园做什么？(Que font les habitants dans le parc chaque matin ?)
\item 居委会的工作人员怎么样？(Comment sont les employés du comité ?)
\item 上个月社区组织了什么活动？ Combien de personnes ont participé ?
\end{enumerate}
\end{exercice}

\begin{corrige}[Exercice « Questions détaillées »]
(1) \zh{我们的社区在市中心。} (Wǒmen de shèqū zài shì zhōngxīn.)\\
(2) \zh{社区里有一所小学、一家医院和一个大市场。} (Shèqū lǐ yǒu yì suǒ xiǎoxué, yì jiā yīyuàn hé yí ge dà shìchǎng.) \emph{Notez les classificateurs : \zh{所} pour école, \zh{家} pour hôpital/institution, \zh{个} pour marché.}\\
(3) \zh{很多居民在公园里锻炼身体。} (Hěn duō jūmín zài gōngyuán lǐ duànliàn shēntǐ.)\\
(4) \zh{居委会的工作人员很热情。} (Jūwěihuì de gōngzuò rényuán hěn rèqíng.)\\
(5) \zh{上个月，社区组织了一次清洁活动，两百多个居民参加了。} (Shàng ge yuè, shèqū zǔzhī le yí cì qīngjié huódòng, liǎng bǎi duō ge jūmín cānjiā le.)
\end{corrige}

\courssub{Collocations lexicales (associations de mots)}

Mémoriser les mots isolés ne suffit pas ; il faut connaître leurs « amis » habituels (collocations). Complétez le tableau ci-dessous en chinois (caractères + pinyin).

\begin{center}
\begin{tabularx}{\linewidth}{|X|X|X|}
\hline
\rowcolor{popblueL} Structure / Verbe + Objet & Caractères + Pinyin & Traduction \\
\hline
Habiter \emph{dans} un quartier & \zh{住在社区} (zhù zài shèqū) & habiter dans le quartier \\
\hline
Organiser \emph{une activité} & \zh{组织活动} (zǔzhī huódòng) & organiser une activité \\
\hline
Participer \emph{à une activité} & \zh{参加活动} (cānjiā huódòng) & participer à une activité \\
\hline
Faire \emph{de l'exercice} & \zh{锻炼身体} (duànliàn shēntǐ) & faire de l'exercice (lit. exercer le corps) \\
\hline
Aider \emph{les personnes âgées} & \zh{帮助老人} (bāngzhù lǎorén) & aider les personnes âgées \\
\hline
Nettoyer / Faire \emph{le nettoyage} & \zh{搞清洁 / 清洁活动} (gǎo qīngjié / qīngjié huódòng) & faire le nettoyage / activité de nettoyage \\
\hline
\end{tabularx}
\end{center}

\begin{attention}[Classificateurs obligatoires pour dénombrer]
En chinois, on ne dit pas « un quartier » (\zh{一个社区} \textbf{obligatoire}), « une école » (\zh{一所学校}), « un hôpital » (\zh{一家医院}), « un marché » (\zh{一个市场}). Le classificateur \zh{个} est générique, mais \zh{所} (bâtiments publics), \zh{家} (institutions, magasins, restaurants) sont très fréquents au HSK 3-4. Erreur fréquente : \zh{*一学校} (incorrect).
\end{attention}

\coursec{Traduction guidée phrase par phrase}

Passons maintenant à la traduction fine du texte principal. Pour chaque phrase, on observe d'abord la structure, puis on justifie les choix de traduction.

\textbf{Phrase 1.} \zh{我住在一个很大的社区。} (Wǒ zhù zài yí ge hěn dà de shèqū.) — « J'habite dans un très grand quartier. » Structure : Sujet + \zh{住在} + lieu. Le verbe \zh{住} (habiter) exige la préposition \zh{在} devant le lieu, comme « habiter \emph{à/au/dans} » en français. Le classificateur général \zh{个} accompagne \zh{社区}.

\textbf{Phrase 2.} \zh{社区里有学校、医院、商店和公园。} (Shèqū lǐ yǒu xuéxiào, yīyuàn, shāngdiàn hé gōngyuán.) — « Dans le quartier, il y a une école, un hôpital, des magasins et un parc. »

\begin{propriete}[La structure d'existence : Lieu + 里 + 有 + Nom]
Pour dire « il y a … dans un lieu », le chinois utilise : \textbf{Lieu + \zh{里} (lǐ, à l'intérieur) + \zh{有} (yǒu, avoir/il y a) + Nom}. Le lieu vient \emph{en tête} de phrase, contrairement au français « il y a … dans … ».
\begin{itemize}
\item \zh{社区里有学校和医院。} (Shèqū lǐ yǒu xuéxiào hé yīyuàn.) — Dans le quartier, il y a une école et un hôpital.
\item \zh{学校里有很多学生。} (Xuéxiào lǐ yǒu hěn duō xuésheng.) — Dans l'école, il y a beaucoup d'élèves.
\item \zh{雅温得有美丽的公园。} (Yǎwēndé yǒu měilì de gōngyuán.) — Yaoundé a de beaux parcs. (Si le lieu est un nom propre de ville, \zh{里} s'omet souvent).
\end{itemize}
Négation : \zh{没有} (méiyǒu) — \zh{社区里没有医院。} (Shèqū lǐ méiyǒu yīyuàn. — Il n'y a pas d'hôpital dans le quartier.)
\end{propriete}

\textbf{Phrase 3.} \zh{居民们都很友好。} (Jūmínmen dōu hěn yǒuhǎo.) — « Les habitants sont tous très amicaux. » Le suffixe \zh{们} marque le pluriel des personnes (animés) ; \zh{很} (hěn) lie le sujet à l'adjectif, souvent sans valeur intensive réelle (« très » affaibli).

\begin{attention}[La place de 都 (dōu) : avant le verbe, jamais après]
\zh{都} signifie « tous » mais se comporte comme un adverbe : il se place \textbf{avant le verbe ou l'adjectif}, après le sujet. On dit \zh{居民们都很友好} (les habitants sont tous amicaux), jamais \zh{居民们很友好都}. Autres exemples : \zh{我们都参加。} (Wǒmen dōu cānjiā. — Nous participons tous.) ; \zh{他们都是学生。} (Tāmen dōu shì xuésheng. — Ils sont tous élèves.) Piège du français : « tous » se déplace librement en français (« ils viennent tous »), pas en chinois.
\end{attention}

\textbf{Phrase 4.} \zh{居委会经常组织活动，大家都参加。} (Jūwěihuì jīngcháng zǔzhī huódòng, dàjiā dōu cānjiā.) — « Le comité de résidents organise souvent des activités, et tout le monde y participe. » L'adverbe de fréquence \zh{经常} (jīngcháng, souvent) se place avant le verbe, comme \zh{都}. Notez la coordination de deux phrases par la virgule chinoise ，sans mot de liaison : le français ajoute « et ».

\textbf{Phrase 5.} \zh{我喜欢我的社区。} (Wǒ xǐhuan wǒ de shèqū.) — « J'aime mon quartier. » Structure classique : Sujet + verbe + COD. Le possessif \zh{我的} (wǒ de, mon) se construit avec \zh{de}.

\begin{methode}[Traduire 的 dans 我住的社区]
Le titre du texte, \zh{我住的社区}, se traduit « le quartier \emph{où j'habite} ». Méthode :
\begin{enumerate}
\item Identifier le nom final (tête) : \zh{社区} (le quartier).
\item Tout ce qui précède \zh{的} est le qualifiant (antécédent) : \zh{我住} (j'habite).
\item Traduire \zh{的} par une relative française : « que », « où », « qui » selon le sens.
\end{enumerate}
Ainsi : \zh{我住的社区} = le quartier où j'habite ; \zh{我学习的学校} (wǒ xuéxí de xuéxiào) = l'école où j'étudie ; \zh{居民参加的活动} (jūmín cānjiā de huódòng) = l'activité à laquelle les habitants participent. En chinois, le qualifiant vient \emph{toujours avant} le nom, relié par \zh{的} ; en français, la relative vient \emph{après} le nom. C'est l'inverse : pensez-y !
\end{methode}

\begin{center}
\begin{tabularx}{\linewidth}{|X|X|X|}
\hline
\rowcolor{popblueL} Structure & Exemple (caractères + pinyin) & Traduction \\
\hline
Lieu + 里 + 有 + Nom & 社区里有公园。Shèqū lǐ yǒu gōngyuán. & Il y a un parc dans le quartier. \\
\hline
Sujet + 都 + Verbe/Adj & 大家都参加。Dàjiā dōu cānjiā. & Tout le monde participe. \\
\hline
Qualifiant + 的 + Nom & 我住的社区。Wǒ zhù de shèqū. & Le quartier où j'habite. \\
\hline
Sujet + 住在 + Lieu & 我住在杜阿拉。Wǒ zhù zài Dù'ālā. & J'habite à Douala. \\
\hline
Sujet + 经常 + Verbe & 居委会经常组织活动。Jūwěihuì jīngcháng zǔzhī huódòng. & Le comité organise souvent des activités. \\
\hline
\end{tabularx}
\end{center}

\courssub{Schéma de la démarche de lecture}

La figure suivante résume l'organigramme de lecture à suivre pour tout texte chinois : on part du titre, on active les mots connus, on dégage le sens global, et seulement ensuite on passe à la traduction fine.

\begin{popfigure}
\begin{center}
\begin{tikzpicture}[
node distance=0.9cm,
boite/.style={rectangle, rounded corners, draw=popblue, fill=popblueL, thick, align=center, minimum width=5.5cm, minimum height=1.1cm, font=\small},
fleche/.style={-{Stealth[length=3mm]}, thick, draw=popdark}
]
\node[boite, align=center] (titre){\textbf{1. Titre} : \zh{我住的社区}\\ deviner le sujet du texte};
\node[boite, below=of titre, align=center] (mots){\textbf{2. Mots connus} : surligner\\ \zh{社区、居民、医院、活动…}};
\node[boite, below=of mots, align=center] (sens){\textbf{3. Sens global} : qui ? où ?\\ quels services ?};
\node[boite, below=of sens, align=center] (trad){\textbf{4. Traduction fine}\\ phrase par phrase, structures justifiées};
\draw[fleche] (titre) -- (mots);
\draw[fleche] (mots) -- (sens);
\draw[fleche] (sens) -- (trad);
\end{tikzpicture}
\end{center}
\legende{Organigramme de lecture : du titre à la traduction fine.}
\end{popfigure}

\coursec{Production guidée : ma collectivité locale}

Maintenant, à vous de produire un court texte oral ou écrit pour présenter votre propre quartier (au Cameroun ou imaginé en Chine).

\begin{methode}[Rédiger / Présenter : « Mon quartier » en 5 phrases]
Suivez ce canevas en réutilisant les structures de la leçon.
\begin{enumerate}
\item \textbf{Localisation :} \zh{我住在\ldots} (J'habite à \ldots) + \zh{社区 / 街道 / 区} + nom propre (ex. \zh{雅温得}, \zh{杜阿拉}, \zh{布埃亚}).
\item \textbf{Taille / Situation :} \zh{这是一个很大/很小的社区。} / \zh{社区在市中心/郊区。} (C'est un grand/petit quartier. / Le quartier est au centre/banlieue.)
\item \textbf{Services (Structure \zh{有}) :} \zh{社区里有\ldots{}、\ldots{}、\ldots 和\ldots{}。} (Il y a \ldots, \ldots, \ldots et \ldots.) Utilisez au moins 3 noms du vocabulaire + classificateurs si dénombrés.
\item \textbf{Vie sociale (Sujet + 都 + Verbe) :} \zh{居民们都很\ldots{}。居委会经常组织\ldots{}活动，大家都参加。} (Les habitants sont tous \ldots Le comité organise souvent des activités \ldots, tout le monde participe.)
\item \textbf{Avis personnel :} \zh{我\ldots{}我的社区。} (J'aime / Je n'aime pas mon quartier.) \zh{因为\ldots} (Parce que \ldots).
\end{enumerate}
\end{methode}

\begin{exemplebox}[Modèle de production (écrit / oral)]
\zh{我住在雅温得的巴斯社区。这是一个很大的社区。社区里有一所中学、一家医院和一个大市场。居民们都很友好。居委会经常组织足球活动，大家都参加。我喜欢我的社区，因为这里很方便、很热闹。}\\
Wǒ zhù zài Yǎwēndé de Bāsī shèqū. Zhè shì yí ge hěn dà de shèqū. Shèqū lǐ yǒu yì suǒ zhōngxué, yì jiā yīyuàn hé yí ge dà shìchǎng. Jūmínmen dōu hěn yǒuhǎo. Jūwěihuì jīngcháng zǔzhī zúqiú huódòng, dàjiā dōu cānjiā. Wǒ xǐhuan wǒ de shèqū, yīnwèi zhèlǐ hěn fāngbiàn, hěn rènào.
\emph{Traduction :} J'habite au quartier Bastos à Yaoundé. C'est un grand quartier. Il y a un lycée, un hôpital et un grand marché. Les habitants sont tous amicaux. Le comité organise souvent des activités foot, tout le monde participe. J'aime mon quartier car c'est pratique et animé.
\end{exemplebox}

\begin{exercice}[Production personnelle]
Rédigez votre paragraphe « Mon quartier » (5 phrases minimum) en caractères, pinyin et français. Utilisez au moins 8 mots du glossaire. Présentez-le à l'oral devant la classe (1 minute).
\end{exercice}

\begin{corrige}[Exercice « Production personnelle »]
\emph{Attendu :} Structure respectée, classificateurs corrects (\zh{所/家/个}), \zh{都} bien placé avant verbe/adjectif, \zh{的} pour possessif ou relative, ponctuation chinoise. Exemple de correction individuelle par l'enseignant.
\end{corrige}

\begin{savaistu}[Le 居委会 : une institution typiquement chinoise]
Le \zh{居委会} (jūwěihuì), « comité de résidents », est une organisation de quartier très présente dans les villes chinoises : il informe les habitants, organise des fêtes, des activités sportives et des campagnes de propreté. Au Cameroun, on peut le comparer aux comités de quartier et aux réunions de chef de quartier qui règlent la vie quotidienne à Yaoundé ou à Douala. Retenir ce mot, c'est comprendre un peu comment la société chinoise s'organise au niveau le plus proche des habitants.
\end{savaistu}

\begin{aretenir}[L'essentiel de la leçon 17]
\begin{itemize}
\item \textbf{Lecture :} 3 étapes (écoute, chorale, individuelle) ; surligner les mots-thème pour le sens global.
\item \textbf{Vocabulaire clé :} \zh{社区、居民、居委会、学校、医院、商店、公园、市场、组织、活动、参加、友好、热情、锻炼、清洁、幸福}.
\item \textbf{Grammaire :}
      \begin{itemize}
      \item Existence : \textbf{Lieu + 里 + 有 + Nom}.
      \item \zh{都} (tous) : \textbf{Sujet + 都 + Verbe/Adj}.
      \item Relative / Qualifiant : \textbf{Proposition + 的 + Nom}.
      \item Localisation : \textbf{Sujet + 住在 + Lieu}.
      \item Fréquence : \textbf{Sujet + 经常 + Verbe}.
      \end{itemize}
\item \textbf{Écriture :} Radicaux \zh{示、匚、尸、禾、巾、氵、力} ; ordre des traits rigoureux.
\item \textbf{Culture :} \zh{居委会} $\approx$ Comité de quartier (Cameroun) ; \zh{社区} = unité de base urbaine.
\end{itemize}
\end{aretenir}

\coursec{Mise en route et découverte du thème}

\begin{experience}[Brainstorming : ma collectivité, mon environnement]
En classe entière, au tableau, construisons une carte mentale autour du mot \zh{社区} (shèqū, quartier / communauté). 
\begin{enumerate}
\item Le professeur écrit \zh{社区} au centre.
\item Chaque élève propose un mot (lieu, personne, action, service) lié à son quartier (Yaoundé, Douala, Buea, village natal…).
\item Le professeur note en chinois (caractères + pinyin) et en français, en regroupant par catégories : \textbf{Lieux} (\zh{医院}, \zh{商店}, \zh{学校}), \textbf{Acteurs} (\zh{居民}, \zh{市长}, \zh{邻居}), \textbf{Actions} (\zh{住}, \zh{参加}, \zh{组织}), \textbf{Services} (\zh{居委会}, \zh{派出所}).
\end{enumerate}
\textbf{Objectif :} Activer le lexique connu, introduire le nouveau vocabulaire et ancrer le thème dans le vécu camerounais.
\end{experience}

\begin{savaistu}[Collectivités locales : Cameroun vs Chine]
\begin{itemize}
\item \textbf{Cameroun :} Organisation en Régions → Départements → Arrondissements → Communes. Le maire (\zh{市长} / \zh{市长}) est élu au suffrage universel direct. Le quartier (\zh{quartier}) est une division de fait, souvent sans structure administrative formelle élue (chef de quartier nommé).
\item \textbf{Chine :} Organisation en Provinces → Villes-préfectures → Districts/Arrondissements (\zh{区}) → \textbf{Rues/Communautés urbaines} (\zh{街道/社区}) → \textbf{Comités de résidents} (\zh{居委会}). La \zh{社区} (shèqū) est la cellule de base de la gouvernance urbaine, gérée par un \zh{居委会} élu par les résidents. En zone rurale : \zh{村} (cūn, village) géré par un \zh{村委会} (cūnwěihuì).
\end{itemize}
\end{savaistu}

\coursec{Texte support : lecture et compréhension globale}

\begin{methode}[Démarche de lecture en 3 temps]
\begin{enumerate}
\item \textbf{Lecture silencieuse :} Lisez le texte ci-dessous sans dictionnaire. Repérez les mots connus (HSK 1-2) et devinez le sens des nouveaux mots grâce au contexte et aux radicaux.
\item \textbf{Lecture à haute voix :} Le professeur lit, puis les élèves répètent en chœur (tons, liaison, rythme). Attention aux polyphones (\zh{长}, \zh{种}) et aux tons neutres (\zh{们}, \zh{了}).
\item \textbf{Compréhension globale :} Répondez oralement : Qui ? Où ? Quand ? Quoi ? Quel événement principal ?
\end{enumerate}
\end{methode}

\begin{textebox}[Texte original — Niveau HSK 2-3]
\zh{我叫马林，住在喀麦隆的首都雅温得。} \\
\zh{雅温得是一个大城市，人口很多。} \\
\zh{我的社区叫姆沃格-阿达，环境很友好，也很方便。} \\
\zh{社区里有一个居委会，一个小医院和几个商店。} \\
\zh{居民们经常在居委会开会，讨论社区的事情。} \\
\zh{上个周末，社区组织了一次大扫除活动，市长也来参加了。} \\
\zh{大家一起打扫街道，种树，种花。} \\
\zh{活动后，市长说：“雅温得是我们的家，我们要爱护它。”} \\
\zh{我觉得很有意义。我爱我的社区，也爱我的城市。}
\end{textebox}

\begin{definition}[Glossaire du texte (mots hors liste de base)]
\begin{center}
\begin{tabularx}{\linewidth}{|X|X|X|X|}
\hline
\rowcolor{popblueL} \textbf{Caractères} & \textbf{Pinyin} & \textbf{Nature} & \textbf{Traduction} \\
\hline
首都 & shǒudū & n. & capitale \\
人口 & rénkǒu & n. & population \\
经常 & jīngcháng & adv. & souvent, fréquemment \\
开会 & kāihuì & v. & tenir/assister à une réunion \\
讨论 & tǎolùn & v. & discuter, débattre \\
上个周末 & shàngge zhōumò & n. & le week-end dernier \\
大扫除 & dàsǎochú & n./v. & grand nettoyage, opération « coup de balai » \\
打扫 & dǎsǎo & v. & balayer, nettoyer (un lieu) \\
种 & zhòng & v. & planter (arbres, fleurs, riz) \\
爱护 & àihù & v. & protéger, prendre soin de, chérir \\
有意义 & yǒu yìyì & adj. & significatif, plein de sens \\
\hline
\end{tabularx}
\end{center}
\end{definition}

\coursec{Glossaire du thème en tableau}

\courssub{Vocabulaire thématique : noms de lieux, acteurs et institutions}

\begin{definition}[Organisation du glossaire]
Le tableau ci-dessous regroupe le vocabulaire de base de la leçon, classé par catégories sémantiques (lieux, acteurs, institutions/services, verbes/adjectifs). Pour chaque entrée : le caractère simplifié, le pinyin avec tons, la nature grammaticale (n. = nom, v. = verbe, adj. = adjectif, pron. = pronom, adv. = adverbe) et la traduction française. Mémorisez les tons avec rigueur : ils distinguent le sens.
\end{definition}

\begin{center}
\begin{tabularx}{\linewidth}{|X|X|X|X|}
\hline
\rowcolor{popblueL} \textbf{Caractères} & \textbf{Pinyin} & \textbf{Nature} & \textbf{Traduction} \\
\hline
城市 & chéngshì & n. & ville \\
农村 & nóngcūn & n. & campagne, zone rurale \\
社区 & shèqū & n. & communauté, quartier (unité urbaine) \\
村 & cūn & n. & village \\
区 & qū & n. & district, arrondissement \\
\hline
居民 & jūmín & n. & habitant, résident \\
市长 & shìzhǎng & n. & maire \\
大家 & dàjiā & pron. & tout le monde, chacun \\
邻居 & línjū & n. & voisin \\
\hline
政府 & zhèngfǔ & n. & gouvernement, administration (mairie : \zh{市政府}) \\
居委会 & jūwěihuì & n. & comité de résidents (comité de quartier) \\
医院 & yīyuàn & n. & hôpital \\
商店 & shāngdiàn & n. & magasin, boutique \\
学校 & xuéxiào & n. & école \\
派出所 & pàichūsuǒ & n. & commissariat de police (poste de quartier) \\
\hline
住 & zhù & v. & habiter, résider \\
组织 & zǔzhī & v. & organiser \\
参加 & cānjiā & v. & participer (à) \\
讨论 & tǎolùn & v. & discuter, débattre \\
打扫 & dǎsǎo & v. & balayer, nettoyer \\
种 & zhòng & v. & planter \\
爱护 & àihù & v. & protéger, prendre soin de \\
\hline
友好 & yǒuhǎo & adj. & amical, aimable \\
方便 & fāngbiàn & adj. & pratique, commode \\
有意义 & yǒu yìyì & adj. & significatif \\
\hline
\end{tabularx}
\end{center}

\begin{exemplebox}[Exemples contextuels]
\begin{itemize}
\item \zh{我住在喀麦隆的城市。} \emph{Wǒ zhù zài Kāmàilóng de chéngshì.} — J'habite dans une ville du Cameroun.
\item \zh{农村的环境很友好。} \emph{Nóngcūn de huánjìng hěn yǒuhǎo.} — L'environnement de la campagne est agréable.
\item \zh{社区组织活动，大家参加。} \emph{Shèqū zǔzhī huódòng, dàjiā cānjiā.} — Le quartier organise une activité, tout le monde participe.
\item \zh{市长在市政府工作。} \emph{Shìzhǎng zài shìzhèngfǔ gōngzuò.} — Le maire travaille à la mairie (gouvernement municipal).
\item \zh{居委会很方便，居民去办事。} \emph{Jūwěihuì hěn fāngbiàn, jūmín qù bànshì.} — Le comité de quartier est pratique, les habitants y vont pour leurs démarches.
\item \zh{邻居们一起打扫街道。} \emph{Línjūmen yìqǐ dǎsǎo jiēdào.} — Les voisins balayent la rue ensemble.
\item \zh{我们要爱护环境。} \emph{Wǒmen yào àihù huánjìng.} — Nous devons protéger l'environnement.
\end{itemize}
\end{exemplebox}

\begin{definition}[社区 shèqū — communauté de quartier]
En Chine, la \cle{社区} (shèqū) est l'unité de base de l'administration urbaine (niveau sous-district \zh{街道}). Elle regroupe un ensemble d'immeubles ou de rues, gère les affaires courantes (état civil, sécurité, activités culturelles, aide aux personnes âgées, médiation) via le \cle{居委会} (jūwěihuì, comité de résidents élu). C'est l'équivalent fonctionnel de notre « quartier » organisé, avec une structure administrative propre et un budget. Au Cameroun, on parle plutôt de « quartier » (ex. Mvog-Ada, Bonapriso, Nlongkak) sans toujours avoir cette structure formelle de comité de résidents élu ; le chef de quartier est souvent nommé par l'administration.
\end{definition}

\begin{attention}[市 vs 是 — homophones à ne pas confondre]
\begin{itemize}
\item \zh{市} \emph{shì} (4e ton) : ville, municipalité (ex. \zh{城市} chéngshì, \zh{市长} shìzhǎng, \zh{超市} chāoshì supermarché). Radical \zh{巾} (jīn, tissu/étoffe → idée de marché, d'échange sous une tente).
\item \zh{是} \emph{shì} (4e ton) : verbe « être » (identité, classification). Radical \zh{日} (rì, soleil/jour → idée de justesse, de vérité).
\end{itemize}
Même ton, même pinyin, caractères et sens totalement différents. Le contexte disambiguïse toujours. \textbf{Astuce mnémotechnique} : Au \zh{市} (marché), on échange du \zh{巾} (tissu) ; pour affirmer la vérité, on regarde le \zh{日} (soleil).
\end{attention}

\begin{savaistu}[Étymologie de 城 chéng et 村 cūn]
\begin{itemize}
\item \zh{城} (chéng) : radical \zh{土} (terre). À l'origine « muraille en terre battue, rempart ». La Chine antique protégeait ses cités par des murailles ; \zh{长城} (Chángchéng) = « Longue Muraille ». Aujourd'hui \zh{城市} (ville moderne), mais la mémoire des remparts reste.
\item \zh{村} (cūn) : radical \zh{木} (arbre) + phonétique \zh{寸}. « Groupe d'arbres » → hameau, village. Au Cameroun, les chefferies traditionnelles (Bafut, Foumban, Bana) avaient aussi des enceintes de protection et des cases regroupées — analogie culturelle forte.
\end{itemize}
\end{savaistu}

\courssub{Prononciation : tons, homophones et paires minimales}

\begin{propriete}[Tons à surveiller — paires minimales]
\begin{itemize}
\item \zh{市} \emph{shì} (4e ton) vs \zh{是} \emph{shì} (4e ton) — homophones parfaits (voir Attention).
\item \zh{村} \emph{cūn} (1er ton, haut plat) vs \zh{存} \emph{cún} (2e ton, montant) — « village » vs « conserver/déposer ». Radical \zh{木} (arbre) vs \zh{宀}+ \zh{子} (enfant sous un toit).
\item \zh{长} polyphone : \emph{zhǎng} (3e ton, verbe grandir / nom chef : \zh{市长} shìzhǎng, \zh{校长} xiàozhǎng) vs \emph{cháng} (2e ton, adj. long : \zh{长城} chángchéng, \zh{长江} chángjiāng).
\item \zh{区} \emph{qū} (1er ton) vs \zh{取} \emph{qǔ} (3e ton, prendre) — même composant phonétique \zh{区}, tons différents.
\item \zh{种} polyphone : \emph{zhòng} (3e ton, verbe planter : \zh{种树}) vs \emph{zhǒng} (3e ton, nom sorte/type : \zh{种类} zhǒnglèi, \zh{种子} zhǒngzi graine).
\end{itemize}
\end{propriete}

\begin{experience}[Entraînement oral — paires minimales et polyphones]
En binôme. L'élève A lit une phrase, l'élève B identifie le sens/le ton. Changez de rôle.
\begin{enumerate}
\item \zh{这是城市。} (Zhè shì chéngshì. — C'est une ville.) \quad vs \quad \zh{这事成事。} (Zhè shì chéngshì. — Cette affaire réussit / devient une affaire.)
\item \zh{农村很美。} (Nóngcūn hěn měi. — La campagne est belle.) \quad vs \quad \zh{浓存不住。} (Nóng cún bù zhù. — [Phrase inventée] La concentration ne tient pas.)
\item \zh{市长来了。} (Shìzhǎng lái le. — Le maire est arrivé.) \quad vs \quad \zh{是张老三。} (Shì Zhāng Lǎosān. — C'est Zhang Laosan.)
\item \zh{居委会开会。} (Jūwěihuì kāihuì. — Le comité de quartier tient une réunion.) \quad vs \quad \zh{举为伟会。} (Jǔ wéi wěi huì. — [Inventé] Élire comme grande réunion.)
\item \zh{种树 (zhòng shù)} — Planter des arbres. \quad vs \quad \zh{种子 (zhǒngzi)} — Graines / semences.
\end{enumerate}
\end{experience}

\courssub{Clés (radicaux) et ordre des traits : caractères clés du thème}

\begin{propriete}[Radicaux, structure et ordre des traits — aide à la mémorisation et recherche dictionnaire]
\begin{itemize}
\item \zh{城} (chéng) : radical \zh{土} (tǔ, terre) \textbf{à gauche}. Structure : haut (\zh{成}) + bas (\zh{土}). Ordre traits \zh{成} : 一 (heng), 丨 (shù), 一 (heng), 丶 (diǎn), 一 (heng) ; puis \zh{土} : 一, 丨, 一. \textit{Logique : une ville se bâtit sur la terre.}
\item \zh{市} (shì) : radical \zh{巾} (jīn, tissu) \textbf{en bas}. Structure : haut (\zh{亠}) + bas (\zh{巾}). Ordre : \zh{亠} (丶, 一) ; \zh{巾} (丨, (trait brisé), 一, 一). \textit{Logique : marché sous une tente/tissu.}
\item \zh{村} (cūn) : radical \zh{木} (mù, arbre) \textbf{à gauche}. Structure : gauche \zh{木} + droite \zh{寸}. Ordre \zh{木} : 一, 丨, 丿 (piě), 丶 (nà) ; \zh{寸} : 一, 丶, 一, 丶. \textit{Logique : village = groupe d'arbres/maisons.}
\item \zh{区} (qū) : radical \zh{匚} (fāng, boîte ouverte) \textbf{à gauche}. Structure : \zh{匚} + \zh{乂}. Ordre \zh{匚} : 丨, (trait brisé), 一 ; \zh{乂} : 丿, 乀 (ou 丶). \textit{Logique : délimitation administrative (boîte).}
\item \zh{居} (jū) : radical \zh{尸} (shī, corps/maison) \textbf{en haut à gauche}. Structure : \zh{尸} + \zh{古}. Ordre \zh{尸} : 一, (trait brisé), 一 ; \zh{古} : 丨, 一, 一, 丨, 一. \textit{Logique : habiter dans une maison ancienne.}
\item \zh{委} (wěi) : radical \zh{禾} (hé, riz/épi) \textbf{à gauche}. Structure : \zh{禾} + \zh{女}. Ordre \zh{禾} : 一, 丨, 一, 一, 丶, 一 ; \zh{女} : 丿 (piě), 一 (héng), 丨 (shù). \textit{Logique : confier la récolte (禾) à une femme (女) → déléguer.}
\end{itemize}
\end{propriete}

\begin{popfigure}
\begin{tikzpicture}[
  mindmap/.style={concept, align=center, text=white},
  level 1 concept/.append style={font=\small\bfseries, sibling angle=90, level distance=3.8cm},
  level 2 concept/.append style={font=\tiny, sibling angle=45, level distance=3.2cm},
  concept color=popblue!70,
  every node/.style={scale=0.85}
]
\node[mindmap] {Collectivités locales}
  child[concept color=poporange!80] {node[mindmap, align=center]{Lieux \\ Lieux}
    child {node[mindmap] {城市 chéngshì}}
    child {node[mindmap] {农村 nóngcūn}}
    child {node[mindmap] {社区 shèqū}}
    child {node[mindmap] {村 cūn}}
    child {node[mindmap] {区 qū}}
  }
  child[concept color=popgreen!80] {node[mindmap, align=center]{Acteurs \\ Acteurs}
    child {node[mindmap] {居民 jūmín}}
    child {node[mindmap] {市长 shìzhǎng}}
    child {node[mindmap] {大家 dàjiā}}
    child {node[mindmap] {邻居 línjū}}
  }
  child[concept color=poppurple!80] {node[mindmap, align=center]{Institutions \\ Services}
    child {node[mindmap] {政府 zhèngfǔ}}
    child {node[mindmap] {居委会 jūwěihuì}}
    child {node[mindmap] {医院 yīyuàn}}
    child {node[mindmap] {商店 shāngdiàn}}
    child {node[mindmap] {学校 xuéxiào}}
    child {node[mindmap] {派出所 pàichūsuǒ}}
  }
  child[concept color=poppink!80] {node[mindmap, align=center]{Actions / Qualités \\ Actions}
    child {node[mindmap] {住 zhù}}
    child {node[mindmap] {组织 zǔzhī}}
    child {node[mindmap] {参加 cānjiā}}
    child {node[mindmap] {讨论 tǎolùn}}
    child {node[mindmap] {打扫 dǎsǎo}}
    child {node[mindmap] {种 zhòng}}
    child {node[mindmap] {爱护 àihù}}
    child {node[mindmap] {友好 yǒuhǎo}}
    child {node[mindmap] {方便 fāngbiàn}}
  };
\end{tikzpicture}
\legende{Carte mentale lexicale : collectivités locales (couleurs par catégorie sémantique).}
\end{popfigure}

\courssub{Collocations fréquentes et structures modèles}

\begin{exemplebox}[Collocations verbes + compléments / adjectifs + noms]
\begin{center}
\begin{tabularx}{\linewidth}{|X|X|X|}
\hline
\rowcolor{popblueL} \textbf{Structure} & \textbf{Exemple chinois (pinyin)} & \textbf{Traduction} \\
\hline
住 + 在 + Lieu & \zh{我住在社区。} (Wǒ zhù zài shèqū.) & J'habite dans le quartier. \\
组织 + [Activité / Réunion] & \zh{居委会组织活动。} (Jūwěihuì zǔzhī huódòng.) & Le comité organise une activité. \\
参加 + [Réunion / Activité / Concours] & \zh{居民参加会议。} (Jūmín cānjiā huìyì.) & Les habitants participent à la réunion. \\
Sujet + 在 + Lieu + Verbe Action & \zh{市长在政府工作。} (Shìzhǎng zài zhèngfǔ gōngzuò.) & Le maire travaille au gouvernement. \\
Adjectif bisyllabique + 的 + Nom & \zh{友好的环境，方便的商店。} (Yǒuhǎo de huánjìng, fāngbiàn de shāngdiàn.) & Environnement amical, magasins pratiques. \\
Verbe + 了 (aspect accompli) & \zh{社区组织了活动。} (Shèqū zǔzhīle huódòng.) & Le quartier a organisé une activité. \\
\hline
\end{tabularx}
\end{center}
\end{exemplebox}

\begin{attention}[Erreurs fréquentes des francophones — Points de vigilance]
\begin{enumerate}
\item \textbf{Ordre des mots (Complément de lieu) :} En chinois, le complément de lieu \textbf{précède} le verbe d'action/état (Sujet + 在 + Lieu + Verbe). \textit{Français :} « J'habite \textbf{dans} le quartier. » \rightarrow \textit{Chinois :} \zh{我 \textbf{在社区} 住。} (Wǒ \textbf{zài shèqū} zhù.) Ne dites pas *\zh{我住在社区里} (redondant \zh{里} souvent inutile avec \zh{住在}) ni *\zh{我社区住} ni *\zh{我在住社区}.
\item \textbf{的 (de) après adjectif bisyllabique :} Adj. (2 syll.) + \textbf{的} + Nom (\zh{友好的居民}, \zh{方便的商店}). L'omission de \zh{的} n'est tolérée que pour adjectifs monosyllabiques très lexicalisés (\zh{好人}, \zh{大山}) ou noms de parenté (\zh{我妈}). \textbf{Règle HSK 2-3 :} Bisyllabique = \zh{的} obligatoire.
\item \textbf{Mots composés indivisibles :} \zh{市长} (maire), \zh{居委会} (comité de résidents), \zh{派出所} (commissariat). Ne séparez pas : *\zh{市很长} (la ville est longue), *\zh{居委会开会} (correct) mais ne traduisez pas « le comité de la réunion des résidents ».
\item \textbf{大家 (dàjiā) = « tout le monde » (sujet collectif), pas « grande famille ».} \zh{大家庭} (dàjiātíng) = grande famille. \zh{大家都去} = Tout le monde y va.
\item \textbf{种 (zhòng / zhǒng) :} Verbe « planter » = 3e ton \emph{zhòng} (\zh{种树}). Nom « sorte, espèce, graine » = 3e ton \emph{zhǒng} (\zh{种类}, \zh{种子}). Même ton, sens radicalement différents selon la place.
\end{enumerate}
\end{attention}

\courssub{Exercice résolu : emploi du vocabulaire en contexte}

\begin{exoresolu}[Compléter, donner le pinyin et traduire]
Complétez les phrases avec le mot approprié du glossaire.
\begin{enumerate}
\item \zh{我们\_\_\_在喀麦隆的\_\_\_。} (ville / habiter)
\item \zh{这个\_\_\_很\_\_\_，大家都喜欢。} (quartier / amical)
\item \zh{\_\_\_每天在\_\_\_工作。} (maire / gouvernement / mairie)
\item \zh{农村的\_\_\_去\_\_\_看病。} (habitants / hôpital)
\item \zh{上个周末，社区\_\_\_了一次大扫除\_\_\_。} (organiser / activité)
\item \zh{市长\_\_\_了这次活动。} (participer)
\end{enumerate}

\tcblower
\textbf{Solution détaillée :}
\begin{enumerate}
\item \zh{我们\textbf{住}在喀麦隆的\textbf{城市}。} \emph{Wǒmen \textbf{zhù} zài Kāmàilóng de \textbf{chéngshì}.} — Nous habitons dans une ville du Cameroun. (Structure : Sujet + 住 + 在 + Lieu)
\item \zh{这个\textbf{社区}很\textbf{友好}，大家都喜欢。} \emph{Zhège \textbf{shèqū} hěn \textbf{yǒuhǎo}, dàjiā dōu xǐhuan.} — Ce quartier est très amical, tout le monde l'aime. (Adj. bisyllabique + 很 ; \zh{大家} = tout le monde sujet)
\item \zh{\textbf{市长}每天在\textbf{市政府}工作。} \emph{\textbf{Shìzhǎng} měitiān zài \textbf{shìzhèngfǔ} gōngzuò.} — Le maire travaille à la mairie tous les jours. (Précision : \zh{市政府} pour mairie ; temps \zh{每天} avant lieu \zh{在...})
\item \zh{农村的\textbf{居民}去\textbf{医院}看病。} \emph{Nóngcūn de \textbf{jūmín} qù \textbf{yīyuàn} kànbìng.} — Les habitants de la campagne vont à l'hôpital se soigner. (\zh{看病} kànbìng = consulter un médecin, locution verbale courante)
\item \zh{上个周末，社区\textbf{组织}了一次大扫除\textbf{活动}。} \emph{Shàngge zhōumò, shèqū \textbf{zǔzhī}le yícì dàsǎochú \textbf{huódòng}.} — Le week-end dernier, le quartier a organisé une activité de grand nettoyage. (\zh{了} après verbe pour accompli ; \zh{一次} classificateur d'événement)
\item \zh{市长\textbf{参加}了这次活动。} \emph{Shìzhǎng \textbf{cānjiā}le zhècì huódòng.} — Le maire a participé à cette activité. (\zh{参加} + \zh{了} + Objet)
\end{enumerate}
\end{exoresolu}

\coursec{Questions de compréhension et collocations}

\begin{exercice}[Compréhension du texte « Ma collectivité locale »]
\begin{enumerate}
\item Répondez en chinois (caractères + pinyin) d'après le texte :
  \begin{enumerate}
  \item 马林住在哪里？
  \item 他的社区叫什么名字？
  \item 社区里有哪些设施 (shèshī, équipements) ?
  \item 居民们经常在居委会做什么？
  \item 上个周末社区组织了什么活动？谁参加了？
  \item 市长说了什么？马林怎么觉得 (juéde) ?
  \end{enumerate}
\item Reliez le verbe à son complément d'objet / lieu habituel (collocation) :
  \begin{center}
  \begin{tabular}{lll}
  \zh{住} & a. \zh{活动} & 1. \_\_\_ \\
  \zh{组织} & b. \zh{在社区} & 2. \_\_\_ \\
  \zh{参加} & c. \zh{会议} & 3. \_\_\_ \\
  \zh{打扫} & d. \zh{街道} & 4. \_\_\_ \\
  \zh{种} & e. \zh{树 / 花} & 5. \_\_\_ \\
  \zh{爱护} & f. \zh{环境 / 家} & 6. \_\_\_ \\
  \end{tabular}
  \end{center}
\item Traduisez en chinois (caractères + pinyin) :
  \begin{enumerate}
  \item Les habitants du village participent à la réunion.
  \item Le maire travaille à la mairie (gouvernement municipal).
  \item Ce magasin est très pratique.
  \item Nous aimons notre quartier, nous devons le protéger.
  \item Le week-end dernier, j'ai planté des arbres avec mes voisins.
  \end{enumerate}
\end{enumerate}
\end{exercice}

\begin{corrige}[Exercice Compréhension et Collocations]
\begin{enumerate}
\item \begin{enumerate}
  \item \zh{马林住在喀麦隆的首都雅温得。} (Mǎlín zhù zài Kāmàilóng de shǒudū Yǎwēndé.)
  \item \zh{他的社区叫姆沃格-阿达。} (Tā de shèqū jiào Mǔwògé-Ādá.)
  \item \zh{有一个居委会，一个小医院和几个商店。} (Yǒu yíge jūwěihuì, yíge xiǎo yīyuàn hé jǐge shāngdiàn.)
  \item \zh{经常开会，讨论社区的事情。} (Jīngcháng kāihuì, tǎolùn shèqū de shìqing.)
  \item \zh{社区组织了一次大扫除活动，市长也来参加了。} (Shèqū zǔzhīle yícì dàsǎochú huódòng, shìzhǎng yě lái cānjiāle.)
  \item \zh{市长说：“雅温得是我们的家，我们要爱护它。” 马林觉得很有意义。} (Shìzhǎng shuō: “Yǎwēndé shì wǒmen de jiā, wǒmen yào àihù tā.” Mǎlín juéde hěn yǒu yìyì.)
  \end{enumerate}
\item \zh{住} — b (\zh{住在社区}) ; \zh{组织} — a (\zh{组织活动}) ; \zh{参加} — c (\zh{参加会议}) ; \zh{打扫} — d (\zh{打扫街道}) ; \zh{种} — e (\zh{种树/花}) ; \zh{爱护} — f (\zh{爱护环境/家}).
\item \begin{enumerate}
  \item \zh{村里的居民参加会议。} (Cūnli de jūmín cānjiā huìyì.)
  \item \zh{市长在市政府工作。} (Shìzhǎng zài shìzhèngfǔ gōngzuò.)
  \item \zh{这个商店很方便。} (Zhège shāngdiàn hěn fāngbiàn.)
  \item \zh{我们爱我们的社区，要爱护它。} (Wǒmen ài wǒmen de shèqū, yào àihù tā.)
  \item \zh{上个周末，我和邻居一起种树。} (Shàngge zhōumò, wǒ hé línjū yìqǐ zhòng shù.)
  \end{enumerate}
\end{enumerate}
\end{corrige}

\coursec{Production guidée : ma collectivité locale}

\begin{methode}[Rédiger une courte présentation de sa collectivité (80-100 caractères)]
\begin{enumerate}
\item \textbf{Identité :} Nom, lieu de résidence (\zh{我叫...，住在...}).
\item \textbf{Type :} Ville / Campagne / Quartier (\zh{这是...城市/农村/社区}).
\item \textbf{Environnement :} Adjectifs + \zh{的} (\zh{环境很友好，很方便}).
\item \textbf{Équipements :} \zh{有...，有...，还有...} (liste : \zh{居委会}, \zh{医院}, \zh{商店}, \zh{学校}, \zh{派出所}).
\item \textbf{Vie collective :} Activité récente (\zh{上个周末，社区组织了...活动，大家参加了}).
\item \textbf{Sentiment / Engagement :} \zh{我觉得...，我爱..., 要爱护...}.
\end{enumerate}
\end{methode}

\begin{experience}[Exposé oral : « Mon quartier / Mon village »]
En binôme. Chaque élève prépare 1 minute de présentation orale (sans lire) en suivant la méthode ci-dessus. L'écouteur note 3 informations clés en français. Échange de rôles. 3 binômes présentent devant la classe. Évaluation : prononciation (tons), fluidité, richesse lexicale (min. 8 mots du glossaire), structure correcte (lieu avant verbe, \zh{的}).
\end{experience}

\begin{exercice}[Expression écrite : Ma collectivité locale]
Rédigez un paragraphe de 80 à 100 caractères chinois (simplifiés) pour présenter votre quartier (Yaoundé, Douala, Buea, village natal...) ou une collectivité locale idéale. Utilisez au moins 10 mots du glossaire de la leçon. Respectez l'ordre des mots (Sujet + Temps + Lieu + Verbe) et la particule \zh{的} après adjectifs bisyllabiques.
\begin{itemize}
\item \textbf{Consigne :} Écrivez les caractères, le pinyin (avec tons) et la traduction française.
\item \textbf{Mots-outils suggérés :} \zh{也} (aussi), \zh{还} (en plus), \zh{因为...所以...} (parce que... donc), \zh{虽然...但是...} (bien que... mais).
\end{itemize}
\end{exercice}

\begin{aretenir}[Bilan de la leçon 17 — Collectivités locales]
\begin{itemize}
\item \textbf{Lexique clé (25+ mots) :} Lieux (\zh{城市, 农村, 社区, 村, 区}), Acteurs (\zh{居民, 市长, 大家, 邻居}), Institutions (\zh{政府, 居委会, 医院, 商店, 学校, 派出所}), Actions (\zh{住, 组织, 参加, 讨论, 打扫, 种, 爱护}), Qualités (\zh{友好, 方便, 有意义}).
\item \textbf{Grammaire :} Structure \zh{住在 + Lieu} ; Complément de lieu \textbf{avant} verbe (\zh{在...工作}) ; \zh{的} obligatoire après adj. bisyllabiques ; Aspect accompli \zh{了} ; Polyphones \zh{长} (zhǎng/cháng), \zh{种} (zhòng/zhǒng).
\item \textbf{Culture :} \zh{社区} + \zh{居委会} (Chine, base gouvernance urbaine) vs Quartier + Chef de quartier (Cameroun, administration décentralisée). Analogie historique : \zh{城} (remparts) et chefferies traditionnelles camerounaises.
\item \textbf{Écriture :} Radicaux \zh{土, 巾, 木, 匚, 尸, 禾} ; Ordre des traits standard (haut→bas, gauche→droite, horizontal→vertical).
\item \textbf{Compétences :} Lire/Comprendre un texte HSK 2-3 ; Répondre à des questions ; Produire un court exposé/écrit structuré.
\end{itemize}
\end{aretenir}

\coursec{Clés (radicaux) et ordre des traits}

Dans la section précédente, nous avons appris le vocabulaire des collectivités locales : \zh{城市} (chéngshì, ville), \zh{区} (qū, quartier, district), \zh{村子} (cūnzi, village), \zh{居民} (jūmín, habitant), \zh{住} (zhù, habiter). Mais un bon sinisant ne se contente pas de reconnaître les caractères : il sait les \cle{écrire} correctement. Or, écrire le chinois n'est pas un dessin libre : c'est un geste codifié, régi par des règles précises. Cette section vous donne les deux clés de voûte de l'écriture : les \cle{clés} (radicaux, \zh{部首} bùshǒu), qui portent le sens, et l'\cle{ordre des traits} (\zh{笔顺} bǐshùn), qui garantit une écriture lisible, rapide et belle.

\courssub{Pourquoi les clés et l'ordre des traits sont-ils indispensables ?}

\begin{definition}[Clé ou radical (\zh{部首} bùshǒu)]
La \cle{clé} (ou radical) est la partie d'un caractère qui indique sa \textbf{famille de sens}. C'est aussi l'élément sous lequel le caractère est classé dans les dictionnaires. Exemple : dans \zh{城} (chéng, ville), la clé est \zh{土} (tǔ, terre) : la ville est bâtie sur la terre.
\end{definition}

\begin{definition}[Trait (\zh{笔画} bǐhuà)]
Le \cle{trait} est le plus petit geste d'écriture : chaque fois que l'on pose le stylo, qu'on trace et qu'on relève, cela compte pour \textbf{un} trait. Les traits fondamentaux sont : le trait horizontal \zh{一} (héng), le vertical \zh{丨} (shù), le tombant à gauche \zh{丿} (piě), le tombant à droite \zh{乀} (nà), le point \zh{丶} (diǎn) et le crochet (gōu).
\end{definition}

\begin{propriete}[Règles générales de l'ordre des traits]
L'écriture chinoise obéit à cinq lois fondamentales :
\begin{enumerate}
\item \textbf{De haut en bas} : on trace d'abord le haut du caractère, puis le bas. Exemple : \zh{三} (sān, trois) se trace du trait supérieur vers le trait inférieur.
\item \textbf{De gauche à droite} : on trace d'abord la partie gauche, puis la partie droite. Exemple : \zh{村} : d'abord \zh{木} à gauche, puis \zh{寸} à droite.
\item \textbf{L'horizontal avant le vertical} quand ils se croisent : dans \zh{十} (shí, dix), on trace d'abord \zh{一} puis \zh{丨}.
\item \textbf{L'extérieur avant l'intérieur} : pour les caractères à enclos, on trace d'abord le contour, puis le contenu.
\item \textbf{On ferme l'enclos en dernier} : le trait qui « ferme la boîte » se trace \textbf{après} le contenu. Exemple : dans \zh{区}, le trait final en « L couché » ferme la boîte \zh{匚} une fois \zh{乂} écrit à l'intérieur.
\end{enumerate}
\end{propriete}

\begin{attention}[Erreur fréquente des francophones]
En français, nous écrivons les lettres de gauche à droite sans règle de geste stricte pour chaque trait. Beaucoup d'élèves « dessinent » les caractères chinois en commençant par où bon leur semble. Résultat : caractères déformés, illisibles, impossibles à mémoriser. \textbf{Respecter l'ordre des traits, c'est mémoriser deux fois plus vite}, car la mémoire du geste aide la mémoire visuelle. Ne « dessinez » jamais un caractère : \textbf{écrivez-le}, trait après trait, en comptant à voix haute.
\end{attention}

\courssub{Étude caractère par caractère}

\textbf{1. \zh{城} chéng — la ville}\par
Clé : \zh{土} (tǔ, la terre), placée à \textbf{gauche}. Partie droite : \zh{成} (chéng, accomplir), qui donne le son. Total : \textbf{9 traits}.
Ordre : on trace d'abord la clé \zh{土} en 3 traits (horizontal court, vertical, horizontal long qui remonte légèrement en \emph{tí}), puis \zh{成} en 6 traits : 1) horizontal haut, 2) tombant gauche (piě), 3) horizontal milieu, 4) vertical, 5) crochet horizontal-vertical (héng zhé gōu), 6) point (diǎn) en haut à droite.
Sens de la clé : les anciennes villes chinoises étaient entourées de murs de \textbf{terre} battue — d'où la clé 土.

\begin{popfigure}
\begin{center}
\begin{tikzpicture}[scale=0.9, font=\small]
% Grille du caractère
\draw[popline, dashed] (0,0) rectangle (4,4);
\draw[popline, dashed] (2,0) -- (2,4);
\draw[popline, dashed] (0,2) -- (4,2);
% Partie gauche : 土
\node[popblue, align=center] at (1,3.6) {\zh{土}};
\draw[-{Stealth}, popblue, thick] (0.5,3.1) -- (1.5,3.1) node[midway, above, popdark]{\tiny 1};
\draw[-{Stealth}, popblue, thick] (1.0,3.35) -- (1.0,2.35) node[midway, right, popdark]{\tiny 2};
\draw[-{Stealth}, popblue, thick] (0.35,2.3) -- (1.65,2.45) node[midway, below, popdark]{\tiny 3};
% Partie droite : 成
\node[poporange, align=center] at (3,3.6) {\zh{成}};
\draw[-{Stealth}, poporange, thick] (2.35,3.15) -- (3.55,3.15) node[midway, above, popdark]{\tiny 4};
\draw[-{Stealth}, poporange, thick] (2.5,3.05) -- (2.35,1.9) node[midway, left, popdark]{\tiny 5};
\draw[-{Stealth}, poporange, thick] (2.55,2.75) -- (3.15,2.75) node[midway, above, popdark]{\tiny 6};
\draw[-{Stealth}, poporange, thick] (3.15,2.75) -- (3.15,2.05) node[midway, right, popdark]{\tiny 7};
\draw[-{Stealth}, poporange, thick] (3.15,2.05) -- (3.6,2.05) -- (3.6,0.9) node[midway, right, popdark]{\tiny 8};
\draw[-{Stealth}, poporange, thick] (3.3,3.5) -- (3.6,3.25) node[midway, right, popdark]{\tiny 9};
\node[align=center, popdark] at (2,-0.5) {D'abord \zh{土} (traits 1--3), puis \zh{成} (traits 4--9)};
\end{tikzpicture}
\end{center}
\legende{Décomposition du caractère \zh{城} (chéng, ville) en 9 traits : la clé \zh{土} à gauche, puis \zh{成} à droite.}
\end{popfigure}

\textbf{2. \zh{区} qū — le quartier, le district}\par
Clé : \zh{匚} (fāng, la boîte, le récipient). Total : \textbf{4 traits}.
Ordre : c'est l'exemple parfait de la règle « on ferme la boîte à la fin ». Trait 1 : l'horizontal supérieur de la boîte. Traits 2 et 3 : le contenu \zh{乂} (le tombant à gauche, puis le tombant à droite). Trait 4 : le grand trait en « L couché » (vertical descendant puis horizontal) qui \textbf{ferme} la boîte par la gauche et le bas.
Piège : ne tracez jamais tout le contour d'abord ! Le côté gauche et le bas se ferment \textbf{après} le contenu.

\textbf{3. \zh{村} cūn — le village}\par
Clé : \zh{木} (mù, le bois, l'arbre), à gauche. Total : \textbf{7 traits}.
Ordre : d'abord \zh{木} en 4 traits (horizontal, vertical, tombant à gauche, point à droite — attention : à gauche d'un caractère, le dernier trait de \zh{木} devient un \textbf{point} (diǎn), pas un grand tombant), puis \zh{寸} en 3 traits (horizontal, vertical avec crochet, point).
Sens de la clé : les villages traditionnels étaient construits en \textbf{bois}, au milieu des arbres.

\textbf{4. \zh{民} mín — le peuple}\par
Clé : \zh{氏} (shì, clan, famille). Total : \textbf{5 traits}.
Ordre : du haut vers le bas. Trait 1 : le petit trait horizontal-crochet (héng zhé) en haut. Trait 2 : le vertical court. Trait 3 : l'horizontal. Trait 4 : l'horizontal relevé (tí). Trait 5 : le grand tombant oblique (nà) qui descend vers la droite.
On retrouve \zh{民} dans \zh{居民} (jūmín, habitant), \zh{人民} (rénmín, le peuple), \zh{农民} (nóngmín, paysan).

\textbf{5. \zh{住} zhù — habiter}\par
Clé : \zh{亻} (rén, l'homme debout), variante gauche de \zh{人}. Total : \textbf{7 traits}.
Ordre : d'abord la clé \zh{亻} en 2 traits (tombant gauche, puis vertical), puis \zh{主} (zhǔ, maître) en 5 traits : 1) point, 2) horizontal, 3) horizontal, 4) vertical traversant, 5) horizontal final (bas).
Sens : un \textbf{homme} (\zh{亻}) chez le \textbf{maître} de maison (\zh{主}) = habiter, résider.

\courssub{Tableau récapitulatif}

\begin{center}
\begin{tabularx}{\linewidth}{|X|X|X|X|}
\hline
\rowcolor{popblueL} Caractère (pinyin, sens) & Clé (radical) & Nombre de traits & Ordre essentiel \\
\hline
\zh{城} chéng, ville & \zh{土} tǔ, terre & 9 & \zh{土} à gauche, puis \zh{成} \\
\hline
\zh{区} qū, quartier & \zh{匚} fāng, boîte & 4 & contenu \zh{乂}, boîte fermée en dernier \\
\hline
\zh{村} cūn, village & \zh{木} mù, bois & 7 & \zh{木} à gauche (dernier trait = point), puis \zh{寸} \\
\hline
\zh{民} mín, peuple & \zh{氏} shì, clan & 5 & du haut vers le bas \\
\hline
\zh{住} zhù, habiter & \zh{亻} rén, homme & 7 & \zh{亻} d'abord, puis \zh{主} \\
\hline
\end{tabularx}
\end{center}

\begin{methode}[Mémoriser un caractère par sa clé : la méthode du sens]
Pour retenir durablement un caractère, ne l'apprenez jamais comme un dessin : reliez la \textbf{clé au sens}.
\begin{enumerate}
\item \textbf{Identifiez la clé} et demandez-vous : quel rapport avec le sens ? \zh{城} a la clé \zh{土} (terre) : la ville est bâtie sur la terre, entourée de murs de terre. \zh{村} a la clé \zh{木} (bois) : le village est fait de maisons en bois, au milieu des arbres.
\item \textbf{Identifiez la partie phonétique} : dans \zh{城}, \zh{成} (chéng) donne le son ; dans \zh{住}, \zh{主} (zhǔ) donne presque le son (zhù).
\item \textbf{Fabriquez une petite histoire} : « Un homme (\zh{亻}) chez le maître (\zh{主}) : il \textbf{habite} là » pour \zh{住}.
\item \textbf{Écrivez 5 fois} en comptant les traits à voix haute : « un, deux, trois… » (\zh{一、二、三}…).
\end{enumerate}
\end{methode}

\begin{attention}[Ne pas confondre \zh{住} et \zh{注} : une clé différente change tout le sens]
\zh{住} (zhù, habiter) a la clé de l'\textbf{homme} \zh{亻}. \zh{注} (zhù, noter, verser, concentrer) a la clé de l'\textbf{eau} \zh{氵}. Même prononciation, même partie droite \zh{主}, mais sens totalement différents ! Comparez : \zh{我住在雅温得。} (Wǒ zhù zài Yǎwēndé. « J'habite à Yaoundé. ») et \zh{请注意！} (Qǐng zhùyì ! « Faites attention ! », littéralement « versez/concentrez votre attention »). En chinois, \textbf{la clé est le gardien du sens} : un seul trait d'écart, et le mot change de famille.
\end{attention}

\courssub{Application : lire les clés dans un texte}

\begin{textebox}[Mon quartier]
我住在杜阿拉的一个区。这个区不大，但是居民很多。\\
Wǒ zhù zài Dù'ālā de yí ge qū. Zhège qū bù dà, dànshì jūmín hěn duō.\\
« J'habite dans un quartier de Douala. Ce quartier n'est pas grand, mais il y a beaucoup d'habitants. »\\
区里有一个小村子，村子旁边有很多树。\\
Qū lǐ yǒu yí ge xiǎo cūnzi, cūnzi pángbiān yǒu hěn duō shù.\\
« Dans le district, il y a un petit village ; à côté du village, il y a beaucoup d'arbres. »\\
我们的城市很美，人民很友好。\\
Wǒmen de chéngshì hěn měi, rénmín hěn yǒuhǎo.\\
« Notre ville est belle et les gens sont très amicaux. »
\end{textebox}

Observez : dans ce texte, retrouvez \zh{住} (clé homme), \zh{区} (clé boîte), \zh{村} (clé bois), \zh{城} (clé terre), \zh{民} dans \zh{居民} et \zh{人民}. Chaque clé vous donne un indice de sens avant même de connaître le mot.

\begin{exoresolu}[Compter les traits et identifier les clés]
Énoncé : pour chaque caractère, donne la clé et le nombre de traits : \zh{城}, \zh{区}, \zh{村}, \zh{民}, \zh{住}. Puis recopie \zh{城} en numérotant les trois premiers traits.
\tcblower
Solution détaillée :
\begin{itemize}
\item \zh{城} : clé \zh{土} (terre), 9 traits. Les trois premiers traits sont ceux de \zh{土} : 1) horizontal court, 2) vertical, 3) horizontal long du bas (relevé).
\item \zh{区} : clé \zh{匚} (boîte), 4 traits. Attention : le contenu \zh{乂} se trace avant la fermeture de la boîte.
\item \zh{村} : clé \zh{木} (bois), 7 traits (4 pour \zh{木}, 3 pour \zh{寸}).
\item \zh{民} : clé \zh{氏}, 5 traits, tracés du haut vers le bas.
\item \zh{住} : clé \zh{亻} (homme), 7 traits (2 pour \zh{亻}, 5 pour \zh{主}).
\end{itemize}
Vérification : si tu as trouvé 8 traits pour \zh{城}, tu as probablement fusionné deux traits de \zh{成} ; si tu as fermé la boîte de \zh{区} avant d'écrire \zh{乂}, relis la règle « on ferme l'enclos en dernier ».
\end{exoresolu}

\begin{exercice}[Tracé sur ardoise ou cahier]
Sur ton ardoise ou ton cahier, écris chaque caractère \textbf{5 fois} en comptant les traits à voix haute en chinois (\zh{一、二、三、四}…) : \zh{城} (9 traits), \zh{区} (4 traits), \zh{村} (7 traits), \zh{民} (5 traits), \zh{住} (7 traits). Puis écris la phrase : \zh{我住在一个村子里。} (Wǒ zhù zài yí ge cūnzi lǐ. « J'habite dans un village. ») en respectant l'ordre des traits de chaque caractère.
\end{exercice}

\begin{aretenir}[L'essentiel de la section]
\begin{itemize}
\item La \cle{clé} (radical) porte le sens : \zh{土} terre dans \zh{城}, \zh{木} bois dans \zh{村}, \zh{亻} homme dans \zh{住}, \zh{匚} boîte dans \zh{区}, \zh{氏} clan dans \zh{民}.
\item Ordre des traits : de haut en bas, de gauche à droite, l'extérieur avant l'intérieur, et \textbf{on ferme les enclos en dernier} (\zh{区}).
\item Une clé différente = un sens différent : \zh{住} (habiter) contre \zh{注} (noter/verser).
\item Écrire en comptant les traits est la meilleure façon de mémoriser.
\end{itemize}
\end{aretenir}

\coursec{Questions de compréhension et collocations}

Dans la section précédente, nous avons appris à écrire les caractères clés du thème des collectivités locales en respectant leurs clés (radicaux) et l'ordre des traits. Savoir écrire, c'est bien ; savoir \emph{employer} ces mots dans des phrases complètes pour répondre à des questions, c'est encore mieux. C'est exactement ce que l'on vous demandera au baccalauréat : lire un texte, puis répondre par des phrases entières, en chinois, en reprenant la formulation de la question. Cette section vous donne la méthode, les réponses modèles et les collocations (associations de mots) indispensables pour parler de votre quartier, de votre village ou de votre ville.

\courssub{Observer avant de répondre : la méthode des questions de compréhension}

\begin{methode}[Répondre à une question de compréhension au Bac]
Pour obtenir tous les points en compréhension de texte, suivez toujours ces quatre étapes :
\begin{enumerate}
\item \textbf{Repérer le mot interrogatif} : 哪里 (où), 什么 (quoi), 怎么样 (comment), 谁 (qui), 为什么 (pourquoi), 什么时候 (quand).
\item \textbf{Retrouver dans le texte la phrase qui contient la réponse} : soulignez-la mentalement.
\item \textbf{Reprendre les mots de la question} dans votre réponse : si la question est \zh{他住在哪里？}, votre réponse doit commencer par \zh{他住在...}.
\item \textbf{Rédiger une phrase complète} (sujet + verbe + complément + point 。), jamais un mot isolé.
\end{enumerate}
\end{methode}

\begin{attention}[L'erreur qui coûte des points]
Ne répondez jamais par un seul mot. À la question \zh{他住在哪里？}, répondre seulement \zh{雅温得} est incomplet. La réponse attendue est une phrase entière : \zh{他住在雅温得。} (\emph{Tā zhù zài Yǎwēndé.} — Il habite à Yaoundé). Au baccalauréat, une réponse sans verbe est sanctionnée, même si l'information est juste.
\end{attention}

\courssub{Questions de compréhension sur le texte support}

Reprenons le texte de la section 2 (le jeune Camerounais qui présente son quartier et son comité de résidents) et répondons aux quatre questions officielles.

\textbf{Question 1 :} \zh{他住在哪里？} (\emph{Tā zhù zài nǎli?} — Où habite-t-il ?)

Réponse complète : \zh{他住在雅温得的一个社区里。} (\emph{Tā zhù zài Yǎwēndé de yí ge shèqū li.} — Il habite dans un quartier de Yaoundé.)

\textbf{Question 2 :} \zh{社区里有什么？} (\emph{Shèqū li yǒu shénme?} — Qu'y a-t-il dans le quartier ?)

Réponse complète : \zh{社区里有学校、医院、市场和小公园。} (\emph{Shèqū li yǒu xuéxiào, yīyuàn, shìchǎng hé xiǎo gōngyuán.} — Dans le quartier, il y a une école, un hôpital, un marché et un petit parc.)

\textbf{Question 3 :} \zh{居民们怎么样？} (\emph{Jūmínmen zěnmeyàng?} — Comment sont les habitants ?)

Réponse complète : \zh{居民们很友好，也很团结。} (\emph{Jūmínmen hěn yǒuhǎo, yě hěn tuánjié.} — Les habitants sont très amicaux et aussi très solidaires.)

\textbf{Question 4 :} \zh{居委会经常做什么？} (\emph{Jūwěihuì jīngcháng zuò shénme?} — Que fait souvent le comité de résidents ?)

Réponse complète : \zh{居委会经常组织活动，比如打扫卫生和庆祝节日。} (\emph{Jūwěihuì jīngcháng zǔzhī huódòng, bǐrú dǎsǎo wèishēng hé qìngzhù jiérì.} — Le comité de résidents organise souvent des activités, comme le nettoyage et la fête des festivals.)

\begin{exemplebox}[Observer la mécanique de la réponse]
Comparez question et réponse : la réponse \emph{recopie} le début de la question et remplace seulement le mot interrogatif par l'information du texte.\\
\zh{他住在\cle{哪里}？} → \zh{他住在\cle{雅温得的一个社区里}。}\\
\zh{社区里有\cle{什么}？} → \zh{社区里有\cle{学校、医院、市场和小公园}。}\\
C'est la même logique qu'en français : « Où habite-t-il ? » → « Il habite \emph{à Yaoundé} ».
\end{exemplebox}

\courssub{Les cinq collocations clés du thème}

Une \cle{collocation} est une association habituelle de mots : on ne dit pas n'importe quel verbe avec n'importe quel nom. En chinois comme en français, certaines combinaisons sont fixes. Voici les cinq collocations à connaître par cœur pour parler des collectivités locales.

\begin{propriete}[Collocation 1 : 住在 + lieu = habiter à / dans]
Structure : \textbf{Sujet + 住在 + lieu (+ 里/这儿/那儿)}. La préposition \zh{在} est \textbf{obligatoire} entre le verbe \zh{住} (habiter) et le lieu.\\
\zh{我住在雅温得。} (\emph{Wǒ zhù zài Yǎwēndé.} — J'habite à Yaoundé.)\\
\zh{他们住在一个小村子里。} (\emph{Tāmen zhù zài yí ge xiǎo cūnzi li.} — Ils habitent dans un petit village.)
\end{propriete}

\begin{attention}[Le piège n° 1 des francophones : oublier 在]
En français, « habiter » est transitif : « j'habite \emph{Yaoundé} », sans préposition. Le francophone traduit donc mot à mot et dit *\zh{我住雅温得} : c'est \textbf{faux}. En chinois, \zh{住} ne peut pas être suivi directement du lieu ; il faut \zh{在}. On dit \zh{住在社区} (habiter dans le quartier), mais jamais *\zh{住社区里} sans \zh{在}. Retenez : \textbf{住 + 在 + lieu}, toujours.
\end{attention}

\begin{propriete}[Collocation 2 : lieu + 里有 + choses = il y a... dans...]
Structure : \textbf{Lieu + 里有 + nom(s)}. C'est la tournure d'existence : \zh{有} exprime « il y a ».\\
\zh{社区里有一个市场。} (\emph{Shèqū li yǒu yí ge shìchǎng.} — Il y a un marché dans le quartier.)\\
\zh{我们村有很多树。} (\emph{Wǒmen cūn yǒu hěn duō shù.} — Il y a beaucoup d'arbres dans notre village.)\\
Négation : \zh{村里没有医院。} (\emph{Cūn li méiyǒu yīyuàn.} — Il n'y a pas d'hôpital dans le village.) Attention : la négation de \zh{有} est toujours \zh{没有}, jamais *\zh{不有}.
\end{propriete}

\begin{propriete}[Collocation 3 et 4 : 组织活动 / 参加活动]
\zh{组织活动} (\emph{zǔzhī huódòng}) = organiser des activités ; \zh{参加活动} (\emph{cānjiā huódòng}) = participer à des activités. Le nom \zh{活动} (activité) se combine avec ces deux verbes, comme en français « organiser / participer à une activité ».\\
\zh{居委会经常组织活动。} (\emph{Jūwěihuì jīngcháng zǔzhī huódòng.} — Le comité organise souvent des activités.)\\
\zh{居民们都参加活动。} (\emph{Jūmínmen dōu cānjiā huódòng.} — Les habitants participent tous aux activités.)
\end{propriete}

\begin{propriete}[Collocation 5 : 对 + personne + 友好 = être amical envers qqn]
Structure : \textbf{Sujet + 对 + personne + (很) + 友好}. La préposition \zh{对} (envers) introduit la personne visée ; elle se place \emph{avant} l'adjectif, contrairement au français où « envers » vient après l'adjectif.\\
\zh{他们对我们很友好。} (\emph{Tāmen duì wǒmen hěn yǒuhǎo.} — Ils sont très amicaux envers nous.)\\
\zh{邻居对新来的家庭很友好。} (\emph{Línjū duì xīn lái de jiātíng hěn yǒuhǎo.} — Les voisins sont très accueillants envers la famille nouvellement arrivée.)
\end{propriete}

\begin{attention}[Ordre des mots : 对 avant, pas après]
Le français dit « amical \emph{envers nous} » (adjectif puis préposition). Le chinois fait l'inverse : \zh{对我们很友好} (envers-nous très amical). Ne dites jamais *\zh{他们很友好对我们} : le complément introduit par \zh{对} se place obligatoirement \textbf{avant} l'adjectif.
\end{attention}

\begin{center}
\begin{tabularx}{\linewidth}{|X|X|X|X|}
\hline
\rowcolor{popblueL} Collocation & Pinyin & Structure & Traduction \\
\hline
住在 + lieu & zhù zài & Sujet + 住在 + lieu & habiter à / dans \\
\hline
lieu + 里有 & li yǒu & Lieu + 里有 + nom & il y a... dans... \\
\hline
组织活动 & zǔzhī huódòng & Verbe + nom & organiser des activités \\
\hline
参加活动 & cānjiā huódòng & Verbe + nom & participer aux activités \\
\hline
对...友好 & duì... yǒuhǎo & Sujet + 对 + qqn + 友好 & être amical envers qqn \\
\hline
\end{tabularx}
\end{center}

\begin{popfigure}
\begin{tikzpicture}[node distance=0.9cm,
  struct/.style={rectangle, rounded corners, draw=popblue, fill=popblueL, align=center, font=\small, inner sep=5pt},
  var/.style={rectangle, rounded corners, draw=poporange, fill=poporangeL, align=center, font=\small, inner sep=5pt},
  lbl/.style={font=\small\itshape, text=popdark}]
  \node[struct] (a1) {住在};
  \node[var, right=of a1, align=center] (a2){[lieu]\\雅温得 / 社区 / 村};
  \draw[-{Stealth}, thick, popblue] (a1) -- (a2);
  \node[lbl, left=0.15cm of a1] {habiter};

  \node[struct, below=1.1cm of a1] (b1) {[lieu] 里有};
  \node[var, right=of b1, align=center] (b2){[choses]\\学校、医院、市场};
  \draw[-{Stealth}, thick, popblue] (b1) -- (b2);
  \node[lbl, left=0.15cm of b1] {il y a};

  \node[struct, below=1.1cm of b1] (c1) {对};
  \node[var, right=of c1, align=center] (c2){[personne]\\我们 / 邻居};
  \node[struct, right=of c2] (c3) {很友好};
  \draw[-{Stealth}, thick, popblue] (c1) -- (c2);
  \draw[-{Stealth}, thick, popblue] (c2) -- (c3);
  \node[lbl, left=0.15cm of c1] {être amical};
\end{tikzpicture}
\legende{Les trois structures à trous : la partie fixe (bleu) ne change jamais, la partie variable (orange) s'adapte au lieu, aux choses et aux personnes.}
\end{popfigure}

\courssub{Exercice de substitution : parler d'autres localités}

Le principe de la \cle{substitution} est simple : on garde la structure fixe et on remplace l'élément variable. C'est l'exercice le plus efficace pour automatiser les collocations.

\begin{exoresolu}[Substitution guidée]
Énoncé : en partant du modèle \zh{我住在雅温得。社区里有学校和市场。}, remplacez les éléments soulignés pour parler de (a) Douala avec un hôpital et un parc ; (b) votre village avec beaucoup d'arbres et une école.
\tcblower
Solution détaillée :\\
(a) On remplace le lieu \zh{雅温得} par \zh{杜阿拉} (Douala) et les choses par \zh{医院和公园} :\\
\zh{我住在杜阿拉。社区里有医院和公园。} (\emph{Wǒ zhù zài Dù'ālā. Shèqū li yǒu yīyuàn hé gōngyuán.} — J'habite à Douala. Dans le quartier, il y a un hôpital et un parc.)\\
(b) Pour un village, on remplace \zh{社区} (quartier) par \zh{村} (village) et on ajoute le classificateur pour les arbres :\\
\zh{我住在一个村子里。村里有很多棵树和一所学校。} (\emph{Wǒ zhù zài yí ge cūnzi li. Cūn li yǒu hěn duō kē shù hé yì suǒ xuéxiào.} — J'habite dans un village. Dans le village, il y a beaucoup d'arbres et une école.)\\
Vérifiez toujours : \zh{在} après \zh{住}, et \zh{有} (et non \zh{是}) pour « il y a ».
\end{exoresolu}

\begin{exercice}[À vous de substituer]
Produisez trois phrases complètes en remplaçant les éléments variables : 1) dire que vous habitez à Buea et qu'il y a une université dans la ville ; 2) dire que le comité de votre quartier organise des matchs de football et que les jeunes y participent ; 3) dire que les voisins sont amicaux envers les étudiants étrangers.
\end{exercice}

\begin{corrige}[Exercice de substitution]
1) \zh{我住在布埃亚。城里有一所大学。} (\emph{Wǒ zhù zài Bù'āiyà. Chéng li yǒu yì suǒ dàxué.} — J'habite à Buea. Il y a une université dans la ville.)\\
2) \zh{我们社区的居委会经常组织足球比赛，年轻人都参加。} (\emph{Wǒmen shèqū de jūwěihuì jīngcháng zǔzhī zúqiú bǐsài, niánqīngrén dōu cānjiā.} — Le comité de résidents de notre quartier organise souvent des matchs de football, et les jeunes y participent tous.)\\
3) \zh{邻居们对外国学生很友好。} (\emph{Línjūmen duì wàiguó xuéshēng hěn yǒuhǎo.} — Les voisins sont très amicaux envers les étudiants étrangers.)
\end{corrige}

\courssub{Mini-comparaison : 城市 et 农村 (ville et campagne)}

Pour comparer la ville et la campagne, deux paires d'adjectifs suffisent à ce niveau : \zh{大} (grand) / \zh{小} (petit) et \zh{多} (nombreux) / \zh{少} (peu nombreux). On les combine avec les structures déjà vues.

\begin{exemplebox}[Ville et campagne]
\zh{城市很大，人很多。} (\emph{Chéngshì hěn dà, rén hěn duō.} — La ville est grande, il y a beaucoup de gens.)\\
\zh{农村很小，人很少。} (\emph{Nóngcūn hěn xiǎo, rén hěn shǎo.} — La campagne est petite, il y a peu de gens.)\\
\zh{城市里有很多汽车，农村里有很多树。} (\emph{Chéngshì li yǒu hěn duō qìchē, nóngcūn li yǒu hěn duō shù.} — En ville il y a beaucoup de voitures, à la campagne il y a beaucoup d'arbres.)
\end{exemplebox}

\begin{center}
\begin{tabularx}{\linewidth}{|X|X|X|}
\hline
\rowcolor{popblueL} Structure & Exemple (caractères + pinyin) & Traduction \\
\hline
城市 + 大 / 人多 & 城市很大，人很多。Chéngshì hěn dà, rén hěn duō. & La ville est grande et peuplée. \\
\hline
农村 + 小 / 人少 & 农村很小，人很少。Nóngcūn hěn xiǎo, rén hěn shǎo. & La campagne est petite et peu peuplée. \\
\hline
Lieu + 里有 + 多/少 + nom & 村里车很少。Cūn li chē hěn shǎo. & Il y a peu de voitures au village. \\
\hline
\end{tabularx}
\end{center}

\begin{attention}[很 ne veut pas toujours dire « très »]
Devant un adjectif monosyllabique, \zh{很} sert surtout de liaison : \zh{城市很大} signifie simplement « la ville est grande », pas forcément « très grande ». Sans \zh{很}, la phrase \zh{城市大} sonne comme une comparaison (« la ville, elle, est grande » — par opposition à une autre). Pour un francophone, c'est contre-intuitif : retenez que \zh{很} est le lien neutre entre le sujet et l'adjectif.
\end{attention}

\begin{textebox}[Texte d'entraînement : 我的家乡]
\zh{我的家乡在喀麦隆的西部。它不大，是一座小城市。}\\
\emph{Wǒ de jiāxiāng zài Kāmàilóng de xībù. Tā bù dà, shì yí zuò xiǎo chéngshì.}\\
\zh{城里有一个大市场、两所学校和一家医院。}\\
\emph{Chéng li yǒu yí ge dà shìchǎng, liǎng suǒ xuéxiào hé yì jiā yīyuàn.}\\
\zh{这里的人很友好，对外国人也很热情。}\\
\emph{Zhèli de rén hěn yǒuhǎo, duì wàiguórén yě hěn rèqíng.}\\
\zh{每个星期六，社区都组织活动，年轻人和老人都参加。}\\
\emph{Měi ge xīngqīliù, shèqū dōu zǔzhī huódòng, niánqīngrén hé lǎorén dōu cānjiā.}\\
\zh{我很喜欢我的家乡。}\\
\emph{Wǒ hěn xǐhuan wǒ de jiāxiāng.}\\
« Ma ville natale se trouve dans l'ouest du Cameroun. Elle n'est pas grande, c'est une petite ville. Il y a un grand marché, deux écoles et un hôpital. Les gens d'ici sont très amicaux et chaleureux envers les étrangers. Chaque samedi, le quartier organise des activités, et les jeunes comme les personnes âgées y participent. J'aime beaucoup ma ville natale. »
\end{textebox}

Questions rapides sur ce texte d'entraînement : 1) \zh{他的家乡在哪里？} 2) \zh{城里有什么？} 3) \zh{社区什么时候组织活动？} Appliquez la méthode en quatre étapes : repérez le mot interrogatif (哪里, 什么, 什么时候), retrouvez la phrase du texte, reprenez la formulation, rédigez une phrase complète. Réponses attendues : \zh{他的家乡在喀麦隆的西部。} / \zh{城里有一个大市场、两所学校和一家医院。} / \zh{社区每个星期六都组织活动。}

\begin{aretenir}[L'essentiel de la section]
\begin{itemize}
\item Répondre à une question = reprendre la formulation de la question + remplacer le mot interrogatif par l'information du texte + phrase complète.
\item \zh{住在 + lieu} : la préposition \zh{在} est obligatoire après \zh{住}.
\item \zh{Lieu + 里有 + nom} : « il y a... dans... » ; négation avec \zh{没有}.
\item \zh{组织活动} / \zh{参加活动} : organiser / participer à des activités.
\item \zh{对 + personne + 友好} : le complément en \zh{对} se place avant l'adjectif, contrairement au français.
\item Ville vs campagne : \zh{大/小} et \zh{多/少} ; \zh{很} est le lien neutre devant l'adjectif.
\item Classificateurs utiles : \zh{所} (écoles, universités), \zh{家} (hôpitaux, entreprises), \zh{座} (villes, montagnes), \zh{棵} (arbres).
\end{itemize}
\end{aretenir}

\coursec{Production guidée : ma collectivité locale}

Après avoir lu le texte support, vérifié votre compréhension et mémorisé les collocations essentielles (\zh{住在}, \zh{有}, \zh{居民}, \zh{方便}…), il est temps de passer à la production : vous allez à votre tour décrire votre commune, votre quartier ou votre village en chinois. C'est l'objectif final de la leçon : réutiliser le vocabulaire et les structures dans un texte personnel, court mais correct.

\courssub{Observer d'abord : un modèle de description}

Avant d'écrire, lisons attentivement un modèle rédigé par un élève de Terminale. Observez comment les quatre phrases s'enchaînent logiquement : le lieu, ce qu'on y trouve, les habitants, puis l'opinion personnelle.

\begin{textebox}[Modèle : mon quartier à Douala]
我住在杜阿拉的一个社区。\\
Wǒ zhù zài Dù'ālā de yí ge shèqū.\\
那里有学校、医院和市场。\\
Nàli yǒu xuéxiào, yīyuàn hé shìchǎng.\\
居民们很友好，生活很方便。\\
Jūmínmen hěn yǒuhǎo, shēnghuó hěn fāngbiàn.\\
我喜欢我的社区，因为它很热闹。\\
Wǒ xǐhuan wǒ de shèqū, yīnwèi tā hěn rènao.\\[0.5em]
\emph{J'habite dans un quartier de Douala. Il y a une école, un hôpital et un marché. Les habitants sont amicaux, la vie y est pratique. J'aime mon quartier parce qu'il est animé.}
\end{textebox}

\textbf{Observations guidées.} Répondez mentalement avant de lire la suite :

\begin{enumerate}
\item Quel verbe exprime « habiter » ? Où se place le complément de lieu ?
\item Quel mot introduit l'énumération des services (école, hôpital, marché) ?
\item Que remarquez-vous devant les adjectifs \zh{友好} et \zh{方便} ?
\item Quel mot relie l'opinion (\zh{我喜欢我的社区}) à sa justification ?
\end{enumerate}

Vous avez trouvé : \zh{住在} + lieu (le lieu suit le verbe, contrairement au français où la préposition vient avant le nom), \zh{有} pour « il y a », \zh{很} devant chaque adjectif prédicat, et \zh{因为} pour justifier. Ce sont exactement les quatre piliers de votre production.

\courssub{La trame de production}

Votre tâche consiste à rédiger \cle{4 à 6 phrases} en suivant cette trame obligatoire :

\begin{center}
\begin{tabularx}{\linewidth}{|X|X|X|}
\hline
\rowcolor{popblueL} Étape & Trame chinoise + pinyin & Fonction \\
\hline
1. Lieu & \zh{我住在…} (Wǒ zhù zài…) & Situer : ville, commune, quartier, village \\
\hline
2. Services & \zh{那里有…} (Nàli yǒu…) & Énumérer : école, hôpital, marché, mairie \\
\hline
3. Habitants & \zh{居民们很…} (Jūmínmen hěn…) & Décrire : amicaux, travailleurs, nombreux \\
\hline
4. Opinion & \zh{我喜欢…因为…} (Wǒ xǐhuan… yīnwèi…) & Jugement personnel justifié \\
\hline
\end{tabularx}
\end{center}

\begin{propriete}[Rappel des deux structures-clés : \zh{在} et \zh{有}]
\begin{center}
\begin{tabularx}{\linewidth}{|X|X|X|}
\hline
\rowcolor{popblueL} Structure & Exemple (caractères + pinyin) & Traduction \\
\hline
Sujet + Verbe + \zh{在} + Lieu & \zh{我住在雅温得。} (Wǒ zhù zài Yǎwēndé.) & J'habite à Yaoundé. \\
\hline
Lieu + \zh{有} + Chose/Personne & \zh{那里有一个大市场。} (Nàli yǒu yí ge dà shìchǎng.) & Il y a un grand marché là-bas. \\
\hline
\end{tabularx}
\end{center}
\zh{在} (zài) fonctionne comme \textbf{verbe de localisation} ou \textbf{préposition} : le lieu le suit. \zh{有} (yǒu) exprime l'\textbf{existence} : le lieu est le sujet (ou le thème) et précède \zh{有}, contrairement au français « il y a » où le verbe précède.
\end{propriete}

\begin{attention}[很 devant l'adjectif prédicat : quasi obligatoire]
En français, « les habitants sont amicaux » se traduit par \zh{居民们很友好}. Le mot \zh{很} (hěn), littéralement « très », est ici presque \textbf{obligatoire} et ne signifie pas vraiment « très » : il sert de lien (copule) entre le sujet et l'adjectif. Dire \zh{*居民们友好} sans \zh{很} sonne faux ou crée un sens comparatif (« les habitants sont \emph{plus} amicaux »). Retenez : \textbf{Sujet + \zh{很} + Adjectif}. Même chose pour \zh{生活很方便}, \zh{社区很热闹}, \zh{村很干净}.
\end{attention}

\courssub{Banque de mots autorisée}

Vous ne pouvez utiliser que le glossaire de la leçon et le vocabulaire acquis en Première. Voici la banque complète à mobiliser :

\begin{center}
\begin{tabularx}{\linewidth}{|X|X|X|X|}
\hline
\rowcolor{popblueL} Caractères & Pinyin & Nature & Traduction \\
\hline
住 & zhù & verbe & habiter, résider \\
\hline
在 & zài & prép./verbe & être à, dans \\
\hline
有 & yǒu & verbe & avoir, il y a \\
\hline
社区 & shèqū & nom & communauté, quartier résidentiel \\
\hline
居民 & jūmín & nom & habitant, résident \\
\hline
城市 & chéngshì & nom & ville \\
\hline
村 & cūn & nom & village \\
\hline
区 & qū & nom & district, arrondissement \\
\hline
学校 & xuéxiào & nom & école \\
\hline
医院 & yīyuàn & nom & hôpital \\
\hline
市场 & shìchǎng & nom & marché \\
\hline
方便 & fāngbiàn & adj. & pratique, commode \\
\hline
友好 & yǒuhǎo & adj. & amical \\
\hline
热闹 & rènao & adj. & animé, vivant \\
\hline
干净 & gānjìng & adj. & propre \\
\hline
因为 & yīnwèi & conj. & parce que \\
\hline
\end{tabularx}
\end{center}

\courssub{Écriture des caractères : radicaux et ordre des traits (sélection)}

\begin{definition}[Clés et ordre des traits — personnages du thème]
Maîtriser l'écriture passe par la connaissance du radical (clé) et de l'ordre des traits (haut → bas, gauche → droite, horizontal avant vertical, trait central avant les ailes, cadre avant contenu, fermer le cadre en dernier).
\begin{center}
\begin{tabularx}{\linewidth}{|X|X|X|X|}
\hline
\rowcolor{popblueL} Caractère & Radical (Clé) & Nombre de traits & Ordre des traits (résumé) \\
\hline
\zh{住} & \zh{亻} (rén, homme) & 7 & 亻 → 主 (point, horizontale, verticale, horizontale, verticale, horizontale, horizontale) \\
\hline
\zh{在} & \zh{土} (tǔ, terre) & 6 & 土 → 才 (horizontale, verticale, horizontale ; puis 才 : verticale, horizontale, horizontale) \\
\hline
\zh{有} & \zh{月} (yuè, viande/lune) & 6 & 月 → 丶 (cadre 月 : verticale gauche, horizontales pliées, horizontale ; point final à droite) \\
\hline
\zh{区} & \zh{匚} (fāng, boîte) & 4 & 匚 → 乂 (verticale gauche, horizontale pliée, horizontale ; puis 乂 à l'intérieur) \\
\hline
\zh{村} & \zh{木} (mù, bois) & 7 & 木 → 寸 (arbre complet ; puis 寸 : horizontale, verticale, point, horizontale, point, horizontale) \\
\hline
\zh{民} & \zh{氏} (shì, clan) & 5 & 氏 (horizontale, verticale, horizontale, horizontale, horizontale) \\
\hline
\zh{医} & \zh{匚} (fāng, boîte) & 7 & 匚 → 矢 (cadre ; puis flèche 矢 : horizontale, verticale, horizontale, horizontale, verticale, horizontale) \\
\hline
\zh{市} & \zh{巾} (jīn, étoffe) & 5 & 巾 (verticale, horizontale pliée, verticale, horizontale, horizontale) \\
\hline
\zh{热} & \zh{灬} (huǒ, feu) & 10 & 执 (tenue) → 灬 (quatre points feu en bas) \\
\hline
\zh{因} & \zh{囗} (wéi, enceinte) & 6 & 囗 → 大 (cadre fermé en dernier ; grand 大 à l'intérieur) \\
\hline
\end{tabularx}
\end{center}
\emph{Attention :} \zh{市} (simplifié, 5 traits, clé 巾) ne s'écrit pas \zh{巿} (traditionnel, clé 巿). \zh{住} (clé 亻) ne se confond pas avec \zh{往} (clé 彳).
\end{definition}

\courssub{Organigramme de la production}

La figure suivante résume le cheminement de votre texte : chaque phrase correspond à une étape, et chaque étape appelle une structure précise.

\begin{popfigure}
\begin{tikzpicture}[
node distance=0.55cm and 1.6cm,
etape/.style={rectangle, rounded corners=4pt, draw=popblue, fill=popblueL, thick, align=center, minimum width=2.6cm, minimum height=1.1cm, font=\small},
struct/.style={rectangle, rounded corners=4pt, draw=poporange, fill=poporangeL, thick, align=center, minimum width=3.6cm, minimum height=1.1cm, font=\small},
fleche/.style={-{Stealth[length=3mm]}, thick, popdark}
]
\node[etape, align=center] (lieu){1. Lieu\\我住在…};
\node[etape, below=of lieu, align=center] (serv){2. Services\\那里有…};
\node[etape, below=of serv, align=center] (hab){3. Habitants\\居民们很…};
\node[etape, below=of hab, align=center] (op){4. Opinion\\我喜欢…因为…};
\node[struct, right=of lieu] (s1) {Sujet + 住在 + Lieu};
\node[struct, right=of serv] (s2) {Lieu + 有 + Nom};
\node[struct, right=of hab] (s3) {Sujet + 很 + Adjectif};
\node[struct, right=of op] (s4) {喜欢 + 因为 + Cause};
\draw[fleche] (lieu) -- (serv);
\draw[fleche] (serv) -- (hab);
\draw[fleche] (hab) -- (op);
\draw[fleche, popmuted, dashed] (lieu) -- (s1);
\draw[fleche, popmuted, dashed] (serv) -- (s2);
\draw[fleche, popmuted, dashed] (hab) -- (s3);
\draw[fleche, popmuted, dashed] (op) -- (s4);
\end{tikzpicture}
\legende{Organigramme de la production guidée : du lieu à l'opinion, chaque étape a sa structure.}
\end{popfigure}

\courssub{Enrichir son texte : la justification avec 因为}

\begin{methode}[Ajouter 因为 pour valoriser sa copie]
Une opinion sans justification reste pauvre. À l'examen, la structure \zh{因为} est valorisée car elle montre que vous savez construire une phrase complexe. Procédez ainsi :
\begin{enumerate}
\item Formulez d'abord votre opinion simple : \zh{我喜欢我的社区。} (Wǒ xǐhuan wǒ de shèqū.) « J'aime mon quartier. »
\item Demandez-vous « pourquoi ? » et choisissez un adjectif de la banque : \zh{热闹}, \zh{方便}, \zh{干净}, \zh{友好}.
\item Reliez les deux avec \zh{因为} : \zh{我喜欢我的社区，因为它很热闹。} « J'aime mon quartier parce qu'il est animé. »
\item Vérifiez : le \zh{很} est-il présent devant l'adjectif ? La virgule chinoise \zh{，} sépare-t-elle bien les deux propositions ?
\end{enumerate}
Variante possible (structure corrélative) : \zh{因为那里很方便，所以我喜欢我的社区。} (Yīnwèi nàli hěn fāngbiàn, suǒyǐ wǒ xǐhuan wǒ de shèqū.) « Comme c'est pratique là-bas, j'aime mon quartier. » La paire \zh{因为…所以…} (« parce que… donc… ») est correcte en chinois, alors que le français n'emploie qu'un seul des deux mots.
\end{methode}

\courssub{Exercice résolu : compléter puis rédiger}

\begin{exoresolu}[Compléter la trame]
Énoncé. Complétez le texte suivant avec les mots de la banque (\zh{在}, \zh{有}, \zh{很}, \zh{因为}), puis traduisez-le :\\
我住＿＿布埃亚的一个村。那里＿＿学校和市场。居民们＿＿友好。我喜欢我的村，＿＿它很干净。
\tcblower
Solution détaillée.\\
1. \zh{我住在布埃亚的一个村。} (Wǒ zhù zài Bù'āiyà de yí ge cūn.) — Après le verbe \zh{住}, il faut la préposition de lieu \zh{在}. « J'habite dans un village de Buea. »\\
2. \zh{那里有学校和市场。} (Nàli yǒu xuéxiào hé shìchǎng.) — Pour exprimer « il y a », on utilise \zh{有} après le lieu \zh{那里}. « Il y a une école et un marché. »\\
3. \zh{居民们很友好。} (Jūmínmen hěn yǒuhǎo.) — L'adjectif prédicat exige \zh{很}. « Les habitants sont amicaux. »\\
4. \zh{我喜欢我的村，因为它很干净。} (Wǒ xǐhuan wǒ de cūn, yīnwèi tā hěn gānjìng.) — La justification s'introduit par \zh{因为}. « J'aime mon village parce qu'il est propre. »\\
Vérification finale : quatre phrases, quatre structures différentes, aucun adjectif sans \zh{很}, ponctuation chinoise respectée.
\end{exoresolu}

\courssub{Déroulement de l'activité en classe}

\textbf{Phase 1 : rédaction individuelle (10 minutes).} Chaque élève rédige seul, au brouillon, ses 4 à 6 phrases en suivant la trame. Choisissez votre vrai lieu de vie : votre quartier de Yaoundé ou de Douala, votre village d'origine, votre commune. Interdiction de copier le modèle mot pour mot : changez au moins le lieu, deux services et l'adjectif final.

\textbf{Phase 2 : relecture en binôme (5 minutes).} Échangez votre brouillon avec votre voisin et contrôlez trois points précis, dans cet ordre :

\begin{enumerate}
\item \textbf{Les tons} : lisez le texte à voix haute ; votre binôme vérifie le pinyin (par exemple \zh{zhù} 4\textsuperscript{e} ton, \zh{yǒu} 3\textsuperscript{e} ton, \zh{hěn} 3\textsuperscript{e} ton, \zh{yīnwèi} 4-4).
\item \textbf{Les caractères} : chaque caractère est-il complet et lisible ? Attention aux confusions : \zh{住} (clé « homme » 亻) ne s'écrit pas comme \zh{往} ; \zh{市} (marché, clé 巾) ne prend pas le trait supplémentaire de \zh{巿} (traditionnel).
\item \textbf{在 et 有} : \zh{在} sert-il bien à localiser (\zh{住在} + lieu) et \zh{有} à exprimer l'existence (\zh{那里有…}) ? C'est l'erreur la plus fréquente des francophones, qui traduisent « il y a » par \zh{*在} : on ne dit pas \zh{*那里有在…}.
\end{enumerate}

\textbf{Phase 3 : présentation orale et rétroaction collective (5 minutes).} Deux ou trois volontaires lisent leur texte devant la classe. La classe écoute avec une grille simple : a-t-on entendu les quatre étapes (lieu, services, habitants, opinion) ? Le \zh{很} est-il prononcé ? La rétroaction commence toujours par un point positif avant la correction.

\begin{experience}[Présentation orale : mon quartier en une minute]
Présentez votre texte sans le lire mot pour mot : regardez la classe, parlez lentement, marquez les tons. Vos camarades lèvent la main pour citer une phrase qu'ils ont comprise et la traduisent en français. L'enseignant note au tableau les bonnes structures entendues et corrige collectivement une seule erreur récurrente (souvent l'oubli de \zh{很} ou l'inversion \zh{在}/\zh{有}).
\end{experience}

\courssub{Complément examen (hors-programme utile)}

\begin{aretenir}[Phrase d'ouverture type pour un paragraphe descriptif au Bac]
Au baccalauréat, une production écrite gagne des points dès sa première phrase si elle annonce clairement le sujet. Mémorisez ce modèle d'ouverture :\\
\zh{我想介绍一下我的家乡。} (Wǒ xiǎng jièshào yíxià wǒ de jiāxiāng.) « Je voudrais présenter brièvement ma région natale. »\\
Variante pour un quartier : \zh{我给大家介绍一下我住的社区。} (Wǒ gěi dàjiā jièshào yíxià wǒ zhù de shèqū.) « Je vous présente le quartier où j'habite. »\\
La structure \zh{给 + quelqu'un + 介绍} (« présenter à quelqu'un ») et la nuance de brièveté \zh{一下} sont des marqueurs de niveau appréciés du correcteur. Enchaînez ensuite avec la trame de la leçon : \zh{那里有…}, \zh{居民们很…}, \zh{我喜欢…因为…}.
\end{aretenir}

\begin{exercice}[À vous d'écrire]
Rédigez votre propre description en 4 à 6 phrases : votre commune, votre quartier ou votre village. Contraintes obligatoires : utiliser au moins une fois \zh{住在}, \zh{有}, \zh{很} + adjectif, et une justification avec \zh{因为}. Ajoutez si possible un deuxième service (\zh{医院}, \zh{市场}, \zh{学校}) relié par \zh{和}. Relisez-vous avec la grille de la phase 2 avant de rendre votre copie.
\end{exercice}

\begin{corrige}[Exercice : éléments de correction]
Il n'existe pas une seule bonne réponse, mais toute copie correcte contient : (1) une phrase de localisation du type \zh{我住在…的一个社区。} ; (2) une phrase d'existence \zh{那里有…和…。} ; (3) une description des habitants avec \zh{很} : \zh{居民们很友好。} ; (4) une opinion justifiée : \zh{我喜欢…，因为…很…。}. Exemple de copie attendue : \zh{我住在雅温得的一个区。那里有学校和医院。居民们很友好，生活很方便。我喜欢我的区，因为它很热闹。} (Wǒ zhù zài Yǎwēndé de yí ge qū. Nàli yǒu xuéxiào hé yīyuàn. Jūmínmen hěn yǒuhǎo, shēnghuó hěn fāngbiàn. Wǒ xǐhuan wǒ de qū, yīnwèi tā hěn rènao.) « J'habite dans un quartier de Yaoundé. Il y a une école et un hôpital. Les habitants sont amicaux, la vie y est pratique. J'aime mon quartier parce qu'il est animé. » Retirez un point pour chaque adjectif sans \zh{很} et pour chaque confusion \zh{在}/\zh{有}.
\end{corrige}

\newpage

\coursec{Activité d'intégration}

\courssub{Objectifs de l'activité}

À l'issue de cette activité d'intégration, tu seras capable de :
\begin{itemize}
\item réutiliser le vocabulaire des \cle{collectivités locales} (commune, quartier, mairie, services, habitants) dans une tâche de communication réelle ;
\item lire et exploiter un document authentique simple (un message d'invitation chinois) ;
\item rédiger 5 phrases descriptives correctes en chinois (ordre des mots, tons, ponctuation pleine chasse) ;
\item tracer 3 caractères du thème en respectant les clés et l'ordre des traits ;
\item présenter oralement une affiche en 1 minute, avec une prononciation claire et des tons maîtrisés.
\end{itemize}

\courssub{Situation}

Une délégation de maires et d'étudiants de la ville chinoise de \cle{Hangzhou} (杭州), ville partenaire de ta région, visitera ton lycée le mois prochain. Le proviseur souhaite décorer le hall d'accueil avec des affiches réalisées par les élèves de Terminale : chaque affiche présente une commune ou un quartier du Cameroun. Le chef de la délégation a envoyé ce message au lycée :

\begin{textebox}[Message du chef de la délégation chinoise]
亲爱的同学们，你们好！\\
Qīn'ài de tóngxuémen, nǐmen hǎo !\\
Chers élèves, bonjour !\\[4pt]
下个月我们要访问你们的城市。\\
Xià ge yuè wǒmen yào fǎngwèn nǐmen de chéngshì.\\
Le mois prochain, nous visiterons votre ville.\\[4pt]
我们很想了解你们的社区：那里有多少人？有什么服务？居民喜欢做什么？\\
Wǒmen hěn xiǎng liǎojiě nǐmen de shèqū : nàli yǒu duōshao rén ? Yǒu shénme fúwù ? Jūmín xǐhuan zuò shénme ?\\
Nous aimerions beaucoup découvrir votre communauté : combien y a-t-il d'habitants ? Quels services y a-t-il ? Qu'aiment faire les habitants ?\\[4pt]
请给我们介绍你们的街区吧！\\
Qǐng gěi wǒmen jièshào nǐmen de jiēqū ba !\\
Présentez-nous votre quartier, s'il vous plaît !
\end{textebox}

Ta classe répond à cette invitation : par groupes de 3, vous préparez une affiche « \zh{欢迎来到我的社区！} » (Huānyíng lái dào wǒ de shèqū ! — « Bienvenue dans ma communauté ! ») et vous la présentez oralement en 1 minute devant la classe, qui joue le rôle de la délégation.

\courssub{Consignes}

\begin{enumerate}
\item \textbf{Comprendre le message.} Lis le message du chef de délégation. Relève dans le texte : (a) deux mots du vocabulaire des collectivités locales ; (b) les trois questions posées par les visiteurs. Ton affiche devra répondre à ces trois questions.
\item \textbf{Choisir le lieu.} En groupe, choisissez votre commune ou votre quartier (exemples : un quartier de Yaoundé ou de Douala, la commune de Buea, un arrondissement de ta ville). Notez en français 5 informations : le lieu, un service public, les habitants, une activité, votre opinion.
\item \textbf{Rédiger l'affiche.} Écris en chinois : un \cle{titre} (ex. \zh{欢迎来到……！}) et \cle{5 phrases} : une phrase sur le lieu, une sur les services, une sur les habitants, une sur les activités, une phrase d'opinion avec \zh{我觉得}. Utilise au moins 6 mots du glossaire de la leçon et la structure \zh{有} pour « il y a ».
\item \textbf{Tracer les caractères.} Choisis 3 caractères du thème (par exemple \zh{区}, \zh{市}, \zh{民}) et trace-les en grand sur l'affiche, en respectant l'ordre des traits ; indique la clé de chacun (ex. la clé 囗 « enceinte » pour 区, la clé 亠/巾 pour 市, la clé 氏 pour 民).
\item \textbf{Présenter oralement.} Chaque groupe présente son affiche en 1 minute : un élève lit le titre et les phrases, un deuxième montre les caractères tracés, un troisième conclut avec la phrase d'opinion. Parle lentement et soigne les tons.
\item \textbf{Voter et évaluer.} La classe vote pour la présentation la plus claire. L'enseignant évalue avec la grille : exactitude du vocabulaire (5 pts), tons et prononciation (5 pts), ordre des traits (5 pts), fluidité et présentation (5 pts).
\end{enumerate}

\begin{attention}[Pièges à éviter]
\begin{itemize}
\item Ne dis pas \zh{我的社区有很多人们} : \zh{人} n'a pas de pluriel en \zh{们} après un nom de groupe ; écris simplement \zh{我的社区有很多人} (Wǒ de shèqū yǒu hěn duō rén).
\item Respecte l'ordre Sujet + Verbe + Complément : \zh{居民喜欢踢足球} (Jūmín xǐhuan tī zúqiú, « les habitants aiment jouer au football »), jamais le verbe à la fin comme en français parfois.
\item Attention aux tons de \zh{shèqū} (社区, 4\textsuperscript{e} + 1\textsuperscript{er} ton) et \zh{jūmín} (居民, 1\textsuperscript{er} + 2\textsuperscript{e} ton) : des tons faux changent le sens.
\item La ponctuation chinoise est pleine chasse : utilise \zh{，} et \zh{。}, pas la virgule et le point français.
\end{itemize}
\end{attention}

\courssub{Ressources à mobiliser}

\begin{aretenir}[Boîte à outils de l'affiche]
\textbf{Vocabulaire utile :} \zh{社区} shèqū (communauté), \zh{城市} chéngshì (ville), \zh{街区} jiēqū (quartier), \zh{居民} jūmín (habitant), \zh{服务} fúwù (service), \zh{医院} yīyuàn (hôpital), \zh{学校} xuéxiào (école), \zh{市场} shìchǎng (marché), \zh{干净} gānjìng (propre), \zh{友好} yǒuhǎo (amicaux).\\[4pt]
\textbf{Structures :}
\begin{itemize}
\item Il y a : Sujet (lieu) + \zh{有} + nom → \zh{我们街区有一个大市场。}
\item Les habitants : \zh{居民} + \zh{很} + adjectif → \zh{居民很友好。}
\item Activité : \zh{喜欢} + verbe → \zh{年轻人喜欢踢足球。}
\item Opinion : \zh{我觉得} + phrase → \zh{我觉得我的社区很漂亮。}
\end{itemize}
\textbf{Caractères :} ordre des traits = de haut en bas, de gauche à droite, l'horizontal avant le vertical, l'enceinte (囗) se referme en dernier.
\end{aretenir}

\courssub{Proposition de production et corrigé commenté}

\begin{exoresolu}[Affiche modèle : le quartier de Mvolyé, Yaoundé]
Énoncé : rédige l'affiche complète d'un groupe (titre + 5 phrases) pour le quartier Mvolyé à Yaoundé, puis justifie tes choix de langue.
\tcblower
\textbf{Affiche modèle :}\\[4pt]
\zh{欢迎来到我们的街区！}\\
Huānyíng lái dào wǒmen de jiēqū !\\
Bienvenue dans notre quartier !\\[4pt]
\zh{我们的街区在雅温得，不太大，但是很干净。}\\
Wǒmen de jiēqū zài Yǎwēndé, bú tài dà, dànshì hěn gānjìng.\\
Notre quartier se trouve à Yaoundé ; il n'est pas très grand, mais il est très propre.\\[4pt]
\zh{街区里有一个市场、一所学校和一个医院。}\\
Jiēqū lǐ yǒu yí ge shìchǎng, yì suǒ xuéxiào hé yí ge yīyuàn.\\
Dans le quartier, il y a un marché, une école et un hôpital.\\[4pt]
\zh{居民很友好，年轻人喜欢踢足球。}\\
Jūmín hěn yǒuhǎo, niánqīngrén xǐhuan tī zúqiú.\\
Les habitants sont très amicaux ; les jeunes aiment jouer au football.\\[4pt]
\zh{我觉得我们的街区是雅温得最好的社区！}\\
Wǒ juéde wǒmen de jiēqū shì Yǎwēndé zuì hǎo de shèqū !\\
Je trouve que notre quartier est la meilleure communauté de Yaoundé !\\[8pt]
\textbf{Commentaire des choix de langue :}
\begin{itemize}
\item Le titre reprend la formule de bienvenue \zh{欢迎来到……}, polie et adaptée à l'accueil d'une délégation.
\item La phrase 1 utilise \zh{在} + lieu pour situer le quartier, et \zh{不太……但是……} pour nuancer (« pas très grand, mais propre »).
\item La phrase 2 mobilise la structure \zh{有} avec trois classificateurs corrects : \zh{个} pour le marché et l'hôpital, \zh{所} pour l'école (classificateur des bâtiments d'enseignement).
\item La phrase 3 répond à la question « qu'aiment faire les habitants ? » avec \zh{喜欢} + verbe.
\item La phrase 4 conclut par une opinion avec \zh{我觉得} et le superlatif \zh{最} : un point de grammaire de niveau Terminale bien réinvesti.
\item Caractères tracés sur l'affiche : \zh{区} (clé 匚, 4 traits), \zh{市} (clé 亠, 5 traits), \zh{民} (clé 氏, 5 traits) — ordre des traits vérifié par le groupe avant le collage.
\end{itemize}
\end{exoresolu}

\courssub{Bilan de l'activité}

Grâce à cette tâche, tu as réinvesti l'ensemble des savoir-faire de la leçon : tu as \cle{compris} un message authentique en chinois, \cle{sélectionné} le vocabulaire des collectivités locales, \cle{rédigé} des phrases correctes (ordre des mots, classificateurs, ponctuation chinoise), \cle{tracé} des caractères en respectant clés et ordre des traits, et \cle{présenté oralement} ton quartier à des visiteurs chinois. Tu es maintenant capable de décrire ta commune en chinois simple — une compétence concrète de médiation entre le Cameroun et la Chine. Pour aller plus loin : enregistre ta présentation en audio et compare tes tons avec ceux du modèle.

\newpage

\coursec{Méthodes et exercices résolus}

\coursec{Leçon 17 — Vocabulaire et compréhension de texte : collectivités locales}

\courssub{Mise en route et découverte du thème}

\begin{experience}[Brainstorming : ma collectivité locale]
\begin{enumerate}
\item \textbf{Observation d'images (projetées ou imprimées) :} une rue de Yaoundé avec le bâtiment de la \zh{社区居委会} (comité de quartier), le marché central de Douala (\zh{市场}), un village du Nord-Ouest (\zh{村}) avec sa chefferie, un panneau \zh{市政府} (mairie/hôtel de ville) à Buea.
\item \textbf{Mise en mots (oral, par paires) :} « Que voyez-vous ? Quels mots chinois connaissez-vous déjà pour nommer ces lieux et ces personnes ? » L'enseignant note au tableau les mots émergents : \zh{城市}、\zh{村庄}、\zh{市长}、\zh{居民}、\zh{医院}、\zh{学校}…
\item \textbf{Annonce de l'objectif :} « Aujourd'hui, nous allons structurer ce lexique, lire un texte authentique sur un quartier de Yaoundé, maîtriser les clés d'écriture des caractères-clés et produire une présentation de votre propre collectivité. »
\end{enumerate}
\end{experience}

\courssub{Texte support : lecture et compréhension globale}

\begin{methode}[Lire un texte sur les collectivités locales en 5 étapes]
\begin{enumerate}
\item \textbf{Repérer le thème au titre et aux mots-clés.} Les mots comme \zh{社区} (shèqū, communauté/quartier), \zh{城市} (chéngshì, ville), \zh{村} (cūn, village), \zh{政府} (zhèngfǔ, gouvernement/mairie) annoncent le sujet. Souligne-les mentalement.
\item \textbf{Faire une première lecture globale, sans dictionnaire.} Ne bloque pas sur un caractère inconnu : devine son sens grâce au contexte et à sa clé (radical). Exemple : \zh{河} a la clé de l'eau \zh{氵}, donc il parle d'un cours d'eau.
\item \textbf{Identifier les informations essentielles :} qui (qui agit ?), où (quel lieu ?), quoi (quelle action ?), pourquoi (quel but ?). Pour un texte sur une collectivité, cherche : le nom du lieu, les habitants, les services (école, hôpital, marché), les projets.
\item \textbf{Repérer les mots-outils grammaticaux} : \zh{的} (de, lien de possession ou de description), \zh{在} (zài, être situé / marque de lieu), \zh{有} (yǒu, il y a), \zh{很} (hěn, très / liaison adjectif). Ils structurent les phrases.
\item \textbf{Reformuler en français} phrase par phrase, en respectant l'ordre chinois : Sujet + Temps/Lieu + Verbe + Complément. Ne traduis jamais mot à mot.
\end{enumerate}
\end{methode}

\begin{textebox}[我的社区 — Mon quartier (Yaoundé)]
我住在雅温得的一个社区。我们的社区不太大，但是很干净。社区里有一所学校、一个医院和一个大市场。周末，居民们常常一起打扫街道，也一起唱歌跳舞。我很喜欢我的社区。\\
\emph{Wǒ zhù zài Yǎwēndé de yí ge shèqū. Wǒmen de shèqū bú tài dà, dànshì hěn gānjìng. Shèqū lǐ yǒu yì suǒ xuéxiào, yí ge yīyuàn hé yí ge dà shìchǎng. Zhōumò, jūmínmen chángcháng yìqǐ dǎsǎo jiēdào, yě yìqǐ chànggē tiàowǔ. Wǒ hěn xǐhuan wǒ de shèqū.}\\
« J'habite dans un quartier de Yaoundé. Notre quartier n'est pas très grand, mais il est très propre. Dans le quartier, il y a une école, un hôpital et un grand marché. Le week-end, les habitants nettoient souvent les rues ensemble, et chantent et dansent aussi ensemble. J'aime beaucoup mon quartier. »
\end{textebox}

\begin{exoresolu}[Compréhension globale du texte]
\textbf{Énoncé.} Lis le texte ci-dessus et réponds en français : (a) Où habite l'auteur ? (b) Quels services y a-t-il dans son quartier ? (c) Que font les habitants le week-end ?
\tcblower
\textbf{Solution détaillée.}
\begin{enumerate}
\item \textbf{Repérage du thème :} le titre \zh{我的社区} (wǒ de shèqū, « mon quartier ») annonce une présentation de collectivité locale.
\item \textbf{Question (a) :} la première phrase \zh{我住在雅温得的一个社区} contient la structure \zh{住在} + lieu (habiter à). Réponse : l'auteur habite dans un quartier de \textbf{Yaoundé} (\zh{雅温得}).
\item \textbf{Question (b) :} la structure \zh{社区里有} (dans le quartier, il y a) introduit une liste avec \zh{、} (virgule chinoise de liste) : \zh{一所学校} (une école, classificateur \zh{所} pour les établissements), \zh{一个医院} (un hôpital, classificateur \zh{个}), \zh{一个大市场} (un grand marché). Réponse : une école, un hôpital et un grand marché.
\item \textbf{Question (c) :} le mot-temps \zh{周末} (zhōumò, week-end) est placé en tête de phrase (règle : temps avant le verbe). Les actions : \zh{打扫街道} (dǎsǎo jiēdào, nettoyer les rues), \zh{唱歌跳舞} (chànggē tiàowǔ, chanter et danser), avec l'adverbe \zh{一起} (yìqǐ, ensemble). Réponse : le week-end, les habitants nettoient les rues ensemble et chantent et dansent ensemble.
\end{enumerate}
\textbf{Conclusion :} le texte présente un quartier propre et solidaire ; les structures-clés à retenir sont \zh{住在} + lieu, \zh{……里有} + liste, et Temps + Sujet + Verbe.
\end{exoresolu}

\courssub{Glossaire du thème : vocabulaire des collectivités locales}

\begin{definition}[Vocabulaire thématique — Collectivités locales (HSK 2-3)]
\begin{center}
\begin{tabularx}{\linewidth}{|X|X|X|X|}
\hline
\rowcolor{popblueL} \textbf{Caractères} & \textbf{Pinyin} & \textbf{Nature} & \textbf{Traduction} \\ \hline
城市 & chéngshì & nom & ville \\ \hline
村 & cūn & nom & village \\ \hline
社区 & shèqū & nom & quartier, communauté (urbaine) \\ \hline
区 & qū & nom & district, zone, arrondissement \\ \hline
市场 & shìchǎng & nom & marché \\ \hline
超市 & chāoshì & nom & supermarché \\ \hline
学校 & xuéxiào & nom & école \\ \hline
医院 & yīyuàn & nom & hôpital \\ \hline
政府 & zhèngfǔ & nom & gouvernement, mairie (niveau local) \\ \hline
市长 & shìzhǎng & nom & maire \\ \hline
居民 & jūmín & nom & habitant, résident \\ \hline
邻居 & línjū & nom & voisin \\ \hline
家乡 & jiāxiāng & nom & ville natale, région d'origine \\ \hline
环境 & huánjìng & nom & environnement \\ \hline
街道 & jiēdào & nom & rue, avenue \\ \hline
河 & hé & nom & rivière, fleuve \\ \hline
住 & zhù & verbe & habiter, résider \\ \hline
建设 & jiànshè & verbe & construire, bâtir, développer \\ \hline
打扫 & dǎsǎo & verbe & balayer, nettoyer (sol, rue) \\ \hline
保护 & bǎohù & verbe & protéger \\ \hline
参观 & cānguān & verbe & visiter (lieu officiel) \\ \hline
干净 & gānjìng & adj. & propre \\ \hline
脏 & zāng & adj. & sale \\ \hline
友好 & yǒuhǎo & adj. & amical, aimable \\ \hline
著名 & zhùmíng & adj. & célèbre \\ \hline
在 & zài & prép. & à, dans (marque le lieu d'action) \\ \hline
一起 & yìqǐ & adv. & ensemble \\ \hline
常常 & chángcháng & adv. & souvent \\ \hline
因为 & yīnwèi & conj. & parce que \\ \hline
所以 & suǒyǐ & conj. & donc, par conséquent \\ \hline
\end{tabularx}
\end{center}
\end{definition}

\begin{aretenir}[Collocations usuelles (associations fixes à mémoriser)]
\begin{center}
\begin{tabularx}{\linewidth}{|X|X|X|}
\hline
\rowcolor{popgreenL} \textbf{Chinois (Caractères + Pinyin)} & \textbf{Structure} & \textbf{Français} \\ \hline
\zh{住在城市} (zhù zài chéngshì) & Verbe + Préposition + Lieu & habiter en ville \\ \hline
\zh{建设家乡} (jiànshè jiāxiāng) & Verbe + Objet & développer sa région natale \\ \hline
\zh{打扫卫生 / 打扫街道} (dǎsǎo wèishēng / jiēdào) & Verbe + Objet & faire le ménage / nettoyer les rues \\ \hline
\zh{保护环境} (bǎohù huánjìng) & Verbe + Objet & protéger l'environnement \\ \hline
\zh{去市场买菜} (qù shìchǎng mǎi cài) & Verbe + Lieu + Verbe + Objet & aller au marché faire des courses \\ \hline
\zh{社区里有……} (shèqū lǐ yǒu…) & Lieu + \zh{有} + Liste & dans le quartier il y a… \\ \hline
\zh{一所学校 / 一家医院 / 一个市场} & Chiffre + Classifieur + Nom & une école / un hôpital / un marché \\ \hline
\zh{因为……所以……} (yīnwèi… suǒyǐ…) & Conjonction corrélative & parce que… donc… \\ \hline
\end{tabularx}
\end{center}
\end{aretenir}

\begin{popfigure}
\begin{tikzpicture}[mindmap, grow cyclic, every node/.style=concept, concept color=popblue!60,
    level 1/.append style={level distance=3.5cm, sibling angle=90},
    level 2/.append style={level distance=2.8cm, sibling angle=45}]
\node[align=center]{Collectivités\\ locales}
    child[concept color=poporange!60] { node {Lieux}
        child { node {城市} }
        child { node {村} }
        child { node {社区} }
        child { node {市场} }
        child { node {学校} }
        child { node {医院} }
    }
    child[concept color=popgreen!60] { node {Personnes}
        child { node {居民} }
        child { node {市长} }
        child { node {邻居} }
    }
    child[concept color=poppurple!60] { node {Actions}
        child { node {住在} }
        child { node {建设} }
        child { node {打扫} }
        child { node {保护环境} }
    }
    child[concept color=poppink!60] { node {Description}
        child { node {干净} }
        child { node {友好} }
        child { node {著名} }
    };
\end{tikzpicture}
\legende{Carte mentale du vocabulaire : familles de sens (lieux, personnes, actions, description).}
\end{popfigure}

\courssub{Clés (radicaux) et ordre des traits}

\begin{methode}[Utiliser les clés pour lire, écrire et classer les caractères]
\begin{enumerate}
\item \textbf{Identifie la clé (radical)} : elle donne souvent le sens général ou la catégorie sémantique. Exemples du thème : \zh{土} (terre) dans \zh{城} ; \zh{亻} (personne) dans \zh{住} ; \zh{氵} (eau) dans \zh{河} et \zh{洗} ; \zh{宀} (toit) dans \zh{家} ; \zh{匚} (boîte) dans \zh{区} ; \zh{巾} (étoffe/borne) dans \zh{市}.
\item \textbf{Repère la partie phonétique} (souvent à droite ou en bas) qui indique la prononciation approximative : \zh{城} = 土 + \zh{成} (chéng) ; \zh{河} = 氵+ \zh{可} (kě/ké).
\item \textbf{Respecte l'ordre des traits universel} : de haut en bas (从上到下), de gauche à droite (从左到右), horizontal avant vertical (先横后竖), extérieur avant intérieur (先外后内), fermer la boîte en dernier (最后封口).
\item \textbf{Entraîne-toi} : écris chaque caractère 5 fois en énonçant le numéro du trait, puis utilise-le dans une phrase personnelle.
\end{enumerate}
\end{methode}

\begin{exemplebox}[Analyse de caractères clés du thème]
\begin{itemize}
\item \zh{城} (chéng, ville / rempart) — 6 traits. Clé \zh{土} (terre, 3 traits) à gauche : les anciennes villes étaient des enceintes de terre. Partie phonétique \zh{成} (chéng, accomplir). Ordre : traits de 土 (横, 竖, 横) puis traits de 成.
\item \zh{市} (shì, ville / marché) — 5 traits. Clé \zh{巾} (perche/borne, 3 traits) en bas : à l'origine, borne marquant la zone de marché. Haut : \zh{亠} (point + horizontal). Ordre : 亠 (点, 横) puis 巾 (竖, 横折钩, 竖).
\item \zh{区} (qū, district) — 4 traits. Clé \zh{匚} (boîte ouverte, 2 traits). Intérieur : \zh{ㄨ} (2 traits). Ordre : 1. 横 (haut boîte), 2. 横折 (côté gauche + bas), 3. 横 (travers intérieur), 4. 竖 (vertical central qui ferme).
\item \zh{河} (hé, fleuve) — 8 traits. Clé \zh{氵} (eau, 3 traits) à gauche. Partie phonétique \zh{可} (5 traits) à droite. Ordre : 氵(点, 点, 提) puis 可.
\item \zh{住} (zhù, habiter) — 7 traits. Clé \zh{亻} (personne, 2 traits) à gauche. Partie phonétique \zh{主} (zhǔ, maître, 5 traits) à droite. Ordre : 亻(撇, 捺) puis 主 (点, 横, 横, 横, 竖).
\end{itemize}
\end{exemplebox}

\begin{exoresolu}[Clés et ordre des traits — Exercices]
\textbf{Énoncé.} (a) Donne la clé (radical), son sens et le nombre de traits pour : \zh{城}, \zh{河}, \zh{家}, \zh{住}, \zh{区}, \zh{市}. (b) Décris précisément l'ordre des traits de \zh{区} et \zh{市}.
\tcblower
\textbf{Solution détaillée.}
\begin{enumerate}
\item \zh{城} : clé \zh{土} (tǔ, terre), sens « mur de terre / enceinte », 6 traits.
\item \zh{河} : clé \zh{氵} (shuǐ, eau), sens « cours d'eau », 8 traits.
\item \zh{家} : clé \zh{宀} (mián, toit), sens « foyer / maison », 10 traits.
\item \zh{住} : clé \zh{亻} (rén, personne), sens « action humaine / résidence », 7 traits.
\item \zh{区} : clé \zh{匚} (fāng, boîte / enceinte), sens « zone délimitée », 4 traits.
\item \zh{市} : clé \zh{巾} (jīn, borne / étoffe), sens « zone de marché délimitée », 5 traits.
\item \textbf{Ordre \zh{区} (4 traits) :} 1. \zh{一} (horizontal haut), 2. \zh{(trait brisé)} (horizontal-plié vertical descendant, forme le côté gauche et le bas de la boîte), 3. \zh{一} (horizontal intérieur), 4. \zh{丨} (vertical central descendant, ferme la boîte).
\item \textbf{Ordre \zh{市} (5 traits) :} 1. \zh{丶} (point du \zh{亠}), 2. \zh{一} (horizontal du \zh{亠}), 3. \zh{丨} (vertical du \zh{巾}), 4. \zh{(trait brisé)} (horizontal-plié-crochet du \zh{巾}), 5. \zh{丨} (vertical central final du \zh{巾}).
\end{enumerate}
\end{exoresolu}

\courssub{Questions de compréhension et collocations}

\begin{methode}[Répondre à une question de compréhension en chinois]
\begin{enumerate}
\item \textbf{Analyse le mot interrogatif :} \zh{谁} (shéi, qui), \zh{什么} (shénme, quoi), \zh{哪儿/哪里} (nǎr/nǎli, où), \zh{什么时候} (shénme shíhou, quand), \zh{为什么} (wèishénme, pourquoi), \zh{怎么样} (zěnmeyàng, comment).
\item \textbf{Localise dans le texte la phrase qui contient la réponse} grâce aux mots-clés de la question.
\item \textbf{Réponds en reprenant la structure de la question} : en chinois, le mot interrogatif occupe la place exacte de la réponse. Question : \zh{你住在哪儿？} → Réponse : \zh{我住在杜阿拉。} (même ordre, on remplace \zh{哪儿} par le lieu).
\item \textbf{Pour \zh{为什么}, réponds obligatoirement avec \zh{因为}} (yīnwèi) : structure \zh{因为……, 所以……} (parce que…, donc…). En français on évite « parce que… donc », en chinois c'est la norme.
\item \textbf{Vérifie la phrase finale :} sujet présent, verbe correct, tons corrects, ponctuation chinoise \zh{。}.
\end{enumerate}
\end{methode}

\begin{textebox}[杜阿拉的市场 — Le marché de Douala]
杜阿拉是喀麦隆的大城市。那里的市场非常有名，每天有很多人去买东西。市场附近有一条河，居民们在河边洗衣服、聊天。市政府想建设一个新的市场，因为旧市场太小了。\\
\emph{Dùālā shì Kāmàilóng de dà chéngshì. Nàlǐ de shìchǎng fēicháng yǒumíng, měitiān yǒu hěn duō rén qù mǎi dōngxi. Shìchǎng fùjìn yǒu yì tiáo hé, jūmínmen zài hé biān xǐ yīfu, liáotiān. Shì zhèngfǔ xiǎng jiànshè yí ge xīn de shìchǎng, yīnwèi jiù shìchǎng tài xiǎo le.}\\
« Douala est une grande ville du Cameroun. Son marché est très célèbre, chaque jour beaucoup de gens y vont faire des achats. Près du marché, il y a une rivière ; les habitants lavent le linge et bavardent au bord de la rivière. La mairie veut construire un nouveau marché, parce que l'ancien marché est trop petit. »
\end{textebox}

\begin{exoresolu}[Questions de compréhension détaillées]
\textbf{Énoncé.} Lis le texte \zh{杜阿拉的市场}, puis réponds aux questions en chinois, par des phrases complètes avec ponctuation.
\begin{enumerate}
\item 杜阿拉的市场怎么样？
\item 居民们在河边做什么？
\item 为什么市政府想建设新市场？
\end{enumerate}
\tcblower
\textbf{Solution détaillée.}
\begin{enumerate}
\item \textbf{Q1 :} \zh{怎么样} demande une évaluation → adjectif \zh{非常有名}. Réponse : \zh{杜阿拉的市场非常有名。} \emph{Dùālā de shìchǎng fēicháng yǒumíng.}
\item \textbf{Q2 :} \zh{做什么} demande l'action → texte : \zh{洗衣服、聊天}. Complément de lieu \zh{在河边} placé \textbf{avant} le verbe. Réponse : \zh{居民们在河边洗衣服、聊天。} \emph{Jūmínmen zài hé biān xǐ yīfu, liáotiān.}
\item \textbf{Q3 :} \zh{为什么} demande la cause → texte : \zh{因为旧市场太小了}. Structure corrélative \zh{因为……所以……}. Réponse : \zh{因为旧市场太小了，所以市政府想建设一个新市场。} \emph{Yīnwèi jiù shìchǎng tài xiǎo le, suǒyǐ shì zhèngfǔ xiǎng jiànshè yí ge xīn de shìchǎng.}
\end{enumerate}
\textbf{Point de grammaire :} noter \zh{太……了} (tài… le, « trop… » / excès) et \zh{附近} (fùjìn, « près de / aux environs ») utilisé comme nom de lieu après le référence (\zh{市场附近}).
\end{exoresolu}

\begin{exercice}[Entraînement : vocabulaire et structures]
\textbf{1. Complète les phrases avec le mot convenable (caractères) :} \zh{居民 / 市长 / 建设 / 社区 / 医院 / 打扫 / 环境 / 邻居}
\begin{enumerate}
\item 我们\underline{\hspace{3cm}}很友好，经常一起活动。
\item 妈妈每天早上去\underline{\hspace{3cm}}买菜。
\item \underline{\hspace{3cm}}先生昨天来我们学校视察。
\item 年轻人应该\underline{\hspace{3cm}}自己的家乡。
\item 为了保护\underline{\hspace{3cm}}，大家不乱扔垃圾。
\item 周末我们全家一起\underline{\hspace{3cm}}卫生。
\end{enumerate}

\textbf{2. Relie chaque nom à son classificateur correct :}
\begin{center}
\begin{tabular}{cc}
\zh{学校} & \zh{个} \\
\zh{医院} & \zh{所} \\
\zh{市场} & \zh{家} \\
\zh{市长} & \zh{位} \\
\zh{河} & \zh{条}
\end{tabular}
\end{center}

\textbf{3. Traduis en chinois (caractères + pinyin) :}
\begin{enumerate}
\item Les habitants de mon village protègent l'environnement.
\item Le week-end, nous nettoyons les rues ensemble.
\item Parce que l'hôpital est loin, la mairie veut en construire un nouveau.
\end{enumerate}
\end{exercice}

\begin{corrige}[Exercice Leçon 17]
\textbf{1. Complétion :}
1. \zh{居民} (jūmín) — « Les habitants de notre quartier sont amicaux. » 2. \zh{市场} (shìchǎng) — « Maman va au marché… » 3. \zh{市长} (shìzhǎng) — « Monsieur le maire… » 4. \zh{建设} (jiànshè) — « Les jeunes doivent construire… » 5. \zh{环境} (huánjìng) — « Pour protéger l'environnement… » 6. \zh{打扫} (dǎsǎo) — « …nettoyons la maison/les lieux. »

\textbf{2. Classificateurs :}
\zh{一所学校} (yì suǒ xuéxiào) — \zh{所} pour écoles/établissements.
\zh{一家医院} (yì jiā yīyuàn) — \zh{家} pour hôpitaux, banques, restaurants, entreprises.
\zh{一个市场} (yí ge shìchǎng) — \zh{个} classifieur général (ou \zh{家} pour grand marché établi).
\zh{一位市长} (yí wèi shìzhǎng) — \zh{位} classifieur respectueux pour personnes.
\zh{一条河} (yì tiáo hé) — \zh{条} pour choses longues (rivières, routes, pantalons).

\textbf{3. Traduction :}
1. \zh{我村的居民保护环境。} \emph{Wǒ cūn de jūmín bǎohù huánjìng.} (Ou \zh{我家乡的居民…})
2. \zh{周末，我们一起打扫街道。} \emph{Zhōumò, wǒmen yìqǐ dǎsǎo jiēdào.}
3. \zh{因为医院太远，所以市政府想建设一个新医院。} \emph{Yīnwèi yīyuàn tài yuǎn, suǒyǐ shì zhèngfǔ xiǎng jiànshè yí ge xīn yīyuàn.}
\end{corrige}

\courssub{Élément socioculturel : organisation administrative Cameroun / Chine}

\begin{savaistu}[Collectivités locales : regards croisés Cameroun — Chine]
\begin{itemize}
\item \textbf{Cameroun (décentralisation) :} Le pays est divisé en 10 \zh{大区} (régions), subdivisées en \zh{部} (départements), puis en \zh{区} (arrondissements) et \zh{市} (communes/villes). L'exécutif local est le \zh{市长} (maire) élu, assisté d'adjoints. Le \zh{社区} (quartier) est une subdivision informelle ou administrative de base (comité de quartier \zh{居委会}).
\item \textbf{Chine (administration territoriale) :} 3 niveaux principaux : \zh{省} (province), \zh{市} (ville de niveau préfecture), \zh{县} (district/canton). En zone urbaine : \zh{市辖区} (district urbain), \zh{街道} (sous-district/rues), \zh{社区} (communauté de quartier — niveau de base, géré par \zh{居委会} comité de résidents). En zone rurale : \zh{乡} (canton), \zh{镇} (bourg), \zh{村} (village — géré par \zh{村委会} comité villageois).
\item \textbf{Point commun :} le \zh{社区} / \zh{村} est la cellule de base où s'exerce la démocratie participative (budget participatif, entretien, sécurité, entraide). Au Cameroun comme en Chine, le \zh{市长} / \zh{区长} / \zh{镇长} est l'élu local clé.
\item \textbf{Différence notable :} la Chine a un maillage très dense de \zh{街道} et \zh{社区} en ville (quelques km², quelques milliers d'habitants), avec un personnel administratif professionnel (\zh{社区工作者}) salarié par l'État. Au Cameroun, le quartier repose davantage sur le bénévolat des notables et chefs de quartier.
\end{itemize}
\end{savaistu}

\courssub{Production guidée : ma collectivité locale}

\begin{methode}[Rédiger un paragraphe sur sa collectivité locale (niveau Bac)]
\begin{enumerate}
\item \textbf{Plan en 4-5 phrases types :}
\begin{enumerate}
\item Localisation : \zh{我住在……} (J'habite à…)
\item Description : \zh{这里……，有……} (Ici…, il y a…) — utiliser \zh{很/不太} + adj., \zh{有} + liste avec classificateurs.
\item Vie quotidienne / Services : \zh{居民们常常……} (Les habitants souvent…) — verbe d'action + lieu (\zh{在} + lieu + verbe).
\item Activité collective : \zh{周末/节日，大家一起……} (Week-end/fête, tous ensemble…)
\item Opinion justifiée : \zh{我喜欢……，因为……} (J'aime…, parce que…)
\end{enumerate}
\item \textbf{Règles d'or de l'ordre des mots :}
\begin{itemize}
\item Sujet + \zh{在} + Lieu + Verbe (jamais Verbe + \zh{在} + Lieu).
\item Temps (周末) + Sujet + Verbe.
\item \zh{很} obligatoire devant adjectif attribut (pas de verbe « être »).
\item \zh{因为}…, \zh{所以}… pour cause/conséquence.
\end{itemize}
\item \textbf{Check-list avant de rendre :} 5+ mots du thème, 2+ classificateurs corrects, 1 structure \zh{因为……所以……}, ponctuation chinoise, pinyin mental vérifié.
\end{enumerate}
\end{methode}

\begin{exoresolu}[Production guidée — Modèle de rédaction]
\textbf{Énoncé.} Rédige un paragraphe de 5-6 phrases en chinois pour présenter ta collectivité locale (ville, village ou quartier) : lieu, description, services, activité collective, opinion justifiée.
\tcblower
\textbf{Solution détaillée — Modèle de rédaction (niveau Bac Terminale).}

\zh{我住在布埃亚，它是喀麦隆西部的一个城市。我的城市不太大，但是很美丽，也很干净。城市里有一所大学和一个大市场，居民们常常去市场买东西。周末，大家一起打扫街道，保护环境。我非常喜欢我的城市，因为这里的居民很友好，环境也好。}

\emph{Wǒ zhù zài Bùāiyà, tā shì Kāmàilóng xībù de yí ge chéngshì. Wǒ de chéngshì bú tài dà, dànshì hěn měilì, yě hěn gānjìng. Chéngshì lǐ yǒu yì suǒ dàxué hé yí ge dà shìchǎng, jūmínmen chángcháng qù shìchǎng mǎi dōngxi. Zhōumò, dàjiā yìqǐ dǎsǎo jiēdào, bǎohù huánjìng. Wǒ fēicháng xǐhuan wǒ de chéngshì, yīnwèi zhèlǐ de jūmín hěn yǒuhǎo, huánjìng yě hǎo.}

« J'habite à Buea, qui est une ville de l'ouest du Cameroun. Ma ville n'est pas très grande, mais elle est belle et propre. Dans la ville, il y a une université et un grand marché ; les habitants vont souvent au marché faire des achats. Le week-end, tout le monde nettoie les rues ensemble et protège l'environnement. J'aime beaucoup ma ville, parce que ses habitants sont très amicaux et l'environnement est aussi bon. »

\textbf{Justifications grammaticales et lexicales :}
\begin{itemize}
\item \zh{住在布埃亚} : structure \zh{住在} + lieu propre (transcription phonétique \zh{布埃亚}).
\item \zh{一所大学} : classificateur \zh{所} (suǒ) obligatoire pour établissements d'enseignement supérieur ; \zh{一个市场} : \zh{个} (gè) classifieur général.
\item \zh{不太大，但是很美丽} : \zh{不太} (bú tài, « pas très ») + adj. ; \zh{很} (hěn) lie le sujet à l'adjectif (pas de \zh{是}).
\item \zh{周末，大家一起……} : Temps en tête ; \zh{大家} (dàjiā, « tout le monde ») sujet collectif ; \zh{一起} (yìqǐ) avant le verbe.
\item \zh{因为这里的居民很友好，环境也好} : cause double coordonnée ; \zh{也} (yě, « aussi ») placé avant l'adjectif.
\end{itemize}
\textbf{Évaluation :} Ce paragraphe valide les compétences « Expression écrite » (cohérence, vocabulaire HSK 3, structures ciblées) et « Caractères » (écriture correcte, classificateurs).
\end{exoresolu}

\begin{experience}[Production orale interactive : « Mon quartier idéal »]
\textbf{Consigne (par binômes, 5 min préparation + 2 min passage) :}
\begin{enumerate}
\item L'élève A est un journaliste de \zh{喀麦隆电视台} (CRTV), l'élève B est le \zh{市长} ou un \zh{居民代表} d'une collectivité (Yaoundé, Douala, Buea, village natal…).
\item Le journaliste pose 3 questions obligatoires : \zh{你们社区/城市有什么特点？} (Quelles sont les particularités ?) — \zh{居民们平时做什么活动？} (Quelles activités ?) — \zh{未来有什么建设计划？} (Quels projets futurs ?).
\item L'interviewé répond en phrases complètes, en utilisant \textbf{au moins 5 mots du glossaire} et la structure \zh{因为……所以……}.
\item \textbf{Évaluation par les pairs (grille) :} Prononciation/tons (5 pts), Vocabulaire/Structures (5 pts), Fluidité/Interaction (5 pts), Contenu/Intérêt (5 pts).
\end{enumerate}
\end{experience}

\begin{attention}[Les 5 erreurs qui coûtent des points au Bac — Rappel final]
\begin{enumerate}
\item \textbf{Oublier \zh{很} devant l'adjectif attributif.} \ding{55} \zh{我的社区大} \quad \ding{51} \zh{我的社区\cle{很}大} (wǒ de shèqū \textbf{hěn} dà). \zh{很} est la liaison standard (copule adjectivale), même sens affaibli « très ».
\item \textbf{Placer le complément de lieu \textbf{après} le verbe (calque français).} \ding{55} \zh{我买东西在市场} \quad \ding{51} \zh{我\cle{在市场}买东西} (wǒ \textbf{zài shìchǎng} mǎi dōngxi). Règle : \zh{在} + Lieu \textbf{avant} le verbe d'action.
\item \textbf{Confondre caractères proches (un trait change tout).}
\begin{itemize}
\item \zh{市} (shì, ville/marché) vs \zh{币} (bì, monnaie) — clé \zh{巾} vs \zh{贝}.
\item \zh{住} (zhù, habiter) vs \zh{往} (wǎng, vers) — clé \zh{亻} vs \zh{彳}.
\item \zh{河} (hé, rivière) vs \zh{何} (hé, interrogatif) — clé \zh{氵} vs \zh{亻}.
\item \zh{区} (qū, district) vs \zh{欧} (ōu, Europe) — clé \zh{匚} vs \zh{欠}.
\end{itemize}
\item \textbf{Négliger les classificateurs (numéral + classifieur + nom obligatoire).} \ding{55} \zh{一学校} \quad \ding{51} \zh{一\cle{所}学校} ; \ding{55} \zh{一市长} \quad \ding{51} \zh{一\cle{位}市长}.
\item \textbf{Faux tons sur paires minimales — inversion de sens garantie.}
\begin{itemize}
\item \zh{买} (mǎi, 3\textsuperscript{e} ton, acheter) vs \zh{卖} (mài, 4\textsuperscript{e} ton, vendre).
\item \zh{市} (shì, 4\textsuperscript{e} ton, ville) vs \zh{是} (shì, 4\textsuperscript{e} ton, être) vs \zh{事} (shì, 4\textsuperscript{e} ton, affaire) — même ton, sens diffèrent par contexte ; mais \zh{十} (shí, 2\textsuperscript{e} ton, dix) se distingue au ton.
\item \zh{河} (hé, 2\textsuperscript{e} ton) vs \zh{喝} (hē, 1\textsuperscript{er} ton, boire).
\end{itemize}
\textbf{Astuce :} dicte-toi le vocabulaire en enregistrant ton téléphone, réécoute et compare au modèle audio.
\end{enumerate}
\end{attention}

\newpage

\coursec{Exercices}

\courssub{Je vérifie mes connaissances}

\begin{exercice}[QCM : le bon mot]
Pour chaque question, entoure la bonne réponse.

1. Le mot qui signifie « ville » est :
\begin{itemize}
\item a) \zh{村} \quad b) \zh{城} \quad c) \zh{家} \quad d) \zh{区}
\end{itemize}

2. \zh{居民} (jūmín) signifie :
\begin{itemize}
\item a) les touristes \quad b) les habitants \quad c) les élèves \quad d) les voisins
\end{itemize}

3. \zh{市场} (shìchǎng) désigne :
\begin{itemize}
\item a) une école \quad b) un hôpital \quad c) un marché \quad d) une mairie
\end{itemize}

4. Le mot qui signifie « organiser » est :
\begin{itemize}
\item a) \zh{参加} \quad b) \zh{组织} \quad c) \zh{生活} \quad d) \zh{欢迎}
\end{itemize}

5. \zh{友好} (yǒuhǎo) signifie :
\begin{itemize}
\item a) amical \quad b) lointain \quad c) moderne \quad d) célèbre
\end{itemize}
\end{exercice}

\begin{exercice}[Vrai ou faux ? Justifie]
Réponds par V (vrai) ou F (faux) et justifie ta réponse en français.

1. \zh{区} (qū) signifie « quartier, arrondissement ».
2. \zh{住} (zhù) est un nom qui signifie « maison ».
3. \zh{参加} (cānjiā) signifie « participer à ».
4. \zh{政府} (zhèngfǔ) signifie « école ».
5. \zh{活动} (huódòng) signifie « activité ».
\end{exercice}

\begin{exercice}[Vocabulaire : complète la traduction]
Complète le tableau avec la traduction française qui convient.

\begin{center}
\begin{tabularx}{\linewidth}{|X|X|X|X|}
\hline
\rowcolor{popblueL} Caractères & Pinyin & Nature & Traduction \\
\hline
\zh{城市} & chéngshì & nom & \ldots \\
\hline
\zh{农村} & nóngcūn & nom & \ldots \\
\hline
\zh{社区} & shèqū & nom & \ldots \\
\hline
\zh{环境} & huánjìng & nom & \ldots \\
\hline
\zh{干净} & gānjìng & adjectif & \ldots \\
\hline
\end{tabularx}
\end{center}
\end{exercice}

\begin{exercice}[Associe caractères et pinyin]
Relie chaque caractère ou mot à son pinyin.

\begin{center}
\begin{tabularx}{\linewidth}{|X|X|}
\hline
\rowcolor{popblueL} Caractères & Pinyin (dans le désordre) \\
\hline
1. \zh{村} & a) zhù \\
\hline
2. \zh{民} & b) qū \\
\hline
3. \zh{住} & c) cūn \\
\hline
4. \zh{区} & d) chéng \\
\hline
5. \zh{城} & e) mín \\
\hline
\end{tabularx}
\end{center}
\end{exercice}

\courssub{Je m'entraîne}

\begin{exercice}[Phrases à trous]
Complète les phrases avec les mots de la banque : \zh{住在、有、组织、参加、友好}.

1. 我\ldots 雅温得的一个小区。\\
\emph{Wǒ \ldots Yǎwēndé de yí ge xiǎoqū.}

2. 我们村\ldots 一所学校和一个大市场。\\
\emph{Wǒmen cūn \ldots yì suǒ xuéxiào hé yí ge dà shìchǎng.}

3. 社区每个月\ldots 一次清洁活动。\\
\emph{Shèqū měi ge yuè \ldots yí cì qīngjié huódòng.}

4. 很多居民\ldots 了昨天的会议。\\
\emph{Hěn duō jūmín \ldots le zuótiān de huìyì.}

5. 我们村的人非常\ldots 。\\
\emph{Wǒmen cūn de rén fēicháng \ldots .}
\end{exercice}

\begin{exercice}[Remets les mots en ordre]
Remets les mots dans le bon ordre pour faire une phrase correcte, puis traduis-la en français.

1. 在 / 住 / 我 / 杜阿拉 / 的 / 一个 / 区

2. 市场 / 有 / 村 / 我们 / 一个 / 里

3. 居民 / 友好 / 很 / 这里 / 的

4. 活动 / 组织 / 社区 / 了 / 一次 / 文化
\end{exercice}

\begin{exercice}[Appariement : mots et expressions]
Associe chaque mot chinois à la bonne expression française.

\begin{center}
\begin{tabularx}{\linewidth}{|X|X|}
\hline
\rowcolor{popblueL} Chinois & Français (dans le désordre) \\
\hline
1. \zh{市政府} & a) la campagne \\
\hline
2. \zh{农村} & b) les transports \\
\hline
3. \zh{交通} & c) la mairie \\
\hline
4. \zh{医院} & d) l'hôpital \\
\hline
5. \zh{学校} & e) l'école \\
\hline
\end{tabularx}
\end{center}
\end{exercice}

\begin{exercice}[Traduction de phrases]
Traduis les phrases suivantes en chinois (caractères + pinyin).

1. J'habite dans un quartier de Yaoundé.
2. Il y a un marché dans notre village.
3. Les habitants de notre communauté sont très amicaux.
4. La mairie organise une activité culturelle.
\end{exercice}

\courssub{Je me prépare au Bac}

\begin{exercice}[Compréhension de texte : un village chinois]
Lis le texte suivant, puis réponds aux questions.

\begin{textebox}[我的家乡]
我的家乡是中国南方的一个小村庄。村里有八百多个居民，大家都很友好。\\
\emph{Wǒ de jiāxiāng shì Zhōngguó nánfāng de yí ge xiǎo cūnzhuāng. Cūn li yǒu bā bǎi duō ge jūmín, dàjiā dōu hěn yǒuhǎo.}\\
村子中间有一个市场，旁边有一所小学和一个小医院。\\
\emph{Cūnzi zhōngjiān yǒu yí ge shìchǎng, pángbiān yǒu yì suǒ xiǎoxué hé yí ge xiǎo yīyuàn.}\\
以前村里的路不好，交通很不方便。\\
\emph{Yǐqián cūn li de lù bù hǎo, jiāotōng hěn bù fāngbiàn.}\\
去年，政府修了新的路，还组织了很多文化活动。\\
\emph{Qùnián, zhèngfǔ xiū le xīn de lù, hái zǔzhī le hěn duō wénhuà huódòng.}\\
现在我们的生活越来越好，环境也越来越干净。\\
\emph{Xiànzài wǒmen de shēnghuó yuè lái yuè hǎo, huánjìng yě yuè lái yuè gānjìng.}
\end{textebox}

\textbf{Questions de compréhension} (réponds en chinois, en phrases complètes) :

1. 作者的家乡在哪儿？\\
\emph{Zuòzhě de jiāxiāng zài nǎr ?}

2. 村子里有多少居民？\\
\emph{Cūnzi li yǒu duōshao jūmín ?}

3. 市场旁边有什么？\\
\emph{Shìchǎng pángbiān yǒu shénme ?}

4. Vrai ou faux : avant, les routes du village étaient bonnes. Justifie avec une phrase du texte.

5. 去年政府做了什么？\\
\emph{Qùnián zhèngfǔ zuò le shénme ?}

\textbf{Question de langue} :

6. Releve dans le texte deux phrases avec la structure \zh{越来越} et traduis-les en français.
\end{exercice}

\begin{exercice}[Thème et version]
\textbf{A. Thème.} Traduis en chinois (caractères + pinyin) :

« J'habite dans un village près de Bafoussam. Il y a une école et un marché. Les habitants sont très amicaux. »

\textbf{B. Version.} Traduis en français le paragraphe suivant :

\begin{textebox}[我们的社区]
我们的社区在杜阿拉市中心。社区里有学校、医院和一个很大的市场。\\
\emph{Wǒmen de shèqū zài Dùlālā shì zhōngxīn. Shèqū li yǒu xuéxiào, yīyuàn hé yí ge hěn dà de shìchǎng.}\\
居民们常常参加社区组织的活动。\\
\emph{Jūmínmen chángcháng cānjiā shèqū zǔzhī de huódòng.}\\
这里的交通很方便，环境也很干净。\\
\emph{Zhèli de jiāotōng hěn fāngbiàn, huánjìng yě hěn gānjìng.}\\
我很喜欢我的社区。\\
\emph{Wǒ hěn xǐhuan wǒ de shèqū.}
\end{textebox}
\end{exercice}

\begin{exercice}[Expression écrite guidée : ma collectivité locale]
\textbf{Sujet :} Ton école participe à un échange avec un lycée chinois. Ton correspondant chinois te demande de lui présenter ta commune (ville, quartier ou village).

Rédige un texte de \textbf{6 à 8 phrases en chinois} (environ 60 à 80 caractères) dans lequel tu :
\begin{itemize}
\item dis où tu habites ;
\item décris les lieux importants de ta collectivité (école, marché, hôpital, mairie\ldots) ;
\item parles des habitants et des activités ;
\item utilises \textbf{au moins 5 mots du glossaire} de la leçon ;
\item utilises \textbf{au moins une phrase avec} \zh{因为} (yīnwèi, parce que).
\end{itemize}

N'oublie pas d'écrire les caractères avec soin, en respectant l'ordre des traits appris dans la leçon (\zh{城、区、村、民、住}).
\end{exercice}

\newpage

\coursec{Corrigés des exercices}

\begin{corrige}[Exercice 1 — QCM : le bon mot]
1. \textbf{b) \zh{城}} (chéng) signifie « ville ». \zh{村} (cūn) = village ; \zh{家} (jiā) = maison, famille ; \zh{区} (qū) = quartier, arrondissement.

2. \textbf{b) les habitants}. \zh{居民} (jūmín) = habitants, résidents. Les touristes = \zh{游客} (yóukè) ; les élèves = \zh{学生} (xuésheng) ; les voisins = \zh{邻居} (línjū).

3. \textbf{c) un marché}. \zh{市场} (shìchǎng) = marché. Une école = \zh{学校} ; un hôpital = \zh{医院} ; une mairie = \zh{市政府}.

4. \textbf{b) \zh{组织}} (zǔzhī) = organiser. \zh{参加} (cānjiā) = participer à ; \zh{生活} (shēnghuó) = vivre, la vie ; \zh{欢迎} (huānyíng) = accueillir, bienvenue.

5. \textbf{a) amical}. \zh{友好} (yǒuhǎo) = amical. Lointain = \zh{远} (yuǎn) ; moderne = \zh{现代} (xiàndài) ; célèbre = \zh{有名} (yǒumíng).
\end{corrige}

\begin{corrige}[Exercice 2 — Vrai ou faux ? Justifie]
1. \textbf{V}. \zh{区} (qū) signifie bien « quartier, arrondissement, district ». Exemple : \zh{雅温得的一个区} (Yǎwēndé de yí ge qū) « un quartier de Yaoundé ».

2. \textbf{F}. \zh{住} (zhù) est un \textbf{verbe} qui signifie « habiter, résider », pas un nom. « Maison » se dit \zh{家} (jiā) ou \zh{房子} (fángzi). On dit : \zh{我住在杜阿拉。} (Wǒ zhù zài Dùlālā.) « J'habite à Douala. »

3. \textbf{V}. \zh{参加} (cānjiā) signifie bien « participer à ». Attention : c'est un verbe transitif direct en chinois, sans préposition : \zh{参加会议} (cānjiā huìyì) « participer à la réunion » (et non *\zh{参加在会议}).

4. \textbf{F}. \zh{政府} (zhèngfǔ) signifie « gouvernement ». « École » se dit \zh{学校} (xuéxiào).

5. \textbf{V}. \zh{活动} (huódòng) signifie bien « activité ». Exemple : \zh{文化活动} (wénhuà huódòng) « activité culturelle ».
\end{corrige}

\begin{corrige}[Exercice 3 — Vocabulaire : complète la traduction]
\begin{center}
\begin{tabularx}{\linewidth}{|X|X|X|X|}
\hline
\rowcolor{popblueL} Caractères & Pinyin & Nature & Traduction \\
\hline
\zh{城市} & chéngshì & nom & la ville \\
\hline
\zh{农村} & nóngcūn & nom & la campagne, le monde rural \\
\hline
\zh{社区} & shèqū & nom & la communauté, le quartier résidentiel \\
\hline
\zh{环境} & huánjìng & nom & l'environnement, le cadre de vie \\
\hline
\zh{干净} & gānjìng & adjectif & propre \\
\hline
\end{tabularx}
\end{center}
\textbf{Points de vigilance :} ne pas confondre \zh{城市} (ville, en général) et \zh{市} (municipalité, dans \zh{市政府} « mairie ») ; \zh{农村} se compose de \zh{农} (agriculture) + \zh{村} (village) : littéralement « village agricole ».
\end{corrige}

\begin{corrige}[Exercice 4 — Associe caractères et pinyin]
1. \zh{村} → \textbf{c) cūn} (village)

2. \zh{民} → \textbf{e) mín} (peuple, habitant)

3. \zh{住} → \textbf{a) zhù} (habiter)

4. \zh{区} → \textbf{b) qū} (quartier, arrondissement)

5. \zh{城} → \textbf{d) chéng} (ville)

\textbf{Points de vigilance :} attention aux tons : \zh{住} zhù (4\ieme{} ton) ≠ \zh{猪} zhū (1\ier{} ton, « cochon ») ; \zh{区} qū prend ü après q, mais on écrit « qu » sans les deux points (règle orthographique du pinyin : q, j, x + ü → u).
\end{corrige}

\begin{corrige}[Exercice 5 — Phrases à trous]
1. 我\zh{住在}雅温得的一个小区。\\
\emph{Wǒ zhù zài Yǎwēndé de yí ge xiǎoqū.} « J'habite dans un petit quartier résidentiel de Yaoundé. »\\
Justification : le verbe \zh{住} s'emploie avec la préposition \zh{在} pour introduire le lieu : \textbf{habiter + 在 + lieu}.

2. 我们村\zh{有}一所学校和一个大市场。\\
\emph{Wǒmen cūn yǒu yì suǒ xuéxiào hé yí ge dà shìchǎng.} « Notre village a une école et un grand marché. »\\
Justification : \zh{有} exprime la possession ou l'existence (« il y a »). Classificateurs : \zh{所} pour les établissements (école, hôpital), \zh{个} pour marché.

3. 社区每个月\zh{组织}一次清洁活动。\\
\emph{Shèqū měi ge yuè zǔzhī yí cì qīngjié huódòng.} « La communauté organise une activité de nettoyage chaque mois. »\\
Justification : le complément de fréquence \zh{每个月} se place avant le verbe ; \zh{一次} se place après le verbe.

4. 很多居民\zh{参加}了昨天的会议。\\
\emph{Hěn duō jūmín cānjiā le zuótiān de huìyì.} « Beaucoup d'habitants ont participé à la réunion d'hier. »\\
Justification : \zh{了} marque l'action accomplie (la réunion a eu lieu, hier).

5. 我们村的人非常\zh{友好}。\\
\emph{Wǒmen cūn de rén fēicháng yǒuhǎo.} « Les gens de notre village sont très amicaux. »\\
Justification : adjectif attribut précédé de l'adverbe de degré \zh{非常}, sans copule « être ».
\end{corrige}

\begin{corrige}[Exercice 6 — Remets les mots en ordre]
1. 我住在杜阿拉的一个区。\\
\emph{Wǒ zhù zài Dùlālā de yí ge qū.} « J'habite dans un quartier de Douala. »\\
Structure : Sujet + \zh{住在} + lieu + \zh{的} + classificateur + nom.

2. 我们村里有一个市场。\\
\emph{Wǒmen cūn li yǒu yí ge shìchǎng.} « Il y a un marché dans notre village. »\\
Structure : Lieu + \zh{里} + \zh{有} + complément (phrase d'existence).

3. 这里的居民很友好。\\
\emph{Zhèli de jūmín hěn yǒuhǎo.} « Les habitants d'ici sont très amicaux. »\\
Structure : Sujet + adverbe de degré \zh{很} + adjectif.

4. 社区组织了一次文化活动。\\
\emph{Shèqū zǔzhī le yí cì wénhuà huódòng.} « La communauté a organisé une activité culturelle. »\\
Structure : Sujet + verbe + \zh{了} + complément d'objet.

\textbf{Point de vigilance :} en chinois, le lieu introduit par \zh{在} se place \textbf{avant} le verbe principal (sauf avec \zh{住} où \zh{在} + lieu suit le verbe : \zh{住在}).
\end{corrige}

\begin{corrige}[Exercice 7 — Appariement : mots et expressions]
1. \zh{市政府} (shì zhèngfǔ) → \textbf{c) la mairie} (littéralement « gouvernement municipal »)

2. \zh{农村} (nóngcūn) → \textbf{a) la campagne}

3. \zh{交通} (jiāotōng) → \textbf{b) les transports}

4. \zh{医院} (yīyuàn) → \textbf{d) l'hôpital}

5. \zh{学校} (xuéxiào) → \textbf{e) l'école}

\textbf{Point de vigilance :} ne pas confondre \zh{医院} (yīyuàn, hôpital) et \zh{学院} (xuéyuàn, institut, faculté) : mêmes sons à un ton près (yī / xué) mais sens très différents.
\end{corrige}

\begin{corrige}[Exercice 8 — Traduction de phrases]
1. 我住在雅温得的一个区。\\
\emph{Wǒ zhù zài Yǎwēndé de yí ge qū.}\\
(On accepte aussi : 我住在雅温得的一个小区。avec \zh{小区} « quartier résidentiel ».)

2. 我们村里有一个市场。\\
\emph{Wǒmen cūn li yǒu yí ge shìchǎng.}\\
Structure d'existence : lieu + \zh{里} + \zh{有} + nom.

3. 我们社区的居民很友好。\\
\emph{Wǒmen shèqū de jūmín hěn yǒuhǎo.}\\
\zh{的} marque l'appartenance (« les habitants \emph{de} notre communauté »).

4. 市政府组织了一次文化活动。\\
\emph{Shì zhèngfǔ zǔzhī le yí cì wénhuà huódòng.}\\
\zh{了} marque l'accompli ; on peut aussi le mettre en projet sans \zh{了} : \zh{市政府要组织一次文化活动。} « La mairie va organiser une activité culturelle. »

\textbf{Barème indicatif (8 points) :} 2 points par phrase (1 point pour la structure correcte, 1 point pour le vocabulaire et les caractères).
\end{corrige}

\begin{corrige}[Exercice 9 — Compréhension de texte : un village chinois]
\textbf{Barème indicatif (12 points) :} 2 points par question (1 point pour le sens, 1 point pour la correction de la langue).

1. 作者的家乡在中国南方的一个小村庄。\\
\emph{Zuòzhě de jiāxiāng zài Zhōngguó nánfāng de yí ge xiǎo cūnzhuāng.} « Le pays natal de l'auteur est un petit village du sud de la Chine. »

2. 村子里有八百多个居民。\\
\emph{Cūnzi li yǒu bā bǎi duō ge jūmín.} « Il y a plus de huit cents habitants dans le village. »\\
(\zh{多} après le nombre signifie « plus de ».)

3. 市场旁边有一所小学和一个小医院。\\
\emph{Shìchǎng pángbiān yǒu yì suǒ xiǎoxué hé yí ge xiǎo yīyuàn.} « À côté du marché, il y a une école primaire et un petit hôpital. »

4. \textbf{Faux.} Justification avec la phrase du texte : \zh{以前村里的路不好，交通很不方便。} (Yǐqián cūn li de lù bù hǎo, jiāotōng hěn bù fāngbiàn.) « Avant, les routes du village n'étaient pas bonnes, les transports étaient très peu pratiques. »

5. 去年，政府修了新的路，还组织了很多文化活动。\\
\emph{Qùnián, zhèngfǔ xiū le xīn de lù, hái zǔzhī le hěn duō wénhuà huódòng.} « L'année dernière, le gouvernement a construit de nouvelles routes et a aussi organisé beaucoup d'activités culturelles. »

6. Les deux phrases avec \zh{越来越} :
\begin{itemize}
\item \zh{现在我们的生活越来越好。} (Xiànzài wǒmen de shēnghuó yuè lái yuè hǎo.) « Maintenant, notre vie est de mieux en mieux. »
\item \zh{环境也越来越干净。} (Huánjìng yě yuè lái yuè gānjìng.) « L'environnement aussi est de plus en plus propre. »
\end{itemize}
Structure : \textbf{越来越 + adjectif} = « de plus en plus + adjectif ». Attention : on ne dit jamais *\zh{越来越很好} (pas de \zh{很} dans cette structure).
\end{corrige}

\begin{corrige}[Exercice 10 — Thème et version]
\textbf{Barème indicatif (10 points) :} thème 5 points, version 5 points.

\textbf{A. Thème — proposition de traduction :}

我住在巴富萨姆附近的一个村庄。村里有一所学校和一个市场。居民们很友好。\\
\emph{Wǒ zhù zài Bāfùsàmǔ fùjìn de yí ge cūnzhuāng. Cūn li yǒu yì suǒ xuéxiào hé yí ge shìchǎng. Jūmínmen hěn yǒuhǎo.}

Points de vigilance :
\begin{itemize}
\item « près de » = \zh{附近} (fùjìn), placé après le nom de lieu : \zh{巴富萨姆附近} ;
\item classificateur \zh{所} pour l'école, \zh{个} pour le marché ;
\item pas de verbe « être » devant l'adjectif : \zh{居民们很友好} et non *\zh{居民们是很友好}.
\end{itemize}

\textbf{B. Version — proposition de traduction :}

« Notre communauté se trouve au centre-ville de Douala. Dans la communauté, il y a des écoles, un hôpital et un très grand marché. Les habitants participent souvent aux activités organisées par la communauté. Ici, les transports sont très pratiques et l'environnement est très propre. J'aime beaucoup ma communauté. »

Points de vigilance :
\begin{itemize}
\item \zh{市中心} = « centre-ville » (et non « centre de la ville » au sens administratif) ;
\item \zh{社区组织的活动} = « les activités organisées par la communauté » : en chinois, le complément du verbe \zh{组织} précède le nom grâce à \zh{的} ;
\item \zh{常常} = « souvent » ; \zh{方便} = « pratique, commode ».
\end{itemize}
\end{corrige}

\begin{corrige}[Exercice 11 — Expression écrite guidée : ma collectivité locale]
\textbf{Barème indicatif (20 points) :} respect de la tâche et des consignes (5 pts) ; correction grammaticale (5 pts) ; richesse du vocabulaire, dont au moins 5 mots du glossaire (5 pts) ; écriture des caractères et présentation (3 pts) ; présence d'une phrase avec \zh{因为} (2 pts).

\textbf{Proposition de rédaction modèle :}

我住在喀麦隆的杜阿拉。我的社区在市中心，很大也很热闹。\\
\emph{Wǒ zhù zài Kāmàilóng de Dùlālā. Wǒ de shèqū zài shì zhōngxīn, hěn dà yě hěn rènao.}\\
社区里有一所学校、一个医院和一个大市场。\\
\emph{Shèqū li yǒu yì suǒ xuéxiào, yí ge yīyuàn hé yí ge dà shìchǎng.}\\
这里的居民都很友好，我们常常一起参加活动。\\
\emph{Zhèli de jūmín dōu hěn yǒuhǎo, wǒmen chángcháng yìqǐ cānjiā huódòng.}\\
市政府每个月组织一次清洁活动，所以环境越来越干净。\\
\emph{Shì zhèngfǔ měi ge yuè zǔzhī yí cì qīngjié huódòng, suǒyǐ huánjìng yuè lái yuè gānjìng.}\\
我很喜欢我的社区，因为这里的生活很方便。\\
\emph{Wǒ hěn xǐhuan wǒ de shèqū, yīnwèi zhèli de shēnghuó hěn fāngbiàn.}\\
欢迎你来杜阿拉看看！\\
\emph{Huānyíng nǐ lái Dùlālā kànkan !}

Traduction : « J'habite à Douala, au Cameroun. Ma communauté est au centre-ville, elle est grande et très animée. Dans la communauté, il y a une école, un hôpital et un grand marché. Les habitants d'ici sont tous très amicaux, nous participons souvent ensemble aux activités. La mairie organise une activité de nettoyage chaque mois, donc l'environnement est de plus en plus propre. J'aime beaucoup ma communauté, parce que la vie y est très pratique. Bienvenue à Douala pour venir voir ! »

\textbf{Points de vigilance :}
\begin{itemize}
\item Vérifier la présence des 5 mots du glossaire minimum : ici \zh{社区、居民、友好、活动、组织、环境、市场、医院} (8 mots utilisés) ;
\item la phrase avec \zh{因为} : \zh{因为} + cause, souvent en corrélation avec \zh{所以} (donc) ; on peut utiliser \zh{因为} seul ou le couple \zh{因为……所以……} ;
\item ordre des traits : \zh{城} (clé 土 « terre » à gauche, puis 成 à droite) ; \zh{区} (le trait horizontal d'abord, l'intérieur 乂, puis le trait vertical-branlant qui ferme) ; \zh{村} (clé 木 « bois » à gauche, puis 寸) ; \zh{民} (5 traits, de haut en bas) ; \zh{住} (clé 亻 « homme » à gauche, puis 主) ;
\item ne pas oublier les classificateurs : \zh{一所学校}, \zh{一个市场}, \zh{一次活动}.
\end{itemize}
\end{corrige}

\newpage

\coursec{Fiche bilan}

\begin{popfigure}
\begin{tikzpicture}[mindmap, grow cyclic, every node/.style={concept, font=\small}, concept color=popblue, root concept/.append style={font=\small\bfseries, minimum size=2.6cm}, level 1 concept/.append style={sibling angle=51, minimum size=2.1cm, font=\scriptsize, level distance=3.4cm}]
\node[root concept]{Collectivités locales \zh{地方政府}}
  child[concept color=poporange]{ node {Lieux \zh{市政府} \zh{区政府} \zh{村委会} \zh{社区}} }
  child[concept color=popgreen]{ node {Personnes \zh{市长} \zh{区长} \zh{村长} \zh{居民}} }
  child[concept color=poppurple]{ node {Services \zh{服务} \zh{医院} \zh{学校} \zh{市场}} }
  child[concept color=poppink]{ node {Actions \zh{建设} \zh{管理} \zh{解决} \zh{改善}} }
  child[concept color=popgold]{ node {Structures S + 为 + O + V ; 越来越 + Adj} }
  child[concept color=popblue!60]{ node {Radicaux 广 (bâtiment) 土 (terre) 亻 (personne)} }
  child[concept color=popgreen!60]{ node {Culture mairies au Cameroun / 市政府 en Chine} };
\end{tikzpicture}
\legende{Carte mentale de la leçon : les collectivités locales}
\end{popfigure}

\begin{bilan}
\textbf{1. Lexique clé à connaître par cœur}
\begin{itemize}
  \item \zh{地方} dìfāng : local, région ; \zh{政府} zhèngfǔ : gouvernement ; \zh{市政府} shì zhèngfǔ : mairie / gouvernement municipal
  \item \zh{市长} shìzhǎng : maire ; \zh{区长} qūzhǎng : chef de district ; \zh{村长} cūnzhǎng : chef de village
  \item \zh{社区} shèqū : communauté, quartier ; \zh{居民} jūmín : habitant, résident
  \item \zh{建设} jiànshè : construire ; \zh{管理} guǎnlǐ : gérer, administrer ; \zh{解决} jiějué : résoudre ; \zh{改善} gǎishàn : améliorer
  \item \zh{服务} fúwù : service, servir ; \zh{环境} huánjìng : environnement ; \zh{生活} shēnghuó : vie
\end{itemize}

\textbf{2. Règles de grammaire à connaître par cœur}

\begin{center}
\begin{tabularx}{\linewidth}{|X|X|X|}
\hline
\rowcolor{popblueL} Structure & Exemple & Traduction \\
\hline
S + 为 + O + V (faire qqch pour qqn) & \zh{政府为居民服务。} Zhèngfǔ wèi jūmín fúwù. & Le gouvernement sert les habitants. \\
\hline
越来越 + adjectif (de plus en plus) & \zh{我们的生活越来越好。} Wǒmen de shēnghuó yuèláiyuè hǎo. & Notre vie est de plus en plus belle. \\
\hline
S + 把 + O + V + complément & \zh{市政府把路修好了。} Shì zhèngfǔ bǎ lù xiū hǎo le. & La mairie a réparé la route. \\
\hline
Verb + 得 + adjectif (manière) & \zh{村长管理得很好。} Cūnzhǎng guǎnlǐ de hěn hǎo. & Le chef de village gère très bien. \\
\hline
\end{tabularx}
\end{center}

\textbf{3. Clés et écriture}
\begin{itemize}
  \item 广 (clé du « toit / bâtiment ») : \zh{府}, \zh{店} — traits : point, horizontal, puis crochet descendant à gauche.
  \item 土 (clé de la « terre ») : \zh{地}, \zh{城} — horizontal court, vertical, horizontal long.
  \item 亻 (clé de « l'homme », variante de 人) : \zh{们}, \zh{住} — deux traits : tombant puis vertical.
  \item Ordre général : haut → bas, gauche → droite, horizontal avant vertical, extérieur avant intérieur.
\end{itemize}

\textbf{4. Méthodes express}
\begin{itemize}
  \item \textbf{Comprendre un texte} : 1) lire le titre et la première phrase ; 2) repérer les mots connus (lieux, personnes) ; 3) deviner les mots inconnus par le contexte ; 4) répondre en recopiant la structure de la question.
  \item \textbf{Parler de sa collectivité} : lieu (\zh{我住在……}) → qui gère (\zh{市长/村长……}) → services (\zh{有医院、学校……}) → opinion (\zh{我觉得……}).
\end{itemize}
\end{bilan}

\begin{aretenir}[Checklist avant le Bac]
\begin{itemize}
  \item $\square$ Je sais écrire et prononcer \zh{地方政府}, \zh{市长}, \zh{社区}, \zh{居民} avec les bons tons.
  \item $\square$ Je connais au moins 15 mots du thème « collectivités locales » avec leur nature grammaticale.
  \item $\square$ Je maîtrise la structure S + 为 + O + V sans oublier le 为.
  \item $\square$ J'emploie correctement 越来越 + adjectif (jamais 越来越很 + adjectif !).
  \item $\square$ Je sais construire une phrase en 把 avec un complément après le verbe.
  \item $\square$ Je distingue 的, 得, 地 dans les exemples de la leçon.
  \item $\square$ Je reconnais les clés 广, 土, 亻 et je respecte l'ordre des traits.
  \item $\square$ Je sais lire un court texte et repérer les informations essentielles (qui, où, quoi).
  \item $\square$ Je peux présenter ma collectivité locale en 4 ou 5 phrases chinoises.
  \item $\square$ Je peux comparer une mairie camerounaise et un 市政府 chinois en deux phrases.
\end{itemize}
\end{aretenir}

\findecours

\end{document}
